\chapter{PROOFS AND SUPPLEMENTARY DETAILS FOR CHAPTER~\ref{ch:diffusion-policies}}
\label{app:diffusion-policies}

\aigen{This appendix provides the complete numerical experiment tables, coverage examples, rollout sampler specification, hybrid RL extension, score regularity lemmas, and regret proofs for Algorithm~\ref{alg:dppo-vanilla-pg}, Algorithm~\ref{alg:dppo-reinforce}, and Algorithm~\ref{alg:dipo-bc} in Chapter~\ref{ch:diffusion-policies}.}

\section{Tables for Numerical Experiments}
\label{app:tables}
In this section, we provide additional information for the numerical experiments in Section~\ref{ch05:sec:experiments} within Tables~\ref{tab:value}, \ref{tab:normalized-regret}, \ref{tab:regret-exponent} and \ref{tab:sublinear-regret}. We compute empirical regret relative to a reference return $V_{\mathrm{ref}}$, using $3500$, $5400$, and $4000$ for Hopper, HalfCheetah, and Walker2D, respectively, and success rate $1$ for Robomimic. For evaluations $\widehat V(u)$ logged every ten iterations, the plotted cumulative quantity is
\[
\widehat R(t)=10\sum_{u\le t}\frac{V_{\mathrm{ref}}-\widehat V(u)}{V_{\mathrm{ref}}},
\]
including the initial evaluation. This reference need not equal the optimal policy value. We use the first $1000$ iterations for Gym and the first $150$, $80$, and $200$ for Can, Lift, and Square, respectively.

For exponent inference, we use the interval-weighted proxy
\[
\widehat R_{\mathrm{int}}(t_i)=\sum_{j<i}(t_{j+1}-t_j)\frac{V_{\mathrm{ref}}-\widehat V(t_j)}{V_{\mathrm{ref}}},
\]
which is zero at the initial evaluation and assigns each evaluated policy to the following interval. Within the horizons above, we fit $\log\widehat R_{\mathrm{int}}$ against $\log t$ separately for each seed, retaining the last $\lceil n/2\rceil$ of its $n$ positive finite pairs. Tables~\ref{tab:regret-exponent}--\ref{tab:rate-tests} report this late-half sensitivity analysis. Using the full window instead changes the sublinearity count for the three theorem-associated methods from $15/18$ to $13/18$ after the same correction.
\begin{table}[H]
    \centering
    \small
    \setlength{\tabcolsep}{3pt}
    \renewcommand{\arraystretch}{1.15}
    \begin{tabular}{|l|c|c|c|c|c|c|}
        \hline 
        \bf Best Value & \bf \begin{tabular}[c]{@{}c@{}}Vanilla PG\\(Alg. 1)\end{tabular} & \bf DPPO+DDIM & \bf \begin{tabular}[c]{@{}c@{}}Modified DIPO\\(Alg. 3)\end{tabular} & \bf DPPO & \bf DIPO & \bf \begin{tabular}[c]{@{}c@{}}Gaussian\\PPO\end{tabular} \\
        \hline 
        Hopper-v2 & \begin{tabular}[c]{@{}c@{}}2993 \\ (2948, 3037)\end{tabular} & \begin{tabular}[c]{@{}c@{}}3121 \\ (3054, 3187)\end{tabular} & \begin{tabular}[c]{@{}c@{}}2957 \\ (2951, 2964)\end{tabular} & \begin{tabular}[c]{@{}c@{}}3233 \\ (3196, 3269)\end{tabular} & \begin{tabular}[c]{@{}c@{}}2936 \\ (2917, 2955)\end{tabular} & -- \\
        \hline 
        HalfCheetah-v2 & \begin{tabular}[c]{@{}c@{}}4999 \\ (4951, 5047)\end{tabular} & \begin{tabular}[c]{@{}c@{}}5127 \\ (5007, 5247)\end{tabular} & \begin{tabular}[c]{@{}c@{}}5062 \\ (4982, 5142)\end{tabular} & \begin{tabular}[c]{@{}c@{}}5297 \\ (5119, 5476)\end{tabular} & \begin{tabular}[c]{@{}c@{}}5028 \\ (4980, 5077)\end{tabular} & -- \\
        \hline 
        Walker2D-v2 & \begin{tabular}[c]{@{}c@{}}3678 \\ (3628, 3728)\end{tabular} & \begin{tabular}[c]{@{}c@{}}3763 \\ (3668, 3859)\end{tabular} & \begin{tabular}[c]{@{}c@{}}3852 \\ (3814, 3890)\end{tabular} & \begin{tabular}[c]{@{}c@{}}3859 \\ (3714, 4004)\end{tabular} & \begin{tabular}[c]{@{}c@{}}3702 \\ (3671, 3733)\end{tabular} & -- \\
        \hline 
        Can & \begin{tabular}[c]{@{}c@{}}0.979 \\ (0.970, 0.988)\end{tabular} & \begin{tabular}[c]{@{}c@{}}0.987 \\ (0.980, 0.994)\end{tabular} & \begin{tabular}[c]{@{}c@{}}0.868 \\ (0.850, 0.886)\end{tabular} & \begin{tabular}[c]{@{}c@{}}0.995 \\ (0.989, 1.001)\end{tabular} & \begin{tabular}[c]{@{}c@{}}0.841 \\ (0.824, 0.859)\end{tabular} & \begin{tabular}[c]{@{}c@{}}0.945 \\ (0.929, 0.961)\end{tabular} \\
        \hline 
        Lift & \begin{tabular}[c]{@{}c@{}}0.974 \\ (0.967, 0.981)\end{tabular} & \begin{tabular}[c]{@{}c@{}}0.988 \\ (0.978, 0.998)\end{tabular} & \begin{tabular}[c]{@{}c@{}}0.986 \\ (0.975, 0.997)\end{tabular} & \begin{tabular}[c]{@{}c@{}}0.999 \\ (0.997, 1.001)\end{tabular} & \begin{tabular}[c]{@{}c@{}}0.977 \\ (0.967, 0.987)\end{tabular} & \begin{tabular}[c]{@{}c@{}}0.950 \\ (0.927, 0.973)\end{tabular} \\
        \hline 
        Square & \begin{tabular}[c]{@{}c@{}}0.482 \\ (0.423, 0.542)\end{tabular} & \begin{tabular}[c]{@{}c@{}}0.988 \\ (0.982, 0.995)\end{tabular} & \begin{tabular}[c]{@{}c@{}}0.585 \\ (0.554, 0.616)\end{tabular} & \begin{tabular}[c]{@{}c@{}}0.994 \\ (0.992, 0.996)\end{tabular} & \begin{tabular}[c]{@{}c@{}}0.502 \\ (0.413, 0.592)\end{tabular} & \begin{tabular}[c]{@{}c@{}}0.925 \\ (0.907, 0.943)\end{tabular} \\
        \hline
    \end{tabular}

    \vspace{3pt}
    \caption{Average best value of each algorithm across all environments. 95\% confidence intervals over five seeds are shown in parentheses.}
    \label{tab:value}
\end{table}

\begin{table}[H]
    \centering
    \small
    \setlength{\tabcolsep}{3pt}
    \renewcommand{\arraystretch}{1.15}
    \begin{tabular}{|l|c|c|c|c|c|c|}
        \hline 
        \bf \begin{tabular}[c]{@{}c@{}}Normalized\\Regret\end{tabular} & \bf \begin{tabular}[c]{@{}c@{}}Vanilla PG\\(Alg. 1)\end{tabular} & \bf DPPO+DDIM & \bf \begin{tabular}[c]{@{}c@{}}Modified DIPO\\(Alg. 3)\end{tabular} & \bf DPPO & \bf DIPO & \bf \begin{tabular}[c]{@{}c@{}}Gaussian\\PPO\end{tabular} \\
        \hline 
        Hopper-v2 & \begin{tabular}[c]{@{}c@{}}0.23 \\ (0.22, 0.25)\end{tabular} & \begin{tabular}[c]{@{}c@{}}0.19 \\ (0.18, 0.20)\end{tabular} & \begin{tabular}[c]{@{}c@{}}0.25 \\ (0.25, 0.25)\end{tabular} & \begin{tabular}[c]{@{}c@{}}0.17 \\ (0.15, 0.18)\end{tabular} & \begin{tabular}[c]{@{}c@{}}0.26 \\ (0.26, 0.27)\end{tabular} & -- \\
        \hline 
        HalfCheetah-v2 & \begin{tabular}[c]{@{}c@{}}0.12 \\ (0.12, 0.12)\end{tabular} & \begin{tabular}[c]{@{}c@{}}0.10 \\ (0.09, 0.11)\end{tabular} & \begin{tabular}[c]{@{}c@{}}0.14 \\ (0.14, 0.15)\end{tabular} & \begin{tabular}[c]{@{}c@{}}0.08 \\ (0.06, 0.09)\end{tabular} & \begin{tabular}[c]{@{}c@{}}0.16 \\ (0.15, 0.16)\end{tabular} & -- \\
        \hline 
        Walker2D-v2 & \begin{tabular}[c]{@{}c@{}}0.12 \\ (0.12, 0.13)\end{tabular} & \begin{tabular}[c]{@{}c@{}}0.10 \\ (0.08, 0.12)\end{tabular} & \begin{tabular}[c]{@{}c@{}}0.13 \\ (0.11, 0.16)\end{tabular} & \begin{tabular}[c]{@{}c@{}}0.08 \\ (0.06, 0.10)\end{tabular} & \begin{tabular}[c]{@{}c@{}}0.23 \\ (0.23, 0.24)\end{tabular} & --\\
        \hline 
        Can & \begin{tabular}[c]{@{}c@{}}0.08 \\ (0.07, 0.08)\end{tabular} & \begin{tabular}[c]{@{}c@{}}0.08 \\ (0.07, 0.08)\end{tabular} & \begin{tabular}[c]{@{}c@{}}0.31 \\ (0.28, 0.35)\end{tabular} & \begin{tabular}[c]{@{}c@{}}0.04 \\ (0.04, 0.05)\end{tabular} & \begin{tabular}[c]{@{}c@{}}0.35 \\ (0.32, 0.38)\end{tabular} & \begin{tabular}[c]{@{}c@{}}0.18 \\ (0.17, 0.18)\end{tabular} \\
        \hline 
        Lift & \begin{tabular}[c]{@{}c@{}}0.07 \\ (0.06, 0.09)\end{tabular} & \begin{tabular}[c]{@{}c@{}}0.05 \\ (0.05, 0.06)\end{tabular} & \begin{tabular}[c]{@{}c@{}}0.06 \\ (0.06, 0.07)\end{tabular} & \begin{tabular}[c]{@{}c@{}}0.04 \\ (0.04, 0.05)\end{tabular} & \begin{tabular}[c]{@{}c@{}}0.07 \\ (0.06, 0.07)\end{tabular} & \begin{tabular}[c]{@{}c@{}}0.16 \\ (0.14, 0.18)\end{tabular} \\
        \hline 
        Square & \begin{tabular}[c]{@{}c@{}}0.65 \\ (0.61, 0.69)\end{tabular} & \begin{tabular}[c]{@{}c@{}}0.16 \\ (0.16, 0.16)\end{tabular} & \begin{tabular}[c]{@{}c@{}}0.66 \\ (0.63, 0.69)\end{tabular} & \begin{tabular}[c]{@{}c@{}}0.13 \\ (0.13, 0.13)\end{tabular} & \begin{tabular}[c]{@{}c@{}}0.68 \\ (0.66, 0.70)\end{tabular} & \begin{tabular}[c]{@{}c@{}}0.30 \\ (0.27, 0.34)\end{tabular} \\
        \hline
    \end{tabular}

    \vspace{5pt}
    \caption{Normalized average regret of each algorithm across all environments, averaged across seeds. We report $\widehat R(T)/T$, using the reference returns and evaluation horizons defined above. 95\% CIs across five seeds shown in parentheses.}
    \label{tab:normalized-regret}
\end{table}

\begin{table}[H]
    \centering
    \small
    \setlength{\tabcolsep}{3pt}
    \renewcommand{\arraystretch}{1.15}
    \begin{tabular}{|l|c|c|c|c|c|c|}
    \hline 
    \bf \begin{tabular}[c]{@{}c@{}}Fitted regret\\exponent\end{tabular} & \bf \begin{tabular}[c]{@{}c@{}}Vanilla PG\\(Alg. 1)\end{tabular} & \bf DPPO+DDIM & \bf \begin{tabular}[c]{@{}c@{}}Modified DIPO\\(Alg. 3)\end{tabular} & \bf DPPO & \bf DIPO & \bf \begin{tabular}[c]{@{}c@{}}Gaussian\\PPO\end{tabular} \\
    \hline 
    Hopper-v2 & \begin{tabular}[c]{@{}c@{}}0.744 \\ (0.625, 0.863)\end{tabular} & \begin{tabular}[c]{@{}c@{}}0.680 \\ (0.493, 0.868)\end{tabular} & \begin{tabular}[c]{@{}c@{}}0.637 \\ (0.330, 0.944)\end{tabular} & \begin{tabular}[c]{@{}c@{}}0.576 \\ (0.422, 0.730)\end{tabular} & \begin{tabular}[c]{@{}c@{}}0.820 \\ (0.606, 1.034)\end{tabular} & -- \\
    \hline 
    HalfCheetah-v2 & \begin{tabular}[c]{@{}c@{}}0.780 \\ (0.655, 0.905)\end{tabular} & \begin{tabular}[c]{@{}c@{}}0.609 \\ (0.443, 0.776)\end{tabular} & \begin{tabular}[c]{@{}c@{}}0.735 \\ (0.599, 0.870)\end{tabular} & \begin{tabular}[c]{@{}c@{}}0.557 \\ (0.262, 0.851)\end{tabular} & \begin{tabular}[c]{@{}c@{}}0.752 \\ (0.623, 0.881)\end{tabular} & -- \\
    \hline 
    Walker2D-v2 & \begin{tabular}[c]{@{}c@{}}0.756 \\ (0.664, 0.849)\end{tabular} & \begin{tabular}[c]{@{}c@{}}0.633 \\ (0.407, 0.860)\end{tabular} & \begin{tabular}[c]{@{}c@{}}0.852 \\ (0.305, 1.399)\end{tabular} & \begin{tabular}[c]{@{}c@{}}0.765 \\ (0.354, 1.177)\end{tabular} & \begin{tabular}[c]{@{}c@{}}1.253 \\ (1.088, 1.419)\end{tabular} & -- \\
    \hline 
    Can & \begin{tabular}[c]{@{}c@{}}0.401 \\ (0.313, 0.488)\end{tabular} & \begin{tabular}[c]{@{}c@{}}0.183 \\ (0.089, 0.277)\end{tabular} & \begin{tabular}[c]{@{}c@{}}0.621 \\ (0.573, 0.669)\end{tabular} & \begin{tabular}[c]{@{}c@{}}0.147 \\ (0.103, 0.190)\end{tabular} & \begin{tabular}[c]{@{}c@{}}0.686 \\ (0.635, 0.737)\end{tabular} & \begin{tabular}[c]{@{}c@{}}0.444 \\ (0.367, 0.521)\end{tabular} \\
    \hline 
    Lift & \begin{tabular}[c]{@{}c@{}}0.397 \\ (0.265, 0.529)\end{tabular} & \begin{tabular}[c]{@{}c@{}}0.273 \\ (0.205, 0.341)\end{tabular} & \begin{tabular}[c]{@{}c@{}}0.333 \\ (0.218, 0.449)\end{tabular} & \begin{tabular}[c]{@{}c@{}}0.073 \\ (0.035, 0.111)\end{tabular} & \begin{tabular}[c]{@{}c@{}}0.381 \\ (0.317, 0.444)\end{tabular} & \begin{tabular}[c]{@{}c@{}}0.339 \\ (0.209, 0.469)\end{tabular} \\
    \hline 
    Square & \begin{tabular}[c]{@{}c@{}}1.026 \\ (0.965, 1.087)\end{tabular} & \begin{tabular}[c]{@{}c@{}}0.406 \\ (0.060, 0.752)\end{tabular} & \begin{tabular}[c]{@{}c@{}}0.829 \\ (0.748, 0.910)\end{tabular} & \begin{tabular}[c]{@{}c@{}}0.222 \\ (0.056, 0.388)\end{tabular} & \begin{tabular}[c]{@{}c@{}}0.887 \\ (0.847, 0.928)\end{tabular} & \begin{tabular}[c]{@{}c@{}}0.507 \\ (0.257, 0.756)\end{tabular} \\
    \hline
\end{tabular}

    \vspace{5pt}
    \caption{Late-half fits of the interval-weighted regret exponent, with Student-$t$ 95\% confidence intervals across seeds in parentheses. Each cell has five seeds except Can/Gaussian PPO, where one seed has too few tail observations and four remain. These finite-horizon fits do not establish an asymptotic rate.}
    \label{tab:regret-exponent}
\end{table}

\begin{table}[H]
    \centering
    \small
    \setlength{\tabcolsep}{3pt}
    \renewcommand{\arraystretch}{1.15}
    \begin{tabular}{|l|c|c|c|c|c|c|}
        \hline \bf \begin{tabular}[c]{@{}c@{}}Sublinearity\\$p$-value\end{tabular} & \bf \begin{tabular}[c]{@{}c@{}}Vanilla PG\\(Alg. 1)\end{tabular} & \bf DPPO+DDIM & \bf \begin{tabular}[c]{@{}c@{}}Modified DIPO\\(Alg. 3)\end{tabular} & \bf DPPO & \bf DIPO & \bf \begin{tabular}[c]{@{}c@{}}Gaussian\\PPO\end{tabular} \\
        \hline Hopper-v2 & 0.0317 & 0.0444 & 0.0914 & 0.0149 & 0.199 & -- \\
        \hline HalfCheetah-v2 & 0.0441 & 0.0242 & 0.0376 & 0.0489 & 0.0376 & -- \\
        \hline Walker2D-v2 & 0.0165 & 0.0444 & 0.742 & 0.377 & 1 & -- \\
        \hline Can & 0.00061 & 0.000256 & 0.000361 & 1.08e-05 & 0.000873 & 0.00207 \\
        \hline Lift & 0.00232 & 0.00012 & 0.0011 & 4.52e-06 & 0.000167 & 0.00174 \\
        \hline Square & 1 & 0.0444 & 0.0318 & 0.00221 & 0.0145 & 0.0376 \\
        \hline
        \end{tabular}

    \vspace{3pt}
    \caption{One-sided Student-$t$ tests of $H_0:\alpha\ge1$ against $\alpha<1$, with Holm correction across all 33 environment--algorithm cells. Values below $0.05$ reject the null.}
    \label{tab:sublinear-regret}
\end{table}

\begin{table}[H]
\centering
\renewcommand{\arraystretch}{1.3}
\begin{tabular}{|l|c|c|c|}
\hline Environment & Vanilla PG & DPPO+DDIM & Modified DIPO \\
\hline Hopper-v2 & 1 & 1 & 1 \\
\hline HalfCheetah-v2 & 1 & 1 & 0.0644 \\
\hline Walker2D-v2 & 1 & 1 & 1 \\
\hline Can & 1 & 1 & 0.0175 \\
\hline Lift & 1 & 1 & 1 \\
\hline Square & 0.00212 & 1 & 0.00298 \\
\hline
\end{tabular}
\caption{One-sided Student-$t$ tests of $H_0:\alpha\le r$ against $\alpha>r$, with $r=3/4$ for Vanilla PG and DPPO+DDIM and $r=1/2$ for Modified DIPO. Reported $p$-values are Holm-adjusted across these 18 tests. Can/Modified DIPO, Square/Vanilla PG, and Square/Modified DIPO reject at $0.05$. Non-rejection does not establish the theorem rate or its assumptions.}
\label{tab:rate-tests}
\end{table}


\begin{table}[H]
\centering
\renewcommand{\arraystretch}{1.3}
\begin{tabular}{|l|r|r|r|}
\hline Environment & Gaussian PPO & Diffusion & Increase \\
\hline Can & 551,971 & 575,564 & 4.27\% \\
\hline Lift & 549,923 & 573,516 & 4.29\% \\
\hline Square & 2,152,483 & 2,305,212 & 7.10\% \\
\hline
\end{tabular}
\caption{Trainable actor parameters in the released Robomimic configurations. Counts include affine-layer weights and biases and the Gaussian actor's learned log variances, and exclude critics and frozen model copies.}
\label{tab:actor-counts}
\end{table}

To examine whether difficulty fitting the actor accompanies high regret, we compare the empirical denoising losses during training. Table~\ref{tab:actor-loss} divides the longest observed actor-update history in each environment--algorithm pair into deciles and uses the same seeds that reach the final decile throughout, normalizing each seed by its first-decile mean. The losses generally decline early and then plateau, but rise on Can. Square has declining loss and high regret, so the diagnostic does not establish a consistent relationship across tasks. These sampled training losses are distinct from the current-target population errors under the deployed state distribution in Theorem~\ref{thm:dipo-bc-regret}.
\begin{table}[H]
\centering
\renewcommand{\arraystretch}{1.3}
\begin{tabular}{|l|c|c|}
\hline Environment & DIPO & Modified DIPO \\
\hline Hopper-v2 & $100\to79\to80$ (4) & $100\to79\to78$ (2) \\
\hline HalfCheetah-v2 & $100\to80\to75$ (4) & $100\to83\to73$ (4) \\
\hline Walker2D-v2 & $100\to90\to88$ (4) & $100\to90\to86$ (4) \\
\hline Can & $100\to107\to108$ (4) & $100\to107\to108$ (5) \\
\hline Lift & $100\to85\to83$ (4) & $100\to85\to82$ (3) \\
\hline Square & $100\to87\to89$ (2) & $100\to89\to91$ (3) \\
\hline
\end{tabular}
\caption{Empirical denoising loss at the first, 40--50\%, and final training deciles, as percentages of each seed's first-decile mean. Parentheses give complete-seed counts.}
\label{tab:actor-loss}
\end{table}

\section{Further Remarks}
\label{app:remarks}

\paragraph{Discount factors and policy gradient bias.} Algorithm \ref{alg:dppo-sampler} uses geometric acceptance to sample the discounted occupancy and geometric return termination with probability $1-\gamma$ to estimate discounted values without bias, as in \citet{agarwal2021policygradient}. Using the undiscounted on-policy occupancy in place of the discounted occupancy can introduce the bias described in \citet{nota2020policygradientgradient}. In practice, it is possible to use the generalized advantage estimation (GAE) algorithm of \citet{schulman2018gae} and the PPO update of \citet{schulman2017proximalpolicyoptimizationalgorithms}, but the former may introduce bias with an inexact critic, and the latter approximates the trust-region update motivating NPG rather than the regression update analyzed here. Further, one may wonder if simply using a state-value baseline as in the REINFORCE with baseline algorithm allows one to achieve a similar result. A state-only baseline preserves the on-policy score-gradient expectation when conditionally independent of the sampled action and denoising path, even if it is a biased estimate of the state value. Its size still affects the estimator variance. This does not make a return minus an arbitrary baseline conditionally unbiased for the advantage. Appendix~\ref{app:hybrid-reset-dist} gives a paired baseline construction that preserves that stronger identity for the hybrid labels. 

\paragraph{Covering the policy class.} Compared to other actor-critic algorithms within recent literature \citep{zhong2023theoreticalanalysisoptimisticproximal, sherman2024rateoptimalpolicyoptimizationlinear, cassel2024warmupfreepolicyoptimization, tan2025actorcriticsachieveoptimalsample}, our policy-gradient analyses use Monte Carlo returns without fitting an off-policy critic, including importance-weighted logged returns for the hybrid branch in Appendix~\ref{app:hybrid-reset-dist}. Therefore, as in \citet{agarwal2021policygradient}, this allows us to bypass the issue of covering an exponentially growing function class (induced by Bellman backups of the critic with respect to the policy class) encountered by \citep{zhong2023theoreticalanalysisoptimisticproximal, cassel2024warmupfreepolicyoptimization, sherman2024rateoptimalpolicyoptimizationlinear, tan2025actorcriticsachieveoptimalsample}, which is necessary when establishing a PAC guarantee for a critic estimated with off-policy data. The hybrid analysis instead controls importance-weighted return moments and uses reserved independent trajectory blocks to handle adaptive policy updates. 

% Otherwise, as in Remark 6.8 of \citet{agarwal2021policygradient}, if we simply have an unbiased estimate of the Q-function and some arbitrary baseline that only depends on the current state as in the REINFORCE with baseline algorithm, we incur an additional factor of $|\gA|$ in the final sample complexity. 

\paragraph{Where does the complexity of the diffusion model show up in the bound?} The complexity of the policy class, as well as the difficulty of estimating advantage functions corresponding to policies from this class, is represented in the final bound by $\beta_K$, $G$, and $\lambda_{\min}^+(\bLambda_\nu)$. Increasing the number of denoising steps can make the policy class more expressive, but does not by itself imply that $\epsilon_{\text{bias}}$ decreases, since this error concerns approximation of the advantage by policy-score features. 

\paragraph{Why not a $\sqrt{T}$ rate for policy optimization?} Our NPG analysis balances an optimization term of order $T/\sqrt L$ against an estimation term of order $T/\sqrt N$, with $T=LN$. Balancing these terms gives $N=L=\sqrt T$ and the $T^{3/4}$ rate. This analysis does not rule out a $\sqrt T$ rate for other policy optimization methods, including methods with exact NPG updates that only estimate the advantage function. This occurs in e.g. the softmax parameterization case \citep{zhong2023theoreticalanalysisoptimisticproximal, liu2023optimisticnaturalpolicygradient, tan2025actorcriticsachieveoptimalsample, cassel2024warmupfreepolicyoptimization, sherman2024rateoptimalpolicyoptimizationlinear}, though only the last three achieve $\sqrt{T}$ regret.  

\paragraph{Related diffusion-policy analyses.} \citet{vu2026expressivity} study the trade-off between expressivity and statistical estimation for continuous-time SDE action samplers in discounted MDPs. They first show how increasing the allowed drift Lipschitz constant lets the diffusion concentrate around an optimal deterministic policy, with a value-approximation bound that decreases inversely with that constant. For neural drifts, however, this flexibility also increases estimation complexity. They balance these effects by choosing the drift Lipschitz budget and ReLU class size according to the sample size. Their empirical objective evaluates candidate policies using the same sampled initial states and noise sequences for action sampling and MDP transitions. Each candidate therefore induces different trajectories from the same random scenarios, whose number is their sample size. A simulator that replays the action and transition randomness across policies would make these evaluations available, although the paper does not explicitly specify this access model. The finite-sample theorem concerns a global empirical return maximizer under this shared randomness, rather than the policy-gradient algorithm used in their experiments. Their bound assumes a Lipschitz common-noise representation of the transition kernel, a bounded Lipschitz optimal deterministic policy, and a quadratic Bellman-gap bound. \citet{jiao2026sfac} take a more algorithm-specific approach, analyzing a Schr\"odinger--F\"ollmer actor--critic algorithm. They combine Bellman regression with an exponentially tilted policy target, then fit a neural drift and discretize the diffusion to sample from that target. Under their coverage and regularity conditions, they bound the average value gap across policy iterates, with a separate regression sample budget for each update. To account for the cost of learning during interaction, our analysis instead counts regret over sampled rounds or episodes, charging policy-gradient updates for their batch size. We handle exploration through coverage conditions or critic optimism under the respective reset, mixing, and hybrid access conditions.

\paragraph{Diffusion policy learning methods.} DIME \citep{celik2025dime} develops an entropy-lower-bound objective and analyzes exact policy iteration with finite action spaces. BDPO \citep{gao2025bdpo} studies behavior regularization through pathwise KL penalties and exact policy improvement. Reweighted score matching enables diffusion policy mirror descent and soft actor--critic updates \citep{ma2025efficient}, while DoobPO uses Doob's $h$-transform and conditional drift matching to fit reweighted diffusion and flow policies \citep{su2026doob}. These constructions are relevant to our critic-reweighted pretrained target, which retains the critic-cloning regret decomposition and gives an explicit rejection sampler. Appendix~\ref{app:gibbs-neural-fitting} also gives a finite-state neural training construction using fresh synthetic samples and exhaustive optimization. A guarantee for practical replay-based finetuning remains open.

\paragraph{Diffusion policies for imitation learning.} \citet{block2023provable} give finite-sample guarantees for matching expert trajectory marginals using diffusion behavior cloning with stabilizing low-level controllers. \citet{cheng2026multiobjective} analyze a semi-supervised specialist-to-generalist construction and bound the learned policy's value gap relative to the experts using demonstrations and specialist rollouts.

\paragraph{Offline data and statistical efficiency.} We use an offline dataset to bypass the need for optimism \citep{song2023hybrid} or a good reset distribution \citep{zhou2023offlinedataenhancedonpolicy}, under the respective Bellman-error coverage or hybrid NPG conditions. Having access to an offline dataset with good coverage may also allow one to achieve sharper regret guarantees as in \citep{li2023reward, tan2024hybridreinforcementlearningbreaks}. 

\paragraph{Comparison between policy gradient methods and behavior cloning methods.} The policy gradient methods and behavior cloning methods make use of two different aspects of diffusion models. 
We discuss their discrepancy from both algorithmic and analytical perspectives as follows.
\begin{itemize}
    \item The policy gradient methods take the diffusion model as an implicit, hierarchical probabilistic model. Since the model is updated via policy gradient instead of denoising score matching, it may not recover the time-reverse of the forward process in general scenarios. However, policy gradients can still improve the policy by learning a vector field that produces rewarding actions, even when that field differs from the score-based flows in standard diffusion models. 
    This can be reflected in our results, e.g., Theorem~\ref{thm:vanilla-pg-regret} and \ref{thm:npg-regret}, in which no denoising-score estimation error terms are in the regret bounds.
    
    \item The behavior cloning methods focus on the sampling ability of the diffusion model. In each iteration, the behavior cloning is done by the denoising score matching procedure. Thus, accurate score matching allows the sampler to approximate the current greedy target, with initialization and sampling errors also contributing to the approximation. 
    That is the reason of score estimation error popping out in our analysis and final regret bounds. 
\end{itemize}

\paragraph{Distributional expressivity examples.}\label{app:representation}
A stochastic policy can assign mass to several good actions without placing much mass between them. A single Gaussian fit can lose this distinction: fitting two separated rewarding regions increases its variance and spends most of its probability away from either region. We first quantify the resulting reward loss, even with an exact population fit. We then consider many independent two-mode choices, where a compact denoising actor avoids the exponential component count of a flat Gaussian mixture, and show how to learn the unknown choice probabilities from demonstrations. Throughout, $\TV(P,Q)=\sup_B|P(B)-Q(B)|$.

\begin{prop}[Single-Gaussian behavior cloning]\label{prop:gaussian-bc-example}
Consider a one-state bandit with actions in $\R$, parameters $m,\sigma_b>0$, demonstration law
\[
 \beta=\tfrac12\mathcal N(-m,\sigma_b^2)+\tfrac12\mathcal N(m,\sigma_b^2),
 \qquad
 r(a)=\tfrac12\left(e^{-(a-m)^2/(2\sigma_b^2)}+e^{-(a+m)^2/(2\sigma_b^2)}\right).
\]
The population maximum-likelihood Gaussian is $G=\mathcal N(0,m^2+\sigma_b^2)$, and
\begin{align}
 \mathbb E_\beta r
 &=\frac{1+e^{-m^2/\sigma_b^2}}{2\sqrt{2}},\label{eq:bc-demo-value}\\
 \mathbb E_G r
 &=\frac{\sigma_b}{\sqrt{m^2+2\sigma_b^2}}
   \exp\!\left(-\frac{m^2}{2(m^2+2\sigma_b^2)}\right).
   \label{eq:bc-gaussian-value}
\end{align}
Consequently, $\mathbb E_G r\to0$ as $m/\sigma_b\to\infty$, whereas every policy $\pi$ with $\TV(\pi,\beta)\leq\varepsilon$ satisfies $\mathbb E_\pi r\geq1/(2\sqrt{2})-\varepsilon$.
\end{prop}
\begin{proof}
For a candidate $\mathcal N(\mu,v)$, maximizing the expected log density is equivalent to minimizing
\[
 \tfrac12\log v+\frac{\mathbb E_\beta(A-\mu)^2}{2v}.
\]
The minimizers are $\mu=\mathbb E_\beta A=0$ and $v=\mathbb E_\beta A^2=m^2+\sigma_b^2$. To compute their returns, for $Y\sim\mathcal N(\mu,v)$ and $w>0$, completing the square gives
\[
 \mathbb E\exp\!\left(-\frac{(Y-c)^2}{2w^2}\right)
 =\frac{w}{\sqrt{w^2+v}}
  \exp\!\left(-\frac{(\mu-c)^2}{2(w^2+v)}\right).
\]
Applying this identity to each demonstration component and to $G$ proves~\eqref{eq:bc-demo-value}--\eqref{eq:bc-gaussian-value}. Finally, $0\leq r\leq1$, so $|\mathbb E_\pi r-\mathbb E_\beta r|\leq\TV(\pi,\beta)$.
\end{proof}

The loss comes from trying to clone both rewarding modes with one Gaussian. If we instead optimize reward directly, a Gaussian centered at $m$ with vanishing variance achieves return approaching $r(m)\geq1/2$. A two-component mixture also clones $\beta$ exactly. The more substantial representational question is therefore how the required number of components grows when each action coordinate has two possible modes. The next construction has $2^d$ combinations, but its denoising map handles the coordinates separately.

\begin{prop}[A comparison with flat Gaussian mixtures]\label{prop:flat-gmm-separation}
For an integer $d\geq1$, let
\[
 r_d=\frac{1}{100d},\qquad
 \sigma_d=\frac{r_d}{\sqrt{2\log(2000d)}},\qquad
 P_d=\bigotimes_{j=1}^d
 \left[\tfrac12\mathcal N(-1,\sigma_d^2)+\tfrac12\mathcal N(1,\sigma_d^2)\right].
\]
Every mixture of $M$ Gaussian distributions on $\R^d$, with arbitrary covariance matrices, satisfying $\TV(Q,P_d)\leq0.01$ has $M\geq0.79\,2^d$.
In contrast, for independent $Z,\xi\sim\mathcal N(0,I_d)$, the smooth denoising layer
\[
 A=\tanh(\kappa_d Z)+\sigma_d\xi,
 \qquad \kappa_d=100d/\sigma_d,
\]
has $\TV(\mathcal L(A),P_d)<0.003$ and uses $O(d)$ scalar $\tanh$ units and nonzero parameters. Its weights and each fixed positive-order spatial derivative grow at most polynomially in $d$.
\end{prop}

\begin{proof}
For each $v\in\{-1,1\}^d$, define the box
\[
 B_v=\prod_{i=1}^d[v_i-r_d,v_i+r_d],
 \qquad U=\bigcup_{v\in\{-1,1\}^d}B_v.
\]
We first show that a Gaussian component placing substantial mass in $U$ must concentrate near one vertex. This lets us count how many vertices a flat mixture can cover.

Write a sample from $P_d$ as $S+\sigma_d\xi$, where the coordinates of $S$ are independent uniform signs and $\xi$ is standard Gaussian. The Gaussian tail bound and a union bound imply
\[
 P_d(U)\geq1-2d\exp\!\left(-\frac{r_d^2}{2\sigma_d^2}\right)=0.999.
\]
Suppose $Q=\sum_{j=1}^M w_jQ_j$ and $\TV(P_d,Q)\leq0.01$, so $Q(U)\geq0.989$. Call component $j$ good when $Q_j(U)\geq0.9$. If the good components have total weight $w_{\mathrm{good}}$, then
\[
 0.989\leq Q(U)\leq w_{\mathrm{good}}+0.9(1-w_{\mathrm{good}}),
 \qquad w_{\mathrm{good}}\geq0.89.
\]

Fix a good component with coordinate means $\mu_i$ and standard deviations $s_i$. Each marginal assigns probability at least $0.9$ to
\[
 I=[-1-r_d,-1+r_d]\cup[1-r_d,1+r_d].
\]
When $s_i>0$, its density is at most $(s_i\sqrt{2\pi})^{-1}$, giving $0.9\leq4r_d/(s_i\sqrt{2\pi})$ and hence $s_i<2r_d$. The bound $s_i\leq2r_d$ also holds for a degenerate marginal. If $\mu_i$ were at distance at least $1/4$ from both signs, every point in $I$ would be at distance at least $1/4-r_d\geq0.24$ from $\mu_i$. Chebyshev's inequality would then give marginal probability at most
\[
 \frac{s_i^2}{0.24^2}\leq\frac{4r_d^2}{0.24^2}<0.007,
\]
contradicting the lower bound $0.9$. Thus $\mu_i$ lies within $1/4$ of a unique sign $v_i$. The resulting vertex $v$ is defined coordinatewise, without any independence assumption on the component.

A point of $U\setminus B_v$ lies in an opposite-sign interval in at least one coordinate, at distance at least $7/4-r_d\geq1.74$ from that coordinate's mean. Another union bound and Chebyshev's inequality give
\[
 Q_j(U\setminus B_v)\leq\frac{4dr_d^2}{1.74^2}<0.001,
 \qquad Q_j(B_v)\geq0.899.
\]
Let $U_\mathrm{sel}$ be the union of the boxes assigned to the good components. There are at most $M$ such boxes, and
\[
 Q(U_\mathrm{sel})\geq0.89\cdot0.899=0.80011,
 \qquad P_d(U_\mathrm{sel})\geq0.79011.
\]
The boxes are disjoint and sign symmetry gives $P_d(B_v)=P_d(U)/2^d\leq2^{-d}$. Therefore $M2^{-d}\geq0.79011$, which proves the stated lower bound, including when the covariance matrices are singular.

For the denoising construction, couple the target using $S=\operatorname{sign}(Z)$, so $S+\sigma_d\xi\sim P_d$. Conditional on $Z$, the target and constructed laws are Gaussians with common covariance $\sigma_d^2I_d$. Their Gaussian KL divergence, Pinsker's inequality, and convexity of total variation under common mixing yield
\begin{align*}
 \TV(\mathcal L(A),P_d)
 &\leq\frac{\mathbb E\|\operatorname{sign}(Z)-\tanh(\kappa_dZ)\|_2}{2\sigma_d}\\
 &\leq\frac{d}{2\sigma_d}\mathbb E[1-\tanh(\kappa_d|Z_1|)]\\
 &\leq\frac{d\sqrt{2/\pi}}{2\sigma_d\kappa_d}
       \int_0^\infty(1-\tanh u)\,\mathrm du
 =\frac{\sqrt{2/\pi}\log2}{200}<0.003.
\end{align*}
The second line uses $\|x\|_2\leq\|x\|_1$, and the third bounds the half-normal density by $\sqrt{2/\pi}$. The coordinatewise map uses $d$ scalar $\tanh$ units with diagonal weights and $O(d)$ nonzero parameters. Since $\kappa_d=10^4d^2\sqrt{2\log(2000d)}$, its weights and every fixed positive-order spatial derivative are bounded by a polynomial in $d$. These bounds are dimension-dependent, rather than uniform smoothness bounds.
\end{proof}

To relate the construction to our Gaussian denoising actor, fix $\alpha\in(0,1)$ and $\eta\ne0$, and choose
\[
 \epsilon(z)=\frac{\sqrt\alpha\tanh(\kappa_dz)-z}{\eta},
 \qquad \sigma_{\mathrm{in}}=\sqrt\alpha\,\sigma_d.
\]
Then $\alpha^{-1/2}(Z+\eta\epsilon(Z)+\sigma_{\mathrm{in}}\xi)=A$. This uses freely chosen coefficients for one denoising layer. It is not a guarantee for a prescribed reverse-diffusion schedule or for an actor satisfying all the regularity and conditioning assumptions of our optimization theorems.

The same product structure also gives a compact spatial score for the noised target, connecting this construction to score-based diffusion sampling. For $A\sim P_d$ and an independent standard Gaussian $Z$, let
\[
 X_\alpha=\sqrt\alpha A+\sqrt{1-\alpha}Z,
 \qquad v_\alpha=\alpha\sigma_d^2+1-\alpha,
 \qquad \alpha\in[0,1].
\]
Its density $p_\alpha$ is a product of symmetric two-component mixtures with coordinate means $\pm\sqrt\alpha$ and common variance $v_\alpha$. The one-coordinate density is proportional to
\[
 \exp\!\left(-\frac{x_i^2+\alpha}{2v_\alpha}\right)
 \cosh\!\left(\frac{\sqrt\alpha x_i}{v_\alpha}\right),
\]
so direct differentiation gives
\begin{equation}\label{eq:product-mixture-spatial-score}
 [\nabla_x\log p_\alpha(x)]_i
 =\frac{-x_i+\sqrt\alpha\tanh(\sqrt\alpha x_i/v_\alpha)}{v_\alpha}.
\end{equation}
At each fixed $\alpha$, this exact score uses $d$ scalar $\tanh$ units and a linear term. Since $v_\alpha\geq\sigma_d^2$, the coefficients and each fixed positive-order spatial derivative admit bounds polynomial in $d$, uniformly over $\alpha\in[0,1]$. The score itself has linear growth in $x$. Thus the exponential number of modes does not force an exponential score representation either. Learning the score or controlling a prescribed finite reverse-sampling schedule requires additional analysis.

Finally, both examples are on $\R^d$. Applying the same measurable bijection from $\R^d$ onto a bounded open action domain preserves total variation, but changes the comparison class to the pushforwards of the latent Gaussian mixtures. It does not prove the same lower bound for ordinary Gaussian mixtures defined directly in a bounded action domain. The compact product-mixture description of $P_d$ also shows why the exponential lower bound specifically concerns flat mixtures.

So far the construction specifies the target distribution in advance. To see whether its compact structure also helps learning, we now let the probability of each coordinate's two modes be unknown. An observed coordinate's sign estimates that probability after a simple correction for Gaussian noise. The resulting actor learns from the same demonstrations used by a Gaussian-mixture fit.

\begin{prop}[Learning the product mixture from demonstrations]
\label{prop:learned-product-mixture}
Use $r_d,\sigma_d$ from Proposition~\ref{prop:flat-gmm-separation}, and
let the unknown demonstration law be
\[
 P_p=\bigotimes_{j=1}^d\left[(1-p_j)\mathcal N(-1,\sigma_d^2)
                      +p_j\mathcal N(1,\sigma_d^2)\right],
 \qquad p_j\in[1/4,3/4].
\]
Let $\Phi$ denote the standard Gaussian distribution function and set
$h_d=2\Phi(1/\sigma_d)-1$. From $n$ independent demonstrations
$A^{(1)},\ldots,A^{(n)}$, form
\[
 \widehat p_j=\operatorname{clip}_{[1/4,3/4]}
 \left(\frac{n^{-1}\sum_{i=1}^n\mathbf1\{A_j^{(i)}\geq0\}
                   -(1-h_d)/2}{h_d}\right),\qquad
 t_j=\Phi^{-1}(1-\widehat p_j).
\]
For $\varepsilon,\delta\in(0,1)$, if
\[
 n\geq\frac{16d\log(2d/\delta)}{3h_d^2\varepsilon^2},\qquad
 \kappa=\frac{d\sqrt{2/\pi}\log2}{\sigma_d\varepsilon},
\]
the denoising actor
\[
 \widehat A_j=\tanh(\kappa(Z_j-t_j))+\sigma_d\xi_j,
 \qquad Z,\xi\stackrel{\rm iid}{\sim}\mathcal N(0,I_d),
\]
satisfies $\TV(\mathcal L(\widehat A),P_p)\leq\varepsilon$ with
probability at least $1-\delta$ over its training demonstrations.
It uses $O(d)$ scalar tanh units and nonzero parameters, with weights
polynomial in $d$ and $\varepsilon^{-1}$. Every flat mixture of $M$
Gaussians whose total variation error is at most $0.01$ instead satisfies
$M\geq0.79(4/3)^d$, regardless of how many demonstrations it receives.
\end{prop}
\begin{proof}
The sign indicator has mean $(1-h_d)/2+h_dp_j$. Hoeffding's inequality,
a union bound, and nonexpansiveness of clipping give, with probability
at least $1-\delta$,
\[
 |\widehat p_j-p_j|\leq
 u_n:=\sqrt{\frac{\log(2d/\delta)}{2nh_d^2}}
 \quad\text{for every }j.
\]
Generate the mixture using independent Bernoulli signs, then add the
same Gaussian noise kernel. Data processing, the elementary Bernoulli
KL bound, and Pinsker's inequality imply
\begin{align*}
 \KL(P_p\|P_{\widehat p})
 &\leq\sum_{j=1}^d\KL(\operatorname{Ber}(p_j)\|
                              \operatorname{Ber}(\widehat p_j))
 \leq\frac{16}{3}\sum_{j=1}^d(p_j-\widehat p_j)^2,\\
 \TV(P_p,P_{\widehat p})
 &\leq\sqrt{\frac{8d}{3}}u_n\leq\varepsilon/2.
\end{align*}
The sign $\operatorname{sign}(Z_j-t_j)$ is positive with probability
$\widehat p_j$. Couple these signs and the smooth actor using the same
$Z,\xi$. Conditional Gaussian KL and Pinsker, followed by the uniform
upper bound on the normal density, give
\[
\begin{aligned}
 \TV(\mathcal L(\widehat A),P_{\widehat p})
 &\leq\frac{1}{2\sigma_d}\sum_{j=1}^d
   \mathbb E[1-\tanh(\kappa|Z_j-t_j|)]\\
 &\leq\frac{d\sqrt{2/\pi}\log2}{2\sigma_d\kappa}
 \leq\varepsilon/2.
\end{aligned}
\]
The triangle inequality proves the learning guarantee. Since
$|t_j|\leq\Phi^{-1}(3/4)$ and
$\sigma_d^{-1}=100d\sqrt{2\log(2000d)}$, the parameter bounds follow.
Also $h_d\geq1/2$, for example by the Gaussian tail bound, so the
sample requirement is $O(d\log(d/\delta)/\varepsilon^2)$.

For the lower bound, retain the boxes $B_v$ and $U$ from
Proposition~\ref{prop:flat-gmm-separation}. Their union still has
$P_p(U)\geq0.999$, so that proof gives a union $U_{\rm sel}$ of at most
$M$ boxes with $P_p(U_{\rm sel})\geq0.79011$. Each one-dimensional
mixture assigns probability at most $3/4$ to either interval
$[\pm1-r_d,\pm1+r_d]$. Indeed, if $c$ and $o$ are the Gaussian
probabilities of the centered and opposite intervals, respectively,
then $c+o\leq1$ and $c\geq o$, so the mixture probability is at most
$(3/4)c+(1/4)o\leq3/4$. Independence now gives
$P_p(B_v)\leq(3/4)^d$. Thus
$M(3/4)^d\geq0.79011$, proving the claim.
\end{proof}

The product structure makes the statistical problem depend on $d$ unknown probabilities, even though the distribution has $2^d$ possible mode combinations. This gives a finite-sample imitation separation under the same demonstration access and total variation loss. The estimator uses the known product structure and sign statistics, rather than denoising-score-matching optimization, and the comparison concerns distribution matching rather than unregularized reward maximization. It also supplies explicit scores for the estimated law: after Gaussian noising at level $\alpha$, let
$v_\alpha=\alpha\sigma_d^2+1-\alpha$. Coordinate $j$ has score
\[
 \partial_{x_j}\log p_{\alpha,p}(x)
 =\frac{-x_j+\sqrt\alpha\tanh\left(
       \frac{\sqrt\alpha x_j}{v_\alpha}
       +\frac12\log\frac{p_j}{1-p_j}\right)}{v_\alpha}.
\]
Replacing $p_j$ by $\widehat p_j$ therefore gives the exact score of the
estimated product law, using one biased tanh unit per coordinate.

\subsection{Coverage examples}

\paragraph{A finite-state coverage example.}\label{app:finite-state-coverage}
Assumption~\ref{aspt:npg-concentrability} requires the sampling law to put mass wherever the comparator does. State resets alone do not ensure this, because the action distributions can still have disjoint supports. A density floor gives a direct sufficient condition. Let
$\gS$ be finite, let the reset law satisfy $\nu(s)\ge c>0$, and let
$\gA$ have finite positive volume. Suppose the current policies and the
comparator have densities, with
$\pi_\theta(a\mid s)\ge m>0$ and
$\pi_{\rm cmp}(a\mid s)\le M<\infty$ almost everywhere.
The time-zero term of the discounted occupancy gives
\[
 d_\nu^{\pi_\theta}(s,a)\ge(1-\gamma)\nu(s)\pi_\theta(a\mid s)
 \ge(1-\gamma)cm.
\]
Since $d_{\rho,S}^{\pi_{\rm cmp}}(s)\le1$, its joint density is at most $M$.
Consequently,
\[
 C\le\frac{M}{(1-\gamma)cm}.
\]
One way to obtain the density floor is to include a fixed uniform component:
if $\pi_\theta=(1-\zeta)p_\theta+\zeta\operatorname{Unif}(\gA)$,
$0<\zeta\le1$, and $p_\theta\le M_0$, then
$m=\zeta/\operatorname{vol}(\gA)$ and
$M=(1-\zeta)M_0+m$ work for a comparator in the same class.
This calculation applies to the displayed density-bounded policy and comparator classes.

Coverage can fail even with one state. On $\gA=[-1,1]$, a sampling law
uniform on $[-b,b]$, $b<1/2$, assigns no mass to the support of a comparator
uniform on $[1/2,1]$, so its concentrability coefficient is infinite.
A density floor avoids this failure, but its magnitude matters: adding
uniform exploration of mass $\zeta$ to that narrow law gives a ratio
$4/\zeta$ on $[1/2,1]$. For hybrid sampling, the ratio is measured against
the full online/offline mixture, so online mass can cover regions absent
from the logged data. Neither coverage calculation implies Fisher
conditioning or compatible advantage approximation.









% A baseline does not help too much here, as it is simply variance-reducing and the issue is in the estimation of the Q-function. With on-policy Monte-Carlo sampling for $N$ steps, we can only guarantee that the error in estimating the squared loss scales with $\epsilon_{\text{stat}} = O(1/\sqrt{N})$. If one uses the least squares estimator, we can achieve a $O(d/N)$ rate, corresponding to $T^{3/4}$ regret.  

% \begin{aspt}[Policy Smoothness]
%     $\log \pi_\theta$ is $\beta$-smooth in $\theta$. That is, $||\log \pi_\theta(a|s) - \log \pi_{\theta'}(a|s)||_2 \leq \beta ||\theta - \theta'||_2$ for all $\theta, \theta' \in \Theta$.
% \end{aspt}
% \kevin{Ideally this is not an assumption. Work it out in terms of the score function.}


% \zh{Reverse the fake DDPM backward process. Can we find a new forward process for simulation? Or analyze the Doob's h-tranformation?}

% The main technical novelty now lies in our analysis of the NPG update direction estimate resulting from running $K$ denoising iterations. 

\section{\aigen{Diffusion Policy Rollout Sampler}}


\begin{algorithm}[tbp]
    \caption{Monte Carlo labels under reset, continuing, or logged-trajectory access}
    \begin{algorithmic}[1]
        \STATE {\bfseries Input:} frozen policy $\pi_\theta$, discount $\gamma$, access mode, and its state or trajectory input.
        \IF{logged-trajectory mode}
            \STATE Read one unused behavior record $(S_t,A_t,R_t)_{t=0}^{H-1}$ with independent $\mathbb{P}(H>t)=\gamma^t$ and initial law $\mu_S\beta$.
            \STATE Set $W_0=1$, $W_t=\prod_{j=1}^t\pi_\theta(A_j\mid S_j)/\beta(A_j\mid S_j)$, and $\widehat Q=\sum_{t=0}^{H-1}W_tR_t$.
            \IF{$\operatorname{Ber}(1/2)=1$}
                \STATE \textbf{Return} feature pair $(S_0,A_0)$ and label $2\widehat Q$.
            \ELSE
                \STATE Draw $A'\sim\beta(\cdot\mid S_0)$ independently of the record given $S_0$.
                \STATE \textbf{Return} feature pair $(S_0,A')$ and label $-2[\pi_\theta(A_0\mid S_0)/\beta(A_0\mid S_0)]\widehat Q$.
            \ENDIF
        \ENDIF
        \IF{reset mode}
            \STATE Initialize $s\sim\nu$. Draw an action and its denoising path under $\pi_\theta$. Before execution, accept with probability $1-\gamma$. Otherwise execute, move to the next state, and repeat.
        \ELSE[continuing mode]
            \STATE Use the current state, after the batch's burn-in or the preceding complete return. Draw an accepted action under $\pi_\theta$ without an occupancy search.
        \ENDIF
        \STATE Retain the accepted pair $(s,a)$ and its denoising path $a_{1:K}$.
        \IF{$\operatorname{Ber}(1/2)=1$}
            \STATE Execute the accepted action and continue under $\pi_\theta$, stopping with probability $1-\gamma$ after each reward. Set $\widehat A$ to twice the sum of collected rewards.
        \ELSE
            \STATE Replace the accepted action by an independent draw from $\pi_\theta(\cdot\mid s)$ before execution. Run the same geometric return, and set $\widehat A$ to minus twice its reward sum.
        \ENDIF
        \STATE \textbf{Return} $\widehat A,s,a,a_{1:K}$ and retain the terminal state for the next continuing sample.
    \end{algorithmic}
\label{alg:dppo-sampler}
\end{algorithm}
The reset branch adapts Algorithm~3 of \citet{agarwal2021policygradient}. The continuing branch uses one burn-in per batch and no restart at a visited state. Its dependence is controlled in Appendix~\ref{app:vanilla-pg-mixing}. The logged branch requires the absolute-continuity, reserved-block, and moment conditions of Appendix~\ref{app:hybrid-reset-dist}. It uses stored continuations through the geometric stopping time, not isolated transition tuples. Both logged branches have feature law $\mu$, and their conditional mean label is the current advantage. The logged branch evaluates densities and endpoint features under the current DDIM actor.




\section{Computational Efficiency via Hybrid RL}
\label{app:hybrid}

% \begin{algorithm}[H]
%     \caption{\aigen{Diffusion Policy Behavior Cloning with Fitted Q-Iteration}}
%     \begin{algorithmic}[1]
%            \STATE {\bfseries Input:} Diffusion policy class $\left\{\pi_\theta \mid \theta \in \Theta \right\}$, denoising steps $K$. Optional offline dataset $\gD^{\off}$. Function class $\gF$. 
%             \STATE {\bfseries Initialize:} Dataset $\gD_h^{(0)} \gets \emptyset$ for all $h=1,...,H$, initial policy $\pi^{(1)}$.  
%             \FOR{episode $t = 1, 2, \dots, T$}
%             \STATE Execute $\pi^{(t)}$ via sampling from DDPM/DDIM for one episode, collect trajectory $\tau^{(t)}$.
%             \STATE Append trajectory $\tau^{(t)}$ to dataset $\gD_h^{(t)} \gets \gD_h^{(t-1)} \cup \left\{\tau^{(t)}\right\}$.
%         \STATE Update critic $f_h^{(t+1)} \in \argmin_{f_h \in \gF_h} \widehat{\mathcal{L}}_h^{(t)}\left(f_h, f_{h+1}\right)$ for $h=H-1,...,1$.
%         \STATE Construct behavior cloning dataset with new actions for every $i=1,...,t$:
%         $$\widetilde{\gD}_h^{(t+1)} \gets \left\{\left(s_h, \widetilde{a}_h\right)\right\}_{s_h \in \gD_h^{(t)} \cup \gD_h^{\off}}, \text{ where } \widetilde{a}_h \gets \pi_h^{f^{(t+1)}}(s_h), h=1,...,H.$$
%         \STATE Optimize $\pi_h^{(t+1)}$ by score-matching to clone $\pi_h^{f^{(t+1)}}$ on $\widetilde{\gD}_h^{(t+1)}$, where $ Z_k^{(i)} \sim \gN(\mathbf{0}, \mathbf{I})$:
%         $$\theta_h^{(t)} \in \argmin_\theta \sum_{k=1}^K \sum_{s_h, \widetilde{a}_h \in \widetilde{\gD}_h^{(t+1)}} \left\lVert Z^{(i)}_k - \epsilon_{\theta}\left(\sqrt{\bar{\alpha}_k}\widetilde{a}_h + \sqrt{1-\bar{\alpha}_k Z^{(i)}_k}, k, s_h\right)\right\rVert_2^2.$$
%         \ENDFOR
%     \end{algorithmic}
% \label{alg:dipo-fqi}
% \end{algorithm}

Notice that Algorithm \ref{alg:dipo-bc} with only online data can be computationally inefficient. At each episode, it searches over a confidence set of functions in order to implement optimism. As in \citet{song2023hybrid, amortila2024harnessing, zhou2023offlinedataenhancedonpolicy, tan2025actorcriticsachieveoptimalsample}, we show that the need for optimism can be bypassed if one has access to an offline dataset with good coverage. We take the offline data to be $N_{\off}$ independent episodes from a fixed behavior policy, with stage marginals $\mu_h$ and the same reward and transition kernels. Subsequent rollouts use fresh environment randomness conditional on the past data. We use measurable selections by the critic and greedy-action oracles, and treat exact regression and greedy maximization as optimization oracles. The quality of the coverage of the offline dataset can be measured by the following measure of offline RL complexity known as the single-policy concentrability coefficient below, of which we employ the Bellman error version:

\begin{defn}[Single-Policy Concentrability]
    We define the single-policy concentrability coefficient \citep{zhan2022offline} by
    \begin{align*}
    c_{\off}^*(\gF) = \max_{h\in[H]}\sup_{f \in \gF} \frac{\E_{\pi^*}\left[f_h-\gT_h f_{h+1} \right]^2}{\E_{\mu_h}[\left(f_h-\gT_h f_{h+1}\right)^2]}.
\end{align*}
A zero denominator gives ratio zero if its numerator is zero and infinity otherwise.
\label{defn:spc}
\end{defn}

We therefore modify the algorithm of \citet{song2023hybrid} to incorporate a diffusion policy, presented in the offline-data branch within Algorithm \ref{alg:dipo-bc}. This amounts to fitting the critic backwards by TD-error minimization in each round while augmenting online data with an offline dataset. This avoids searching over a confidence set, and can be thought of as fitted-Q iteration (FQI) with a diffusion exploration policy learned by behavior cloning. With exact regression and greedy maximization, Algorithm \ref{alg:dipo-bc} achieves a similar regret to the online-only case:
\begin{thm}[Hybrid Regret Bound for Algorithm \ref{alg:dipo-bc}]
Under the assumptions of Theorem~\ref{thm:dipo-bc-regret} and the preceding offline sampling model, let $N_{\off}\ge1$, $n=N_{\off}+T$, and $\beta=C\log(2H(n+1)\gN_\gF(1/(n+1))^2/\delta)$. With probability at least $1-\delta$, the backward-FQI branch satisfies
\[
 \mathrm{Reg}(T)\le O\!\left(H^2\sqrt{\beta T\mathsf{SEC}(\gF,\Pi,T)}
                +H^2T\sqrt{\frac{\beta c_{\off}^*(\gF)}{N_{\off}}}\right)
                  +2H\sum_{t=1}^T B_t.
\]
Here $B_t$ is the population cloning loss of the hybrid iterates. Lemma~\ref{lem:behavior-cloning-error} gives $B_t\le B_t^{\mathsf{sampler}}$ under Assumption~\ref{aspt:score-matching-error}.
\end{thm}

We defer the proof to Appendix \ref{app:dipo-fqi}. For $N_{\off}=\Theta(T)$, the additional critic term is $O(H^2\sqrt{\beta c_{\off}^*(\gF)T})$. It is $\widetilde O(\sqrt T)$ when $H$ and the coverage coefficient are fixed and $\beta$ grows at most polylogarithmically in $T$. This demonstrates that having access to offline data with good coverage, i.e. to say a bounded single-policy concentrability coefficient $c_{\off}^*(\gF)$, allows one to omit the use of optimism. This replaces the optimistic search by regression with the displayed additional regret. 


\section{The Policy Score for Diffusion Policies}
\label{app:lipschitz}

\paragraph{Policy gradients and denoising scores.}\label{app:nonconservative-score}
For a one-state problem, take $Z\sim N(0,I_2)$ and a one-step DDIM map
\[
 F_\theta(z)=\big((1+\theta)z_1+\theta z_2,z_2\big),\qquad
 r(a)=\tfrac12(1+\cos a_1).
\]
Its expected reward is $J(\theta)=\tfrac12[1+\exp(-((1+\theta)^2+\theta^2)/2)]$, so $J'(0)=-e^{-1/2}/2$. Thus a sufficiently small gradient-ascent step from zero improves reward and gives $\theta\ne0$.
To realize this map in Equation~(\ref{eq:DDIM}), use $\epsilon_\theta(z)=(\sqrt\alpha F_\theta(z)-z)/\eta$ for fixed $0<\alpha<1$ and $\eta\ne0$. At zero this field is proportional to a Gaussian spatial score. After the update, its cross-partials are $\partial_{z_2}\epsilon_{\theta,1}=\sqrt\alpha\theta/\eta$ and $\partial_{z_1}\epsilon_{\theta,2}=0$, so it cannot be a scaled spatial log-density gradient. Nevertheless, $F_\theta$ remains invertible when $\theta\ne-1$, and its induced policy has a valid parameter score. This is why our PG argument controls the policy score directly, without requiring score matching to remain accurate during reward optimization.



\begin{lem}[DDPM and DDIM Gradient of Log-Policy Densities]\label{lem:ddpm-ddim-log-policy-densities}
For DDPM with positive transition variances,
\begin{align*}
\widehat\bphi^\theta(s,a_{0:K})
=\sum_{k=1}^K\frac{\eta_k}{\sigma_k}
[\nabla_\theta\epsilon_\theta(a_k,k,s)]^\top Z_k,
\qquad
\E[\widehat\bphi^\theta\mid s,a_0]=\bphi^\theta(s,a_0).
\end{align*}
For DDIM, fix $s$ and write $z=a_K$,
$\bJ_\theta(z)=\nabla_z\Phi_{K:0}^\theta(z)$ and
$w_i(z)=\bJ_\theta(z)^{-1}\partial_{\theta_i}\Phi_{K:0}^\theta(z)$.
The $i$-th component of the endpoint score, differentiating at fixed $a_0$, is
\begin{align*}
\bphi_i^\theta(s,a_0)
=z^\top w_i(z)-\operatorname{tr}(\bJ_\theta^{-1}\partial_{\theta_i}\bJ_\theta)
+(\nabla_z\log|\det\bJ_\theta|)^\top w_i(z).
\end{align*}
All quantities on the right are evaluated at $z=a_K$, and the displayed
$\partial_{\theta_i}\bJ_\theta$ holds $z$ fixed.
\end{lem}

\begin{proof}
    This is a direct combination of Lemmas \ref{lem:features-ddpm} and \ref{lem:feature-concentration-ddim}.
\end{proof}



% \yuting{move the lemma to the appendix}
\begin{lem}[Policy Score Moments of Diffusion Policies]
For DDPM, if $\|\nabla_\theta\epsilon_\theta(a,k,s)\|_{\mathrm{op}}\leq L_\theta$ uniformly, one may take
\[
(\beta_K^{\mathsf{ddpm}})^2=d_A L_\theta^2
\sum_{k=1}^K\eta_k^2/\sigma_k^2.
\]
For DDIM, let $W_\theta(z)$ have columns
$w_i(z)=\bJ_\theta(z)^{-1}\partial_{\theta_i}\Phi_{K:0}^\theta(z)$ and put
$q_i(z)=\operatorname{div}_z w_i(z)$.
If, uniformly in $\theta,s,z$, $\|W_\theta(z)\|_{\mathrm{op}}\leq U_K$
and $\|q(z)\|_2\leq V_K$, then one may take
\[
(\beta_K^{\mathsf{ddim}})^2=2U_K^2d_A+2V_K^2,
\qquad
\E[\|\bphi^\theta\|^4\mid s]
\leq8U_K^4d_A(d_A+2)+8V_K^4.
\]
The quantities $U_K$ and $V_K$ control the inverse-map sensitivity and its
spatial variation, which together determine the DDIM score moments.
\label{lem:lipschitz-constant-policy}
\end{lem}
% \vspace{-4ex} 
% \subsection{Proof of Lemma \ref{lem:lipschitz-constant-policy}}
\begin{proof}
    Fix a state $s\in \gS$. Consider a sequence $a_K,...,a_0$ from either the DDPM or DDIM samplers. Recall that these are generated by
    $$a_{k-1} = \frac{1}{\sqrt{\alpha_k}}(a_k + \eta_k^{\mathsf{ddim}} \epsilon_\theta(a_k, k, s)) = \Phi_k^\theta(a_k)$$
    for the DDIM sampler, and
    $$a_{k-1} = \frac{1}{\sqrt{\alpha_k}}(a_k + \eta_k^{\mathsf{ddpm}} \epsilon_\theta(a_k, k, s) + \sigma_k^{\mathsf{ddpm}}Z_k) = \Psi_k^\theta(a_k, Z_k)$$
    for the DDPM sampler.
    
    For the following sufficient sensitivity estimates, suppose the denoiser
is $(L_a,L_\theta)$-Lipschitz in $(a,\theta)$, uniformly in $s$ and $k$.
These derivative bounds are not additional premises when the score moments
in Assumption~\ref{aspt:ddpm-bijection} are verified directly.
First, the recursion below holds for both the DDPM and DDIM samplers:
    \begin{equation}
        \nabla_\theta a_{k-1} = \frac{1}{\sqrt{\alpha_k}}\left(\nabla_\theta a_{k} + \eta_k \nabla_\theta \epsilon_\theta(a_k, k, s) + \eta_k\nabla_a \epsilon_\theta(a_k, k, s) \nabla_\theta a_{k} \right),
        \label{eq:lipschitz-gradient-recursion-1}
    \end{equation}
    noting that the Gaussian noise has derivative zero by the reparametrization trick, and the second term arises from the chain rule. Taking norms on both sides, we see that
    \begin{equation}
        \lVert\nabla_\theta a_{k-1}\rVert_2 \leq \frac{1+|\eta_k| L_a}{\sqrt{\alpha_k}}\lVert\nabla_\theta a_{k}\rVert_2 + \frac{|\eta_k| L_\theta}{\sqrt{\alpha_k}}.
        \label{eq:lipschitz-norm-recursion-2}
    \end{equation}
    Using the fact that $a_K \sim \gN(0, \mathbf{I}_d) \implies \nabla_\theta a_K = 0$ to unroll the recursion in Equation \ref{eq:lipschitz-norm-recursion-2}, we have
    \begin{equation}
        \lVert\nabla_\theta a_0 \rVert_2 \leq L_\theta \sum_{k=1}^K \frac{|\eta_k|}{\sqrt{\alpha_k}} \prod_{i=1}^{k-1} 
        \frac{1+|\eta_i| L_a}{\sqrt{\alpha_i}}.
        \label{eq:lipschitz-unrolled-recursion-3}
    \end{equation}

    \paragraph{DDPM.}
We now address the DDPM sampler. We first bound the realized pathwise score estimator
\[
    \widehat{\bphi}^{\theta}(s,a_{0:K})
    :=
    \sum_{k=1}^K \nabla_\theta \log \pi_\theta(a_{k-1}\mid a_k,s),
\]
rather than a uniform bound over all possible denoising paths. Let
\(\mathcal I\) denote the set of DDPM pathwise score evaluations used by the
algorithm, and write \(M:=|\mathcal I|\). For Algorithm \ref{alg:dppo-vanilla-pg}, one may take
\(M=T\), or more generally \(M\) equal to the total number of pathwise score
queries made in the proof.

For the \(m\)-th such query, let \(Z_{m,k}\) be the Gaussian noise used at
denoising step \(k\). Define
\[
    R_{M,K,\delta}
    :=
    \sqrt{d_A + 2\sqrt{d_A \log(MK/\delta)} + 2\log(MK/\delta)} .
\]
By the standard \(\chi^2\) tail bound and a union bound over
\((m,k)\in \mathcal I\times[K]\), with probability at least \(1-\delta\),
\[
    \max_{m\in\mathcal I}\max_{k\in[K]}\|Z_{m,k}\|_2
    \le R_{M,K,\delta}.
\]
This event is over the realized Gaussian noises only, and does not involve any
supremum over unobserved denoising paths.

For a DDPM transition, writing
\[
    \mu_{\theta,k}(a_k,s)
    :=
    \frac{1}{\sqrt{\alpha_k}}
    \left(a_k+\eta_k\epsilon_\theta(a_k,k,s)\right),
\]
we have
\[
    a_{k-1}\mid a_k,s
    \sim
    \mathcal N\left(
        \mu_{\theta,k}(a_k,s),
        \frac{\sigma_k^2}{\alpha_k}I_{d_A}
    \right).
\]
Hence, for the conditional likelihood score, where \(a_{k-1}\) and \(a_k\) are
held fixed when differentiating,
\[
\begin{aligned}
    \nabla_\theta \log \pi_\theta(a_{k-1}\mid a_k,s)
    &=
    \frac{\alpha_k}{\sigma_k^2}
    \left[\nabla_\theta \mu_{\theta,k}(a_k,s)\right]^\top
    \left(a_{k-1}-\mu_{\theta,k}(a_k,s)\right)  \\
    &=
    \frac{\eta_k}{\sigma_k}
    \left[\nabla_\theta \epsilon_\theta(a_k,k,s)\right]^\top Z_k .
\end{aligned}
\]
Therefore, using \(\|\nabla_\theta\epsilon_\theta(a,k,s)\|_{\op}\le L_\theta\),
on the above event,
\[
\begin{aligned}
    \left\|\widehat{\bphi}^{\theta}(s,a_{0:K})\right\|_2
    &\le
    \sum_{k=1}^K
    \left\|
    \nabla_\theta \log \pi_\theta(a_{k-1}\mid a_k,s)
    \right\|_2 \\
    &\le
    L_\theta
    \sum_{k=1}^K
    \frac{|\eta_k|}{\sigma_k}\|Z_k\|_2 \\
    &\le
    L_\theta R_{M,K,\delta}
    \sum_{k=1}^K
    \frac{|\eta_k|}{\sigma_k}.
\end{aligned}
\]
Thus, for the realized DDPM pathwise features used by Algorithm 1,
\[
    \left\|\widehat{\bphi}^{\theta}\right\|_2
    \le
    \beta^{\ddpm}_{K,M,\delta}
    :=
    L_\theta R_{M,K,\delta}
    \sum_{k=1}^K
    \frac{|\eta_k|}{\sigma_k}
\]
simultaneously over all \(M\) score evaluations with probability at least
\(1-\delta\).


For the second moment, reveal the Gaussian increments in reverse denoising
order. Each coefficient $\nabla_\theta\epsilon_\theta(a_k,k,s)$ is measurable
before $Z_k$ is drawn, so cross terms have mean zero. Consequently,
\[
\E[\|\widehat\bphi^\theta\|^2\mid s]
=\sum_{k=1}^K\frac{\eta_k^2}{\sigma_k^2}
\E[\|\nabla_\theta\epsilon_\theta(a_k,k,s)\|_F^2\mid s]
\leq d_A L_\theta^2\sum_{k=1}^K\eta_k^2/\sigma_k^2.
\]
Conditional Jensen gives the same second-moment bound for the endpoint score.

    % \textbf{DDPM. } Now, we first address the DDPM case. Observe that
    % \begin{align}
    %     a_{k-1} | a_k \sim \gN\left(\frac{1}{\sqrt{\alpha_k}}\left(a_k + \eta_k \epsilon_\theta(a_k, k, s)\right), \frac{\sigma_k^2}{\alpha_k}\mathbf{I}_d\right).
    % \end{align}
    % This implies that
    % \begin{align}
    %     &\lVert \nabla_\theta \log \pi_\theta(a_{k-1} | a_k, s) \rVert_2 \\
    %     &\leq \left\lVert \frac{\alpha_k}{\sigma^2_k} \left(a_{k-1} - \frac{1}{\sqrt{\alpha_k}}\left(a_k + \eta_k \epsilon_\theta(a_k, k, s)\right)\right)^\top \nabla_\theta \frac{1}{\sqrt{\alpha_k}}\left(a_k + \eta_k \epsilon_\theta(a_k, k, s)\right) \right\rVert_2 \\
    %     &\leq \left\lVert \frac{\alpha_k}{\sigma^2_k} \left(a_{k-1} - \frac{1}{\sqrt{\alpha_k}}\left(a_k + \eta_k \epsilon_\theta(a_k, k, s)\right)\right)\right\rVert_2 \lVert\nabla_\theta a_{k-1}\rVert_2 \\
    %     &= \left\lVert \frac{\sqrt{\alpha_k}}{\sigma_k} Z_k \right\rVert_2 \lVert\nabla_\theta a_{k-1}\rVert_2,
    %     \label{eq:lipschitz-ddpm-log-pi-recursion-3a}
    % \end{align}
    % where $Z_k \sim \gN(0, \mathbf{I}_{d_A})$. With probability at least $1-\delta$, we can apply a standard Gaussian tail bound to see that for all $k=1,...,K$, we have that $\max_k \lVert Z_k \rVert_2 \leq \sqrt{d_A \log (d_A K /\delta)}$. Plugging the above as well as Equation \ref{eq:lipschitz-ddpm-log-pi-recursion-3a} into Equation \ref{eq:lipschitz-unrolled-recursion-3}, and applying Lemma \ref{lem:diffusion-policy-score} in the first line, we see that
    % \begin{align}
    %     \lVert \bphi^{\theta} \rVert_2  = \lVert \nabla_\theta \log \pi_\theta(a_0 | s) \rVert_2 
    %     &\leq \left\lVert \E \sum_{k=1}^K  \nabla_\theta \log \pi_\theta(a_{k-1} | a_k, s) \right\rVert_2\\
    %     &\leq \sup_{a_1,...,a_K} \left\lVert  \sum_{k=1}^K \nabla_\theta \log \pi_\theta(a_{k-1} | a_k, s) \right\rVert_2 \\
    %     &= \sup_{a_1,...,a_K} \left\lVert  \widehat{\bphi}^{\theta} \right\rVert_2 \\
    %     &\leq \sup_{a_1,...,a_K} \sum_{k=1}^K \lVert \nabla_\theta \log \pi_\theta(a_{k-1} | a_k, s) \rVert_2 \\
    %     &\leq \sup_{a_1,...,a_K} \sum_{k=1}^K \left\lVert \frac{\sqrt{\alpha_k}}{\sigma_k} Z_k \right\rVert_2 \lVert\nabla_\theta a_{k-1}\rVert_2 \\
    %     &\leq L_\theta  \sqrt{d_A \log (d_A K /\delta)} \sum_{k=1}^K  \frac{\sqrt{\alpha_k}}{\sigma_k} \sum_{j=k}^K \frac{\eta_j}{\sqrt{\alpha_j}} \prod_{i=j+1}^K 
    %     \frac{1+\eta_i L_a}{\sqrt{\alpha_i}}.
    % \end{align}
    
    \textbf{DDIM. } We now address the DDIM case. This is somewhat more straightforward, as it amounts to bounding the norm of the feature vectors we found in Lemma \ref{lem:feature-concentration-ddim}. We now bound the operator norms of the forward and inverse Jacobians at each timestep. The Jacobian $\nabla_a \Phi_k^\theta(a)$ of the update is given by
    \begin{equation}
        \nabla_a \Phi_k^\theta(a) = \frac{1}{\sqrt{\alpha_k}}\left(\mathbf{I}_d + \eta_k \nabla_a \epsilon_\theta(a_k, k, s)\right).
        \label{eq:lipschitz-jacobian-4}
    \end{equation}
    Therefore, the operator norms of it and its inverse are bounded by
    \begin{equation}
        \lVert \nabla_a \Phi_k^\theta(a) \rVert_{\text{op}} \leq \frac{1+|\eta_k| L_a}{\sqrt{\alpha_k}}, \qquad
        \lVert \nabla_a \Phi_k^\theta(a)^{-1}(a_k) \rVert_{\text{op}} \leq \frac{\sqrt{\alpha_k}}{1-|\eta_k| L_a},
        \label{eq:lipschitz-jacobian-bound-5}
    \end{equation}
These bounds use $1-\|A\|_{\mathrm{op}}\leq\sigma_{\min}(I+A)\leq\sigma_{\max}(I+A)\leq1+\|A\|_{\mathrm{op}}$, with $|\eta_k|L_a<1$ for the inverse bound. This condition also guarantees global inverse existence: for each output $y\in\mathbb R^{d_A}$, the equation $z=\sqrt{\alpha_k}y-\eta_k\epsilon_\theta(z,k,s)$ is a contraction in $z$, so every smooth DDIM step, and hence their composition before any invertible output squash, is a diffeomorphism.

    Now we wish to apply the change of variables formula. By Assumption \ref{aspt:ddpm-bijection}, the full map is a smooth diffeomorphism, so $a_0 = \Phi_{K:0}^\theta(a_K)$ admits a density $\pi_\theta(\cdot | s)$ given by the change of variables formula. 
    Write $\bJ_\theta(a_K) = \nabla_a \Phi_{K:0}^\theta(a_K)$ for the Jacobian of the entire diffusion map. By the change of variables formula, we have that 
    \begin{equation}
        \log \pi_\theta(a_0 | s) = - \frac{d_A\log(2\pi)}{2} - \frac{\lVert a_K \rVert_2^2}{2} - \log\left|\det\left(\bJ_\theta(a_K)\right)\right|,
        \label{eq:lipschitz-change-of-variables-6}
    \end{equation}
    and the derivative of the log-policy with respect to $\theta$ follows from Lemma \ref{lem:feature-concentration-ddim}:
\begin{equation}
\bphi_i^\theta=z^\top w_i-\operatorname{div}_z w_i,
\qquad z=a_K.
\label{eq:lipschitz-derivative-log-policy-7}
\end{equation}
Indeed, differentiating $w_i=\bJ_\theta^{-1}\partial_{\theta_i}\Phi_{K:0}^\theta$
and using equality of mixed derivatives gives
\[
\operatorname{div}w_i
=\operatorname{tr}(\bJ_\theta^{-1}\partial_{\theta_i}\bJ_\theta)
-(\nabla_z\log|\det\bJ_\theta|)^\top w_i.
\]
Thus $\bphi=W_\theta^\top z-q$, and the sufficient full-map bounds in the
lemma imply $\|\bphi\|\leq U_K\|z\|+V_K$. Since $z\sim\mathcal N(0,I_{d_A})$,
$\E\|z\|^2=d_A$ and $\E\|z\|^4=d_A(d_A+2)$. Applying
$(a+b)^2\leq2a^2+2b^2$ and $(a+b)^4\leq8a^4+8b^4$ proves both claims.
In particular this fourth moment also bounds
$\E\|\bphi\bphi^\top-\bLambda^\theta\|_F^2$.
\end{proof}



\subsection{Regularity of standard MLP and U-Net denoisers}
\label{app:practical-network-regularity}

The score-moment and smoothness assumptions are useful only if they can be checked from a denoiser. Here we do so for MLPs and smooth U-Nets with independently trainable bounded weights. The main difficulty is that their inputs are Gaussian and hence unbounded. Bounded weights give polynomial parameter derivatives, which suffice for sampling moments, but need not give a uniform pointwise bound on the policy-score curvature.

For example, the residual map
\[
 F_w(z_1,z_2)=(z_1,z_2+\tanh(wz_1))
\]
is invertible and satisfies the residual contraction condition for
$|w|\leq w_0<1$. Under a standard Gaussian input its log-density is
\[
 \log p_w(a)=-\frac{a_1^2}{2}
             -\frac{(a_2-\tanh(wa_1))^2}{2}-\log(2\pi),
 \qquad
 \left.\partial_w^2\log p_w(a)\right|_{w=0}=-a_1^2.
\]
This curvature is unbounded over actions, but integrable under a Gaussian law. We therefore verify the sampling moments directly and bound the comparator-averaged curvature used in the NPG proof.

We start with the width-$256$ MLP used as a concrete architecture:
\[
 h_1=\tanh(W_1x+b_1),\quad h_2=\tanh(W_2h_1+b_2),\quad
 \epsilon_\theta(a,k,s)=W_3h_2+b_3,\qquad x=(a,e(k,s)),
\]
Both hidden layers have width $256$. The state/time embedding $e$ is fixed and bounded by $E$, and the trainable weights satisfy $\|W_j\|_{\rm op}\leq M_j$, with biases in a bounded set. We measure parameter changes using the Euclidean norm of all weights and biases, or equivalently the Frobenius norm within each weight block. The action derivative passes through all three weight matrices, whereas a parameter derivative starts at the block being changed. This gives
\begin{align*}
 \|D_a\epsilon_\theta\|_{\rm op}&\leq M_3M_2M_1,\\
 \|D_\theta\epsilon_\theta\|_{\rm op}^2
 &\leq257(1+M_3^2)
 +M_3^2M_2^2(\|a\|^2+E^2+1).
\end{align*}
To see the second bound, split a parameter direction into the three layer blocks $u_1,u_2,u_3$. Since each hidden tanh vector has norm at most $\sqrt{256}$, the contributions from the output, second hidden, and first hidden blocks are bounded, respectively, by
\[
 \sqrt{257}\|u_3\|,
 \quad M_3\sqrt{257}\|u_2\|,
 \quad M_3M_2\sqrt{\|a\|^2+E^2+1}\|u_1\|.
\]
Cauchy--Schwarz combines these three bounds. Each further parameter derivative can introduce another factor from the unbounded input, while the first three derivatives of tanh remain bounded. Repeating the chain rule therefore gives constants $C_{ij}$ with
\begin{equation}\label{eq:network-polynomial-derivatives}
 \|D_a^iD_\theta^j\epsilon_\theta(a,k,s)\|
 \leq C_{ij}(1+\|a\|)^j,
 \qquad 1\leq i+j\leq3,\quad j\leq2.
\end{equation}
Here the derivative tensors use the induced Euclidean multilinear norms. The constants are uniform over the bounded parameter set and over $s,k$, and the network itself has at most linear growth in $a$. These polynomial bounds are what we will integrate under the Gaussian denoising inputs.

The calculation also extends through a finite U-Net computation graph. Affine convolutions and fixed linear downsampling or upsampling transform the incoming derivatives linearly. Skip additions add their derivative bounds, and concatenations combine them in the Euclidean norm. Smooth pointwise activations with bounded first three derivatives preserve polynomial growth. We include normalization with a positive stabilizing constant,
\[
 v\longmapsto g\odot
 \frac{Pv}{\sqrt{\|Pv\|^2/n+\delta}}+b,
 \qquad P=I-\mathbf1\mathbf1^\top/n,\quad \delta>0,
\]
with bounded gains and biases. Its denominator is bounded away from zero, so its input Lipschitz constant is at most $\|g\|_\infty/\sqrt\delta$ and its first three input derivatives are bounded.

We can also compute the full input Lipschitz constant $L_\epsilon$ from the graph. Convolutions and fixed resampling operators multiply the incoming bound by their operator norms, and activations multiply it by their scalar Lipschitz constants. Addition sums the bounds, while concatenation takes the square root of their squared sum. Bounded conditioning gains and shifts can be included when their parameter derivatives through order three are bounded. For a U-Net with these operations, this gives both $L_\epsilon$ and the polynomial derivative bounds in \eqref{eq:network-polynomial-derivatives}. Additional attention or conditioning blocks require their own bounds of the same form.

We now pass from these denoiser bounds to the policy quantities used by PG and NPG. For DDIM, a lower bound on each residual Jacobian first ensures that the endpoint density is well defined. Gaussian input moments then control the derivatives of that density.

\begin{prop}[Network regularity under Gaussian denoising inputs]
\label{prop:practical-network-regularity}
Fix a finite denoising depth and a bounded parameter region containing the
update segments. Suppose the denoiser has the growth and derivative
bounds above. For DDIM assume each step
$F_{k,\theta}(a)=c_ka+b_k\epsilon_\theta(a,k,s)$ satisfies
\[
 m_k:=c_k-|b_k|L_\epsilon>0,\qquad c_k>0.
\]
Then the full DDIM map is a smooth diffeomorphism of $\mathbb R^{d_A}$,
with inverse Lipschitz constant at most $\prod_km_k^{-1}$. Its endpoint
parameter score has uniformly finite second and fourth sampling moments.
For DDPM the same moments hold for the path score whenever all transition
variances are positive and parameter-independent. In either case the
discounted value is smooth for bounded rewards and fixed $\gamma<1$.

For DDIM write $F_\theta$ for the full map and let $S$ be an optional fixed
smooth invertible action squash. Suppose the fixed comparator has
\[
 M_{*,3}:=\E_{(s,a)\sim d_\rho^{\pi_{\rm cmp}}}
 \left[1+\|F_{\theta_0}^{-1}(S^{-1}(a),s)\|\right]^3<\infty
\]
for one reference parameter $\theta_0$ in the region. Then it satisfies
comparator-averaged score smoothness, with constant
$L_*\leq C M_{*,3}$ for a finite architecture-dependent $C$. Consequently,
Theorem~\ref{thm:npg-regret} and its mixing and hybrid variants hold with
$L_*$ in place of $L_\theta^{(K)}$, provided their other hypotheses hold.
\end{prop}

\begin{proof}
We first check that the DDIM density can be differentiated. Solving $F_{k,\theta}(a)=y$ is equivalent to finding a fixed point of
\[
 a\longmapsto c_k^{-1}(y-b_k\epsilon_\theta(a,k,s)).
\]
Its Lipschitz constant is $|b_k|L_\epsilon/c_k<1$, so each $y$ has a unique inverse image. Also,
\[
 \|F_{k,\theta}(a)-F_{k,\theta}(a')\|\geq m_k\|a-a'\|.
\]
The inverse is therefore $m_k^{-1}$-Lipschitz. The inverse function theorem gives smoothness, and composing the steps multiplies their inverse Lipschitz bounds.

Next we bound the parameter derivatives of the composed map. At an affine node the input grows at most linearly. Each parameter differentiation of a nonlinear node introduces at most one additional factor with linear input growth, while its input derivatives have bounded coefficients. Induction through the MLP or U-Net graph gives \eqref{eq:network-polynomial-derivatives}, including the skip connections and regularized normalization above. Continuing through the finite denoising composition gives
\[
 \|D_z^iD_\theta^jF_\theta(z,s)\|
 \leq C(1+\|z\|)^j,
 \qquad 1\leq i+j\leq3,\quad j\leq2.
\]
The inverse Jacobian is uniformly bounded by the first part of the proof. When the output action is held fixed, implicit differentiation of $F_\theta(z,s)=a$ therefore yields
\[
 \|D_\theta z\|\leq C(1+\|z\|),\qquad
 \|D_\theta^2z\|\leq C(1+\|z\|)^2.
\]
These inverse derivatives enter the fixed-action log-density
\[
 \log\pi_\theta(a\mid s)
 =-\tfrac12\|z\|^2-\log|\det D_zF_\theta(z,s)|
 -\tfrac{d_A}{2}\log(2\pi).
\]
Differentiating it once and twice, including the dependence of $z$ on
$\theta$, yields
\begin{equation}\label{eq:network-score-growth}
 \|\nabla_\theta\log\pi_\theta(a\mid s)\|
 \leq C(1+\|z\|)^2,\qquad
 \|\nabla_\theta^2\log\pi_\theta(a\mid s)\|_{\rm op}
 \leq C(1+\|z\|)^3.
\end{equation}
A fixed squash contributes a log-Jacobian with zero parameter derivative, so the same estimates hold after squashing. Under policy sampling, $z$ is standard Gaussian. Integrating the polynomial bounds, using its moments through order eight, gives the second and fourth score moments and an integrable likelihood Hessian.

For DDPM we can work directly with the Gaussian transition densities. The reverse recursion grows at most linearly in its preceding action and Gaussian noise, so all action moments are finite uniformly over the parameter region and state. Each path-score term is a fixed coefficient times $[D_\theta\epsilon_\theta(a_k,k,s)]^\top Z_k$. Its polynomial growth gives the required second and fourth moments, as well as a finite first moment for the conditional log-likelihood Hessian. These integrable bounds justify the likelihood differentiations used below.

These score estimates also control value smoothness, because the environment kernel has no policy-parameter derivative. Let $\widehat\bphi_t$ and $H_t$ be the score and log-likelihood Hessian of the policy sample at environment time $t$, including the denoising path for DDPM, and define
$B_2=\sup_{s,\theta}\E\|\widehat\bphi_t\|^2$ and
$B_H=\sup_{s,\theta}\E\|H_t\|_{\rm op}$. The score increments are
martingale differences. For rewards bounded in absolute value by $R$,
likelihood differentiation of the length-$(t+1)$ trajectory gives
\[
 \|\nabla_\theta^2\E r_t\|_{\rm op}
 \leq R\left(\E\Big\|\sum_{j=0}^t\widehat\bphi_j\Big\|^2
             +\sum_{j=0}^t\E\|H_j\|_{\rm op}\right)
 \leq R(t+1)(B_2+B_H).
\]
Summing with weights $\gamma^t$ gives
$L_V\leq R(B_2+B_H)/(1-\gamma)^2$.

Finally, NPG evaluates log-density increments under the comparator rather than the current sampling law. Uniform inverse Lipschitzness and bounded $F_\theta(0,s)$ let us compare the latent coordinates at two parameters:
\[
 1+\|F_\theta^{-1}(F_{\theta_0}(z,s),s)\|
 \leq C(1+\|z\|).
\]
Thus \eqref{eq:network-score-growth} is integrable under the comparator,
uniformly along each update segment. Taylor's integral formula gives
\[
 \E_{d_\rho^{\pi_{\rm cmp}}}\log\frac{\pi_{\theta+d}(a\mid s)}
                                      {\pi_\theta(a\mid s)}
 \geq \left\langle\E_{d_\rho^{\pi_{\rm cmp}}}\bphi^\theta,d\right\rangle
       -\frac{L_*}{2}\|d\|^2.
\]
This replaces the pointwise Taylor bound followed by integration in
Lemma~\ref{lem:npg-regret-lemma}; the rest of that proof and the estimation
arguments are unchanged.
\end{proof}

\paragraph{Conditioning a pretrained Gaussian-latent policy.}
The small-ball comparison policy can have a much better moment bound than division by its conditioning probability suggests. For $Z\sim N(0,I_{d_A})$ and any measurable event $E$ with $\mathbb{P}(E)\geq p>0$, Gaussian concentration gives
\[
 \mathbb{P}(\|Z\|>\sqrt{d_A}+u\mid E)
 \leq\min\{1,p^{-1}e^{-u^2/2}\}.
\]
Put $a_p=\sqrt{2\log(1/p)}$ and
$U=(\|Z\|-\sqrt{d_A}-a_p)_+$. Then
$\mathbb{P}(U>v\mid E)\leq e^{-v^2/2}$ for $v\geq0$. Integrating this tail and using Minkowski's inequality yields, for every $j\geq1$,
\begin{equation}\label{eq:conditioned-gaussian-moment}
 \bigl(\E[(1+\|Z\|)^j\mid E]\bigr)^{1/j}
 \leq 1+\sqrt{d_A}+\sqrt{2\log(1/p)}
       +\bigl(2^{j/2}\Gamma(1+j/2)\bigr)^{1/j}.
\end{equation}
No independence between $E$ and $Z$ is used.

Apply this statewise to the initial actor
$a=S(F_{\theta^{(1)}}(Z,s))$ and the event that $a$ lies in the chosen radius-$r$ ball. If its mass is at least $p_r$ at every relevant state, then the conditional latent law is exactly the one defining $\pi_{{\rm cmp},r}$. The inverse-coordinate comparison in the proof above is uniform on the same bounded parameter region. Together with the quadratic score and cubic Hessian bounds in \eqref{eq:network-score-growth}, it therefore gives
\begin{align*}
 \sup_\theta\left\|\E_{d_\rho^{\pi_{{\rm cmp},r}}}\bphi^\theta\right\|
 &\leq C\bigl(1+\sqrt{d_A}+\sqrt{\log(1/p_r)}\bigr)^2,\\
 L_{*,r}&\leq C\bigl(1+\sqrt{d_A}+\sqrt{\log(1/p_r)}\bigr)^3.
\end{align*}
Here $C$ depends on the architecture and the common parameter region, but not on $r$ or the conditioning event. The statewise bounds permit averaging under the comparison policy's own state occupancy. They retain their explicit Gaussian dimension dependence, and do not assert that the architecture-dependent constant is dimension-free.

For the mass bound $p_r=c_0(r/R)^m$ and $r=R/\sqrt T$, the comparison mean-score bound is $O(\log(2T))$ and the averaged-curvature bound is $O((\log(2T))^{3/2})$, at fixed architecture, $m$, and confidence level. Combined with $\mathcal K_0=O(\log(2T))$, $L=N=\sqrt T$ and the corresponding balanced stepsize, these give statistical NPG terms of order $O(T^{3/4}(\log(2T))^{5/4})$ and a comparison value-gap term of order $O(\sqrt T)$. This statement requires one common bounded parameter region containing the update segments for these horizons, with the remaining Fisher, coefficient, and sampling-moment conditions uniform there. The transferred advantage error still contributes $N\sum_{i=1}^L\sqrt{\epsilon_i}/(1-\gamma)$; Gaussian conditioning does not make that residual small.

\paragraph{Fisher conditioning for a frozen neural backbone.}\label{app:adapter-fisher}
The preceding calculation controls score moments, but empirical NPG also needs conditioning on the identifiable parameter directions. For a trainable location adapter, we can derive this conditioning from the frozen backbone's Lipschitz constant. Let $X=F_\psi(Z,s)$ with $Z\sim N(0,I_{d_A})$ be frozen, and suppose
$F_\psi$ is $L_F$-Lipschitz uniformly in $s$. Write $q_s$ for its smooth
positive density on $\mathbb R^{d_A}$. Gaussian Poincare gives
$\operatorname{Cov}(X\mid s)\preceq L_F^2I$. Writing
$\ell_s=\nabla_x\log q_s$, integration by parts gives
\[
 \E[(X-\E X)\ell_s(X)^\top]=-I.
\]
Cauchy--Schwarz, applied after whitening $X-\E X$, therefore gives
\[
 I_s=\E[\ell_s(X)\ell_s(X)^\top]
 \succeq\operatorname{Cov}(X\mid s)^{-1}\succeq L_F^{-2}I.
\]
The polynomial bounds above justify the integration by parts and make
$I_s$ finite. Let the actor output $X+B_s\theta$, where
$B_s:\mathbb R^p\to\mathbb R^{d_A}$ is fixed. Its parameter score is
$-B_s^\top\ell_s(X)$. Under reset or mixing sampling, the feature-action
law is the current actor conditional on the sampled state. Write
$q_{\theta,S}$ for this state law and
$M_\theta=\E_{s\sim q_{\theta,S}}[B_s^\top B_s]$. If
$\lambda_{\min}^+(M_\theta)\geq b^2>0$ uniformly over the update
region, then
\[
 \bLambda^\theta=\E_{s\sim q_{\theta,S}}[B_s^\top I_s B_s]
 \succeq L_F^{-2}M_\theta,\qquad
 \lambda_{\min}^+(\bLambda^\theta)\geq b^2/L_F^2.
\]
Indeed, $v\in\ker M_\theta$ exactly when $B_sv=0$ almost surely under the
state law, so $M_\theta$ and the parameter Fisher have the same
kernel. The bound therefore holds on their common identifiable span even
when the parameter Fisher is singular. More generally, the same argument applies to any smooth positive frozen density with finite location Fisher $I_s\succeq\iota I$, giving $\lambda_{\min}^+(\bLambda^\theta)\geq\iota b^2$. If the bounded action-domain
interior is diffeomorphic to $\mathbb R^{d_A}$, a fixed such map can be
applied after translation without changing this parameter Fisher.
The hybrid feature-action law need not be the current actor conditional
on its state, and its conditioning is instead handled by
Appendix~\ref{app:hybrid-reset-dist}.
The resulting Fisher lower bound applies to a trainable location adapter on a frozen MLP or U-Net. As in the regret theorem, the derivative bounds must hold throughout the update region, so this is a condition on the iterates rather than a projection step. For independently trained backbone weights, positive-spectrum conditioning and compatible advantage approximation still require separate checks.

\section{Analysis of Vanilla Diffusion Policy Optimization}
\label{app:vanilla-pg}
For the proofs, abbreviate the comparison-policy value gap by $\epsilon_*:=V^*(\rho)-V^{\pi_{\rm cmp}}(\rho)$. This is distinct from the squared advantage residual $\epsilon_{\rm comp}$. We also write $\Delta:=V^*(\nu)-V^{(1)}(\nu)$, $C_A:=4(1+\gamma)/(1-\gamma)^2$, and $M_g^2:=C_A\beta_K^2/(1-\gamma)^2$ for the quantities substituted explicitly in Theorem~\ref{thm:vanilla-pg-regret}.


\subsection{Proof of Lemma \ref{lem:ddpm-ddim-log-policy-densities}}

%\kevin{To construct reverse reverse process, need transition kernel of reverse Markov chain. If this was real DDPM, suffices to just add Gaussian noise, ignoring discretization error. }
\begin{lem}[Features for DDPM]
    Consider the diffusion model features $\widehat{\bphi}^\theta := \sum_{k=1}^K\nabla_\theta\log\pi_{\theta}(a_{k-1}|a_{k}^{(t)},s)$, and the ground-truth gradient of the log-policy ${\bphi}^\theta = \nabla_\theta\log\pi_{\theta}(a|s)$. For positive transition variances,
    $$\widehat{\bphi}^\theta(s,a_{0:K}) = \sum_{k=1}^K \frac{\eta_k^{\mathsf{ddpm}}}{\sigma_k^{\mathsf{ddpm}}}\left[\nabla_\theta \epsilon_\theta\left(a_k, k, s\right)\right]^{\top} Z_k.$$
    \label{lem:features-ddpm}
\end{lem}
\begin{proof}
% For simplicity, we abbreviate $a_{k}^{(t)}$ as $a_k$ and $s^{(t)}$ as $s$.   
%     \begin{align*}
%         \E\left[\left\|\widehat{\bphi}^{(t)} - \bphi^{(t)}\right\|^2\right] = \sum_{i=1}^K 
%     \end{align*}
    Define $\mu_{\theta}(a,k,s) = \frac{1}{\sqrt{\alpha_k}}\big(a + \eta_k^{\mathsf{ddpm}} \epsilon_{\theta}(a, k, s)\big)$.
    \begin{align*}
        \nabla_\theta\log\pi_{\theta^{(t)}}(a_{k-1}|a_{k},s) 
        &= \frac{\alpha_k}{(\sigma_k^{\mathrm{ddpm}})^2} \left[\nabla_\theta \mu_\theta(a_{k},k,s)\right]^\top \big(a_{k-1} - \mu_\theta(a_{k},k,s)\big)\\
        &= \alpha_k \cdot \left[\frac{\nabla_\theta \mu_\theta(a_{k},k,s)}{\sigma_k^{\mathrm{ddpm}}}\right]^\top \cdot \left[\frac{a_{k-1} - \mu_\theta(a_{k},k,s)}{\sigma_k^{\mathrm{ddpm}}}\right]\\
        &= \frac{\eta_k^{\mathsf{ddpm}}}{\sigma_k^{\mathsf{ddpm}}} [\nabla_\theta \epsilon_{\theta}(a_k, k, s)]^\top Z_k,
    \end{align*}
    where $Z_k$'s are i.i.d. standard Gaussian noise used in DDPM update rule~(\ref{eq:DDPM}). The likelihood calculation requires $\sigma_k^{\mathsf{ddpm}}>0$ at every step, including the final step.
\end{proof}


% \kevin{Brownian bridge with only one end fixed. DPPO authors want to finetune the last few steps, i.e. the more dependent end. Solutions: Condition on the last step and look at the Markov chain. Alternatively, tune learning rate so that the last few steps don't matter. }

% \kevin{If this doesn't work, then we can use DDIM with the assumption that there is a one to one mapping. This means that we have the convolution of a Gaussian and the DDIM map as the unconditional action distribution, and given an action there is only one initial noise that produces it.}

\begin{lem}[Exact Endpoint Features for DDIM]
Under Assumption~\ref{aspt:ddpm-bijection}, the sampled latent $a_K$ is uniquely
determined by $(s,a_0)$. Thus the DDIM feature equals the endpoint score
$\widehat\bphi=\bphi=\nabla_\theta\log\pi_\theta(a_0\mid s)$ and is given by
Lemma~\ref{lem:ddpm-ddim-log-policy-densities}, including the spatial derivative
of $\log|\det\bJ_\theta|$.
\label{lem:feature-concentration-ddim}
\end{lem}

\begin{proof}
    By Assumption \ref{aspt:ddpm-bijection}, the map $ \Phi^{(t)}_{K:0}$ is invertible. We therefore have that for any action $a^{(t)}=a_{\cdot,0}^{(t)} \in \gA$, there exists a $a_{\cdot,K}^{(t)} \in \R^d$ that is a realization of $\gN(0, \mathbf{I}_d)$ such that 
    $$a_{\cdot,0}^{(t)} = \Phi^{(t)}_{K:0}(a_{\cdot,K}^{(t)}), \qquad
    a_{\cdot,K}^{(t)} = \left(\Phi^{(t)}_{K:0}\right)^{-1}(a_{\cdot,0}^{(t)}).$$

    The Jacobian with regard to this map is given by the derivative with regard to the action, not the parameter:
    $$\bJ_{\theta^{(t)}}(a_{\cdot,K}^{(t)}) = \nabla_z \Phi^{(t)}_{K:0}(z) \big|_{z=a_{\cdot,K}^{(t)}}, \text{ which, by Assumption \ref{aspt:ddpm-bijection}, is invertible} .$$

    By the change of variables formula, we have that
    $$\pi_{\theta^{(t)}}(a_{\cdot,0}^{(t)}) = \frac{1}{(2\pi)^{d/2}}\exp\left\{-\frac{||a_{\cdot,K}^{(t)}||_2^2}{2}\right\}\left|\det\left(\bJ_{\theta^{(t)}}(a_{\cdot,K}^{(t)})\right)\right|^{-1},$$
    
    $$\log \pi_{\theta^{(t)}}(a_{\cdot,0}^{(t)}) = -\frac{d\log(2\pi)}{2} - \frac{||a_{\cdot,K}^{(t)}||_2^2}{2} - \log \left|\det\left(\bJ_{\theta^{(t)}}(a_{\cdot,K}^{(t)})\right)\right|,$$

    Differentiation here holds $a_{\cdot,0}^{(t)}$ fixed. Write
$z=a_{\cdot,K}^{(t)}$ and $\ell_\theta(z)=\log|\det\bJ_\theta(z)|$. Then
\[
\partial_{\theta_i}\log\pi_\theta(a_0\mid s)
=-z^\top\partial_{\theta_i}z
-\operatorname{tr}(\bJ_\theta^{-1}\partial_{\theta_i}\bJ_\theta)
-(\nabla_z\ell_\theta)^\top\partial_{\theta_i}z.
\]
    %https://math.stackexchange.com/questions/3057656/derivative-of-logarithm-of-determinant for gradient of log determinant.
    Regarding $\nabla_\theta a_{\cdot,K}^{(t)} = \nabla_\theta \left(\Phi^{(t)}_{K:0}\right)^{-1}(a_{\cdot,0}^{(t)})$, we can write $$\varphi(\theta, z) = \Phi^{\theta}_{K:0}(z) - a_{\cdot,0}^{(t)}, \text{ where $a_{\cdot,0}^{(t)}$ is fixed, and } \varphi(\theta^{(t)}, a_{\cdot,K}^{(t)}) = \Phi^{(t)}_{K:0}(a_{\cdot,K}^{(t)}) - a_{\cdot,0}^{(t)} = 0,$$
    as a function $\varphi : \Theta \times \R^d \to \gA$ of the parameter and noise. By the implicit function theorem, there exists an open set $U \subseteq \Theta$ containing $\theta^{(t)}$ and some function $\psi : U \to \R^d$ such that $\psi(\theta^{(t)}) = a_{\cdot,K}^{(t)}$, and $\varphi(\theta, \psi(\theta)) = 0$ for all $\theta \in U$.
    Moreover, by the implicit function theorem as well, we can write
    $$\nabla_\theta \psi(\theta) = -\left(\nabla_a \varphi(\theta, \psi(\theta))\right)^{-1} \nabla_\theta \varphi(\theta, \psi(\theta)).$$
    This implies that
    \begin{align*}
        \nabla_\theta a_{\cdot,K}^{(t)}
        &= \nabla_\theta \psi(\theta^{(t)}) \\
        &= - \left(\nabla_z \varphi(\theta^{(t)}, a_{\cdot,K}^{(t)})\right)^{-1} \nabla_\theta \varphi(\theta^{(t)}, a_{\cdot,K}^{(t)}) \\
        &= -\left(\nabla_z \Phi^{(t)}_{K:0}(a_{\cdot,K}^{(t)})\right)^{-1} \nabla_\theta \Phi^{(t)}_{K:0}(a_{\cdot,K}^{(t)}) \\
        &= -\left(\bJ_{\theta^{(t)}}(a_{\cdot,K}^{(t)})\right)^{-1} \nabla_\theta \Phi^{(t)}_{K:0}(a_{\cdot,K}^{(t)}).
    \end{align*}

    Substituting $\partial_{\theta_i}z=-w_i(z)$ gives
\[
\bphi_i^\theta=z^\top w_i(z)
-\operatorname{tr}(\bJ_\theta^{-1}\partial_{\theta_i}\bJ_\theta)
+(\nabla_z\log|\det\bJ_\theta|)^\top w_i(z).
\]
Given $(s,a_0)$, the inverse map supplies a unique $z$. We can therefore
evaluate the endpoint score from the sampled DDIM path without averaging over
additional latent paths.
\end{proof}

\subsection{Proof of Unbiased Gradient Estimates}

\begin{lem}[Unbiased Policy Gradient Estimator]
    Let $\widehat{A}^{\pi_\theta}, s, a_0,...,a_K$ be a sample from Algorithm \ref{alg:dppo-sampler} run from a distribution $\nu$ at parameter $\theta$. Then, when $\widehat{\bphi}^{\theta} = \sum_{k=1}^K \nabla_\theta \log \pi_\theta(a_{k-1} | a_k, s)$, we have that $\widehat{A}^{\pi_\theta}\widehat{\bphi}^{\theta}/(1-\gamma)$ is an unbiased estimator for the policy gradient starting from $\nu$:
    $$\nabla_\theta V^{\pi_\theta}(\nu) = \frac{1}{1-\gamma}\E_{s,a_0,...,a_K \sim d_{\nu}^{\pi_\theta}}\left[\widehat{A}^{\pi_\theta}(s,a_0) \cdot \widehat{\bphi}^{\theta}(s,a_{0:K})\right].$$
    Also, when $\bphi^\theta = \nabla_\theta \log \pi_\theta(a_0 | s)$, $\widehat{A}^{\pi_\theta}{\bphi}^{\theta}/(1-\gamma)$ is also an unbiased estimator for the policy gradient starting from $\nu$:
    $$\nabla_\theta V^{\pi_\theta}(\nu) = \frac{1}{1-\gamma}\E_{s,a_0,...,a_K \sim d_{\nu}^{\pi_\theta}}\left[\widehat{A}^{\pi_\theta}(s,a_0) \cdot {\bphi}^{\theta}(s,a_{0:K})\right].$$
    \label{lem:unbiased-vanilla-pg}
Furthermore, conditional on the accepted state and denoising path,
$\E[\widehat A\mid s,a_{0:K}]=A^{\pi_\theta}(s,a_0)$ and
$\E[\widehat A^2\mid s,a_{0:K}]\leq C_A=4(1+\gamma)/(1-\gamma)^2$.
Consequently either normalized gradient estimate has second moment at most
$M_g^2=C_A\beta_K^2/(1-\gamma)^2$.
\end{lem}

\begin{proof}
    The claim that $\widehat{A}^{\pi_\theta}$ is unbiased and $s, a_0 \sim d_{\nu}^{\pi_\theta}$ follows from \citet{agarwal2021policygradient} in their construction of their Algorithm 3 (our Algorithm \ref{alg:dppo-sampler} when adapted to the special case of diffusion policies). As such, in an abuse of notation, we can now write $d_{\nu}^{\pi_\theta}$ for a measure over the state-action space and denoising actions yielded by Algorithm \ref{alg:dppo-sampler} here. 
    
    Unbiasedness of the estimator then follows from plugging this into Equation \ref{eq:diffusion-pg-thm}:
    \begin{align}
        &\E_{s,a_0,...,a_K \sim d_{\nu}^{\pi_\theta}}\left[\widehat{A}^{\pi_\theta}(s, a_0)\widehat{\bphi}^{\theta}(s, a_{0:K})\right] \\
        &=  \E_{s,a_0,...,a_K \sim d_{\nu}^{\pi_\theta}}\left[\E[\widehat{A}^{\pi_\theta}(s, a_0) \cdot \widehat{\bphi}^{\theta}(s, a_{0:K})  \mid s, a_{0:K}]\right] \\
        &=  \E_{s,a_0,...,a_K \sim d_{\nu}^{\pi_\theta}}\left[\E[\widehat{A}^{\pi_\theta}(s, a_0) \mid s, a_0] \cdot \widehat{\bphi}^{\theta}(s, a_{0:K})\right] \\
        &= \E_{s,a_0,...,a_K \sim d_{\nu}^{\pi_\theta}}\left[A^{\pi_\theta}(s, a_0) \cdot \widehat{\bphi}^{\theta}(s, a_{0:K})\right] \\
        &= (1-\gamma)\nabla_\theta V^{\pi_\theta}(\nu).
    \end{align}
    The first equality follows from the tower property, the second from the observation that the fact that $\widehat{\bphi}^{\theta}(s, a_{0:K}) $ is a deterministic measurable function of $s, a_0,...,a_K$ allows us to pull it out of the conditional expectation, the third from the unbiasedness of the estimate $\E[\widehat{A}^{\pi_\theta}(s,a) | s,a] = A^{\pi_\theta}(s,a)$, and the fourth from Equation \ref{eq:diffusion-pg-thm}.

    The case for $\bphi^\theta = \nabla_\theta \log \pi_\theta(a_0 | s)$ follows identically, but by plugging the resulting expression into Equation \ref{eq:policy-gradient-thm} instead of Equation \ref{eq:diffusion-pg-thm}:
    \begin{align}
        &\E_{s,a_0 \sim d_{\nu}^{\pi_\theta}}\left[\widehat{A}^{\pi_\theta}(s, a_0){\bphi}^{\theta}(s, a_0)\right] \\
        &=  \E_{s,a_0 \sim d_{\nu}^{\pi_\theta}}\left[\E[\widehat{A}^{\pi_\theta}(s, a_0) \cdot {\bphi}^{\theta}(s, a_{0})  \mid s, a_{0}]\right] \\
        &=  \E_{s,a_0 \sim d_{\nu}^{\pi_\theta}}\left[\E[\widehat{A}^{\pi_\theta}(s, a_0) \mid s, a_0] \cdot {\bphi}^{\theta}(s, a_0)\right] \\
        &= \E_{s,a_0 \sim d_{\nu}^{\pi_\theta}}\left[A^{\pi_\theta}(s, a_0) \cdot {\bphi}^{\theta}(s, a_0)\right] \\
        &= (1-\gamma)\nabla_\theta V^{\pi_\theta}(\nu).
    \end{align}
The return rollout uses fresh randomness conditional on the accepted state
and path. Its horizon $H$ includes its first reward and obeys
$\mathbb{P}(H=h)=(1-\gamma)\gamma^{h-1}$, so
$\E H^2=(1+\gamma)/(1-\gamma)^2$. Since $0\leq r\leq1$, the coin estimator
satisfies $|\widehat A|\leq2H$ in either branch, proving its conditional
second-moment bound. Multiplying by the score moment proves the gradient
bound. One sampler execution uses an expected $(1+\gamma)/(1-\gamma)$
environment transitions, counting occupancy search and the return rollout.
\end{proof}

\subsection{Gradient Domination, Proof of Lemma \ref{lem:gradient-dom}}

Now we provide a counterpart of the gradient domination lemma proved for tabular and softmax policies from \citet{agarwal2020pcpgpolicycoverdirected}. We establish a variant for a smooth policy class with controlled compatible advantage error, finite comparator coverage, and a positive lower bound on the nonzero Fisher spectrum. While we utilize this lemma in Lemma \ref{lem:gradient-dom} for diffusion policy optimization, it does not rely on diffusion-specific properties, and as such is potentially of independent interest. 

\begin{proof}
    By the performance difference lemma in Lemma \ref{lem:pdl}, as well as Cauchy-Schwarz and a change of measure, we see that
    \begin{align}
        V^{\pi_{\rm cmp}}(\rho) - V^{\pi_\theta}(\rho) 
        &= \frac{1}{1-\gamma}\E_{s,a \sim d_\rho^{\pi_{\rm cmp}}}\left[A^{\pi_\theta}(s,a)\right] \\
        &\leq \frac{1}{1-\gamma} \sqrt{\left\lVert\frac{d_\rho^{\pi_{\rm cmp}}}{d_\nu^{\pi_\theta}}\right\rVert_\infty\E_{s,a \sim d_\nu^{\pi_\theta}}\left[ A^{\pi_\theta}(s,a)^2\right]}.
        % &\leq \frac{1}{1-\gamma}\left\lVert\frac{d_\rho^{\pi_{\rm cmp}}}{d_\nu^{\pi_\theta}}\right\rVert_\infty\E_{s \sim d_\nu^{\pi_\theta}}\left[\max_{a \in \gA} A^\pi(s,a)\right] \\
        % &\leq \frac{1}{1-\gamma}\left\lVert\frac{d_\rho^{\pi_{\rm cmp}}}{d_\nu^{\pi_\theta}}\right\rVert_\infty\E_{s \sim d_\nu^{\pi_\theta}}\left[\max_{\pi'} \left\langle A^\pi(s,\cdot), \pi'(\cdot | s) - \pi(\cdot | s) \right\rangle\right] \\
        % &\leq \frac{1}{1-\gamma}\left\lVert\frac{d_\rho^{\pi_{\rm cmp}}}{d_\nu^{\pi_\theta}}\right\rVert_\infty\E_{s \sim d_\nu^{\pi_\theta}}\left[\max_{\pi'} \left\langle Q^\pi(s,\cdot), \pi'(\cdot | s) - \pi(\cdot | s) \right\rangle\right] 
    \end{align}

    By the policy gradient theorem,
    \[\begin{aligned}
\nabla_\theta V^{\theta}(\nu)&=\frac{1}{1-\gamma}\; \Bbb E_{(s,a)\sim d_\nu^{\pi_\theta}} \bigl[A^{\pi_\theta}(s,a)\,\nabla_\theta\log\pi_\theta(a\mid s)\bigr]\\
&= \frac{1}{1-\gamma}\; \Bbb E_{(s,a)\sim d_\nu^{\pi_\theta}} \bigl[A^{\pi_\theta}(s,a)\, \bphi^\theta(s,a)\bigr].
\end{aligned}\]
    We therefore require an upper bound on $\E_{s,a \sim d_\nu^{\pi_\theta}}\left[A^{\pi_\theta}(s,a)^2\right]$ in terms of the norm of \[\E_{s,a \sim d_\nu^{\pi_\theta}}\left[A^{\pi_\theta}(s,a) \bphi^\theta(s,a)\right].\]
    Write $\bLambda=\E_{d_\nu^{\pi_\theta}}[\bphi\bphi^\top]$ and
$\bb=\E_{d_\nu^{\pi_\theta}}[A\bphi]$. If $v\in\ker\bLambda$, then
$\E(v^\top\bphi)^2=0$, hence $v^\top\bb=0$. Thus
$\bb\in\operatorname{range}(\bLambda)$ and
$\bg=\bLambda^\dagger\bb$ solves the normal equations.
Recall that the NPG algorithm updates $\theta$ in this direction. The lemma holds regardless of the update rule: we are seeking to bound the suboptimality gap by the policy-gradient norm.
The residual $e=A-\bphi^\top\bg$ is orthogonal in $L^2(d_\nu^{\pi_\theta})$
to every score coordinate. Expanding the square therefore gives
\[
\E A^2=\bb^\top\bLambda^\dagger\bb+\E e^2
\leq\frac{\|\bb\|^2}{\lambda_{\min}^+(\bLambda)}+\E e^2.
\]
This also shows $\|\bg\|^2\leq\E A^2/\lambda_{\min}^+(\bLambda)$,
which gives the stated sufficient choice of $G$ since $|A|\leq(1-\gamma)^{-1}$.
Under Assumption~\ref{aspt:approx-realizability}, $\E e^2\leq\epsilon_{\mathrm{comp}}$.
Substituting $\bb=(1-\gamma)\nabla_\theta V^{\pi_\theta}(\nu)$ into the
performance-difference bound above yields
\[
V^*(\rho)-V^{\pi_\theta}(\rho)
\leq\epsilon_*+\frac{\sqrt{C_\theta\epsilon_{\mathrm{comp}}}}{1-\gamma}+\sqrt{\frac{C_\theta}{\lambda_{\min}^+(\bLambda)}}
\|\nabla_\theta V^{\pi_\theta}(\nu)\|.
\]
Under exact compatibility the residual term vanishes, recovering gradient
domination with the positive Fisher spectrum.
\end{proof}

\subsection{Proof of Theorem \ref{thm:vanilla-pg-regret}}
We first prove the common bound for reset batches. Appendix~\ref{app:vanilla-pg-mixing} replaces independence by the successive-return variance bound and adds the burn-in coupling error.
\begin{proof}
Write $J(\theta)=V^{\pi_\theta}(\nu)$ and let
$\widehat g_{i,n}=\widehat A_{i,n}\widehat\bphi_{i,n}/(1-\gamma)$ be the normalized gradient output of sample $n$ in batch $i$, using the DDPM path score or the DDIM endpoint score as appropriate. With reset access, the $N$ outputs are independent conditional on the history before the batch. For their mean $\bar g_i=N^{-1}\sum_{n=1}^N\widehat g_{i,n}$, Lemma~\ref{lem:unbiased-vanilla-pg} gives
\[
\E_i\bar g_i=\nabla J(\theta_i),\qquad
\E_i\|\bar g_i-\nabla J(\theta_i)\|^2\leq M_g^2/N.
\]
Smoothness along the update segment therefore implies, for $\eta\leq L_V^{-1}$,
\begin{align*}
\E_i[J(\theta_{i+1})-J(\theta_i)]
&\geq\eta\left(1-\frac{L_V\eta}{2}\right)\|\nabla J(\theta_i)\|^2
-\frac{L_V\eta^2M_g^2}{2N}\\
&\geq\frac\eta2\|\nabla J(\theta_i)\|^2
-\frac{L_V\eta^2M_g^2}{2N}.
\end{align*}
Summing over the $L$ updates yields
\[
\sum_{i=1}^L\E\|\nabla J(\theta_i)\|^2
\leq\frac{2\Delta}{\eta}+\frac{L_V\eta M_g^2L}{N}.
\]
Each policy is used for $N$ labels. Applying Lemma~\ref{lem:gradient-dom} and Cauchy--Schwarz gives
\[
\E\mathrm{Reg}(T)
\leq T\epsilon_*+\frac{T\sqrt{C\epsilon_{\rm comp}}}{1-\gamma}
+\sqrt{\frac C\lambda}\,N\sqrt L
\left(\frac{2\Delta}{\eta}+\frac{L_V\eta M_g^2L}{N}\right)^{1/2},
\]
which is the theorem with $\tau_{\rm mix}=1$ and $\varepsilon_b=0$.

For either access model, the stepsize that balances the two terms is
\[
\eta=\min\left\{L_V^{-1},
\sqrt{\frac{2\Delta N}{L_V\tau_{\rm mix}M_g^2L}}\right\}.
\]
For this choice, the statistical term is at most
\[
\sqrt{\frac C\lambda}\left[
2^{3/4}(\Delta L_V\tau_{\rm mix}M_g^2)^{1/4}T^{3/4}
+2\sqrt{\Delta L_VNT}\right].
\]
Indeed, when the cap is inactive the first term is the exact value of the bound. When the cap is active, $\tau_{\rm mix}M_g^2L/N<2\Delta L_V$, so the second term bounds it. Thus $N\leq\sqrt T$ gives the stated rate. If a tuning constant vanishes, choose a positive upper bound and use it throughout these estimates.

When $N=1$, a direct second-moment calculation gives a sharper constant and does not require the stepsize cap. Take $\eta=\sqrt{2\Delta/(L_VM_g^2T)}$, again using positive upper bounds for zero constants.
Let $\widehat g_t=\widehat A^{(t)}\widehat\bphi^{(t)}/(1-\gamma)$, with the score
chosen as in Algorithm~\ref{alg:dppo-vanilla-pg}. Conditional on the history,
Lemma~\ref{lem:unbiased-vanilla-pg} gives
$\E_t\widehat g_t=\nabla V^{(t)}(\nu)$ and
$\E_t\|\widehat g_t\|^2\leq M_g^2$.
By $L_V$-smoothness along the update segment,
\[
\E_t V^{(t+1)}(\nu)
\geq V^{(t)}(\nu)+\eta\|\nabla V^{(t)}(\nu)\|^2
-\frac{L_V\eta^2M_g^2}{2}.
\]
Summing and using $V^{(T+1)}(\nu)\leq V^*(\nu)$ yields
\begin{equation}
\sum_{t=1}^T\E\|\nabla V^{(t)}(\nu)\|^2
\leq\frac\Delta\eta+\frac{L_V\eta M_g^2T}{2}
=\sqrt{2\Delta L_VM_g^2T}.
\label{eq:vanilla-pg-1}
\end{equation}
By the gradient domination lemma in Lemma \ref{lem:gradient-dom}, we see that
\[
\E\mathrm{Reg}(T)-T\epsilon_*-\frac{T\sqrt{C\epsilon_{\mathrm{comp}}}}{1-\gamma}
\leq\sqrt{C/\lambda_{\min}^+(\bLambda_\nu)}
\sum_{t=1}^T\E\left\lVert\nabla_\theta V^{(t)}(\nu)\right\rVert_2.
\]
By Cauchy-Schwarz and Equation \ref{eq:vanilla-pg-1},
\begin{align*}
    \sum_{t=1}^T \E\left\lVert \nabla_\theta V^{(t)}(\nu) \right\rVert_2
    &\leq \sqrt{T \sum_{t=1}^T \E \left\lVert \nabla_\theta V^{(t)}(\nu) \right\rVert_2^2} \\
    &\leq T^{3/4}(2\Delta L_VM_g^2)^{1/4}.
\end{align*}
Putting it all together,
\[
\E\mathrm{Reg}(T)\leq
T[V^*(\rho)-V^{\pi_{\rm cmp}}(\rho)]
+\frac{T\sqrt{C\epsilon_{\rm comp}}}{1-\gamma}
+\sqrt{\frac C\lambda}
\left[\frac{8[V^*(\nu)-V^{(1)}(\nu)]L_V(1+\gamma)\beta_K^2}
{(1-\gamma)^4}\right]^{1/4}T^{3/4}.
\]
This gives the sharper reset bound for one label per update.
\end{proof}

\subsection{Policy gradients with mixing access}
\label{app:vanilla-pg-mixing}
% Supporting lemmas for surgical restoration. Retains ONE burn-in per batch.
% Earlier per-label-burn-in construction is superseded.
\paragraph{State mixing and complete return samples.}
Write $K_\theta$ for the state transition kernel under $\pi_\theta$.
Assume its stationary state law $m_\theta$ satisfies
\[
 \sup_{\theta,s}\operatorname{TV}(K_\theta^h(s,\cdot),m_\theta)
 \leq c_{\rm mix}\kappa^h,\qquad c_{\rm mix}\geq1,\quad0<\kappa<1.
\]
For vanilla PG, $m_\theta=m$ is fixed, so
$d_m^{\pi_\theta}=m\pi_\theta$ and the objective is
$J(\theta)=V^{\pi_\theta}(m)$. The fixed state law avoids the additional derivative of a policy-dependent stationary distribution. For NPG the state law can depend on the
frozen policy. The joint regression law is $q_\theta=m_\theta\pi_\theta$. Unlike vanilla PG, the NPG proof uses a performance-difference argument and the stationary law only as the compatible-function-approximation distribution, so it does not differentiate through $m_\theta$.
A fixed stationary joint law would force policies to agree on its state
support, so the fixed object in the vanilla assumption is the state law.

In the mixing case, we do not burn in before every call to Algorithm~\ref{alg:dppo-sampler}. Instead, we use an epochwise implementation. At epoch $i$, freeze the policy at $\theta^{(t_i)}$, run $b$ burn-in transitions under $\pi_{\theta^{(t_i)}}$, collect $N$ complete return samples, and update. At each
sample's starting state draw an action and its denoising path and execute
one geometric-return branch of Algorithm~\ref{alg:dppo-sampler}, omitting
its reset and occupancy search. In the value branch the accepted action
is replaced before any action is executed. Only one trajectory is run and
no reset to a visited state is needed. Start the next sample at the terminal
state of that trajectory. Every return uses $H\geq1$ transitions with
$\mathbb{P}(H=h)=(1-\gamma)\gamma^{h-1}$ independently of its trajectory.

\begin{lem}[Second moments of successive return samples]
Fix a policy and start its first sample at stationarity. Let $Y_n$ be a
Hilbert-space-valued output of its $n$th complete sample, including the
accepted action, denoising path and return. If
$\sigma_Y^2=\E\|Y_n-\E Y_n\|^2<\infty$, then
\[
 \E\left\|N^{-1}\sum_{n=1}^N Y_n-\E Y_1\right\|^2
 \leq\tau_{\rm mix}\sigma_Y^2/N,\qquad
 \tau_{\rm mix}=1+4\sqrt{c_{\rm mix}}/(1-\sqrt\kappa).
\]
\end{lem}
\begin{proof}
The states $X_n$ at sample boundaries have kernel
\[
 R_\theta=\sum_{h\geq1}(1-\gamma)\gamma^{h-1}K_\theta^h.
\]
Both return branches execute the current policy after averaging over their
action choice. Hence this identity also holds when the value branch
replaces the accepted action. The kernel has stationary law $m_\theta$.
Its $j$th power is a mixture of $K_\theta^h$ with $h\geq j$, giving
$\sup_x\operatorname{TV}(R_\theta^j(x,\cdot),m_\theta)
\leq c_{\rm mix}\kappa^j$.
For centered Hilbert-valued $f\in L^2(m_\theta)$, weighted Cauchy--Schwarz
and the total variation bound give
\[
 \|R_\theta^jf(x)\|^2\leq2c_{\rm mix}\kappa^j
 \int\|f(y)\|^2|R_\theta^j(x,dy)-m_\theta(dy)|.
\]
Integrating in $m_\theta(dx)$, using
$|R_\theta^j-m_\theta|\leq R_\theta^j+m_\theta$ and stationarity, gives
\[
 \|R_\theta^jf\|_{L^2(m_\theta)}^2
 \leq4c_{\rm mix}\kappa^j\|f\|_{L^2(m_\theta)}^2.
\]
For $j=0$ the same loose bound holds because $c_{\rm mix}\geq1$.
Set $g(x)=\E[Y_n-\E Y_n\mid X_n=x]$. Then $g$ is centered and
$\|g\|_{L^2(m_\theta)}^2\leq\sigma_Y^2$.
Conditional on the complete $n$th sample, future samples depend on it only
through $X_{n+1}$. Therefore for $h\geq1$,
\[
 \left|\E\langle Y_n-\E Y_n,Y_{n+h}-\E Y_{n+h}\rangle\right|
 =\left|\E\langle Y_n-\E Y_n,R_\theta^{h-1}g(X_{n+1})\rangle\right|
 \leq2\sqrt{c_{\rm mix}}\kappa^{(h-1)/2}\sigma_Y^2.
\]
Expanding the squared norm of the average and summing this geometric bound
proves the result. Thus the bound includes dependence of complete returns.
\end{proof}

\paragraph{Vanilla PG.}
For a stationary-start epoch use the normalized gradient output
$Y=\widehat A\widehat\bphi/(1-\gamma)$, with the DDPM path or DDIM endpoint score.
Its conditional mean is $\nabla J(\theta_i)$ and second moment at most
$M_g^2$. The batch mean therefore has conditional variance at most
$\tau_{\rm mix}M_g^2/N$. Smoothness and $\eta\leq L_V^{-1}$ imply
\[
 \E_i[J(\theta_{i+1})-J(\theta_i)]
 \geq\frac\eta2\|\nabla J(\theta_i)\|^2
       -\frac{L_V\eta^2\tau_{\rm mix}M_g^2}{2N}.
\]
Thus the original telescoping step gives
\[
 \sum_{i=1}^L\E\|\nabla J(\theta_i)\|^2
 \leq2\Delta/\eta+L_V\eta\tau_{\rm mix}M_g^2L/N.
\]
Let $T=LN$, $\mathfrak e=\epsilon_*+
\sqrt{C\epsilon_{\rm comp}}/(1-\gamma)$ and
$\mathfrak c=\sqrt{C/\lambda}$. The original gradient-domination and
Cauchy--Schwarz steps give
\[
 \E\operatorname{Reg}(T)\leq T\mathfrak e+
 \mathfrak cN\sqrt{L(2\Delta/\eta+
               L_V\eta\tau_{\rm mix}M_g^2L/N)}
 +\frac{TL\varepsilon_b}{1-\gamma}.
\]
The last term couples the full adaptive run to stationary-start epochs.
The conditional initial-state coupling fails in at most $L\varepsilon_b$
probability, and regret is at most $T/(1-\gamma)$.
Choose
\[
 \eta=\min\{L_V^{-1},
       \sqrt{2\Delta N/(L_V\tau_{\rm mix}M_g^2L)}\}.
\]
For $N\leq\sqrt T$, the nonapproximation term is $O(T^{3/4})$.
For example $\varepsilon_b\leq T^{-2}$ and
$N\asymp bT^{1/4}\leq\sqrt T$ preserve the original epochwise rate while
making the additive burn-in regret $Lb/(1-\gamma)=O(T^{3/4})$.
Positive upper bounds may replace zero smoothness or gap constants.

\paragraph{Environment interactions.}
There are $Lb$ burn-in transitions and $T$ complete geometric returns,
for expected transition count $Lb+T/(1-\gamma)$. With probability
$1-\delta$ every return uses at most
\[h=\lceil\log(T/\delta)/\log(1/\gamma)\rceil\]
transitions, so the total
is at most $Lb+Th$. Charging each transition its policy's value gap gives
regret at most $h\operatorname{Reg}(T)+Lb/(1-\gamma)$ on this event.
The stated $T$ therefore counts labels, and the interaction accounting
also includes return lengths. This is a discounted-value pseudo-regret,
not a claim about undiscounted reward regret of the continuing chain.

\section{Analysis of DDIM NPG}
\label{app:diffusion-npg}

\subsection{NPG Update Concentration}

\begin{lem}[NPG Update Direction Bound for Oracle Features]
Condition on the history $\mathcal F_i$ before batch $i$. Under the reset,
positive-spectrum, and moment conditions of Theorem~\ref{thm:npg-regret}, let
$\widehat\bg^{(t_i)}$ be the minimum-norm minimizer of
$\widehat\gL_A^{(t_i)}$ over $\gG$ and
$\bg^{(t_i)}=(\bLambda^{(t_i)})^\dagger\bb^{(t_i)}$.
Then
\[
\E[\|\widehat\bg^{(t_i)}-\bg^{(t_i)}\|\mid\mathcal F_i]\leq a_N.
\]
\label{lem:least-squares-npg}
\end{lem}
\begin{proof}


    % By the standard learning theory argument,
    % \begin{align*}
    %     &\gL_A^{(t)}(\widehat{\bg}^{(t)},\theta^{(t)}) - \gL_A^{(t)}(\bg^{(t)},\theta^{(t)}) \\
    %     &= \gL_A^{(t)}(\widehat{\bg}^{(t)},\theta^{(t)}) - \widehat{\gL}_A^{(t)}(\widehat{\bg}^{(t)},\theta^{(t)}) + \widehat{\gL}_A^{(t)}(\widehat{\bg}^{(t)},\theta^{(t)}) - \widehat{\gL}_A^{(t)}(\bg^{(t)},\theta^{(t)}) + \widehat{\gL}_A^{(t)}(\bg^{(t)},\theta^{(t)}) - \gL_A^{(t)}(\bg^{(t)},\theta^{(t)}) \\
    %     &\leq 2\sup_{\bg \in \gG} \left({\gL}_A^{(t)}-\widehat{\gL}_A^{(t)}\right)(g,\theta^{(t)}).
    % \end{align*}

    % Now we evaluate the expectation of the original loss and the diffusion loss. Define
    % $$\widehat{\gL}_A^{(t)}(g,\theta^{(t)}) := \frac{1}{N}\sum_{n=1}^N \left(\widehat{A}^{(t_i,n)} - \left\langle g, \nabla_\theta \log\pi_{\theta^{(t)}}(a^{(t_i,n)} | s^{(t_i,n)}) \right\rangle \right)^2,$$
    % noting that ${\gL}_A^{(t)}(g,\theta^{(t)}) = \E\left[\widehat{\gL}_A^{(t)}(g,\theta^{(t)})\right]$, and consider
    % \begin{align*}
    %     &\E\left[\widehat{\gL}_A^{(t)}(g,\theta^{(t)})\right] \\
    %     &= \E\left[\frac{1}{N}\sum_{n=1}^N \left(\widehat{A}^{(t_i,n)} - \left\langle g, \nabla_\theta \log\pi_{\theta^{(t)}}(a^{(t_i,n)} | s^{(t_i,n)}) \right\rangle \right)^2\right] \\
    %     &= \E\left[\frac{1}{N}\sum_{n=1}^N \left(\widehat{A}^{(t_i,n)} - \left\langle g, \E_{a_{\cdot,1:K}}\left[\sum_{k=0}^K \nabla_\theta \log \pi_{\theta^{(t)}}(a_{\cdot,k}|a_{\cdot,k+1},s^{(t_i,n)})\right]\right\rangle \right)^2\right],
    % \end{align*}
    % where the last line holds by Lemma \ref{lem:diffusion-policy-score}. Continuing the argument, and writing $$\widetilde{\bphi}^{(t_i,n)} := \E_{a_{\cdot,1:K}}\left[\sum_{k=0}^K \nabla_\theta \log \pi_{\theta^{(t)}}(a_{\cdot,k}|a_{\cdot,k+1},s^{(t_i,n)})\right],$$
    % we expand the square to obtain
    % \begin{align*}
    %     \E\left[\frac{1}{N}\sum_{n=1}^N \left(\widehat{A}^{(t_i,n)} - \left\langle g, \widetilde{\bphi}^{(t_i,n)}\right\rangle \right)^2\right]
    %     &= \E\left[\frac{1}{N}\sum_{n=1}^N \left(\left(\widehat{A}^{(t_i,n)}\right)^2 -2\widehat{A}^{(t_i,n)}\left\langle g, \widetilde{\bphi}^{(t_i,n)}\right\rangle + \left\langle g, \widetilde{\bphi}^{(t_i,n)}\right\rangle^2 \right)\right].
    % \end{align*}

    % Similarly, for the diffusion loss,
    % \begin{align*}
    %     \E\left[\widehat{\gL}_A^{(t)}(g,\theta^{(t)})\right]
    %     &= \E\left[\frac{1}{N}\sum_{n=1}^N \left(\widehat{A}^{(t_i,n)} - \left\langle g, \widehat{\bphi}^{(t_i,n)} \right\rangle \right)^2\right]\\
    %     &= \E\left[\frac{1}{N}\sum_{n=1}^N \left(\left(\widehat{A}^{(t_i,n)}\right)^2 -2\widehat{A}^{(t_i,n)}\left\langle g, \widehat{\bphi}^{(t_i,n)}\right\rangle + \left\langle g, \widehat{\bphi}^{(t_i,n)}\right\rangle^2 \right)\right].
    % \end{align*}

    % \kevin{Weird. The diffusion policy gradient seems to be unbiased for the policy gradient proper, but there seems to be some Jensen bias for the squared loss. }

    % Recall the identity $(a-b)^2 - (a-c)^2 = (c-b)(2a-b-c)$. We can then use this to show that
    % \begin{align*}
    %     \Delta_n &= \left(\widehat{A}^{(t_i,n)} - \left\langle g, \widehat{\bphi}^{(t_i,n)}\right\rangle\right)^2 - \left(\widehat{A}^{(t_i,n)} - \left\langle g, \widetilde{\bphi}^{(t_i,n)}\right\rangle\right)^2 \\
    %     &= \left\langle g, \widetilde{\bphi}^{(t_i,n)} - \widehat{\bphi}^{(t_i,n)}\right\rangle \cdot \left(2\widehat{A}^{(t_i,n)} - \left\langle g, {\bphi}^{(t_i,n)}\right\rangle - \left\langle g, \widetilde{\bphi}^{(t_i,n)}\right\rangle\right).
    % \end{align*}
    % \kevin{This is wrong. $\Delta_n$ is not necessarily mean zero.}

    % Now, by construction in Algorithm \ref{alg:dppo-sampler}, note that the advantage estimates have a tail exponentially decaying by $(1-\gamma)^x$. Also note that $||g||_2$ is bounded by $G$, and $||{\bphi}^{(t_i,n)}||_2, ||\widetilde{\bphi}^{(t_i,n)}||_2 \leq R_K$, then for some fixed $g \in \gG$ we have that
    % $$|\Delta_n| \leq (2GR_K)(2A_{\max} + 2GR_K) = (2GR_K)(4\log(T/\delta) + 2GR_K).$$

    % By a standard union bound over an $\epsilon$-net over the $g \in \gG$, 
    % $$\frac{1}{N}\sum_{n=1}^N|\Delta_n| \leq O\left(GR_K(\log(T/\delta) + GR_K)\sqrt{\frac{d+\log(T/\delta)}{N}}\right).$$

    % \kevin{This rate is not fast enough. We need the squared loss to concentrate at a rate $d/N$.}

    Let $\mathcal S_\theta=\operatorname{range}(\bLambda)$ for the frozen policy.
If $P$ projects onto $\mathcal S_\theta^\perp$, then
$\E_i\|P\bphi\|^2=\operatorname{tr}(P\bLambda P)=0$, so every sampled feature
lies in $\mathcal S_\theta$ almost surely. Projecting a candidate coefficient
onto the span of the sampled features preserves its empirical predictions
and cannot increase its norm. The minimum-norm empirical minimizer therefore
lies in that span, and hence in $\mathcal S_\theta$, as does
$\bg=\bLambda^\dagger\bb$.

    Define $\widehat\bLambda=N^{-1}\sum_n\bphi_n\bphi_n^\top$ and
$\widehat\bb=N^{-1}\sum_n\widehat A_n\bphi_n$, suppressing the batch index.
Due to the multi-batched updating rule, we do not update the policy within each batch, and as such the feature vectors are i.i.d. in the reset case conditional on $\mathcal F_i$, allowing us to continue our analysis. We address the mixing case in Appendix \ref{app:npg-mixing}. The sample averages have means $\bLambda$ and $\bb$. The moment assumptions and the sampler give
\[
\E_i\|\widehat\bLambda-\bLambda\|_F^2\leq M_F/N,
\qquad
\E_i\|\widehat\bb-\bb\|^2\leq C_A\beta_K^2/N.
\]
Let $\mathcal E=\{\|\widehat\bLambda-\bLambda\|_{\mathrm{op}}\leq\lambda/2\}$.
Markov's inequality yields $\mathbb{P}_i(\mathcal E^c)\leq4M_F/(\lambda^2N)$.
On $\mathcal E$, the sample covariance restricted to $\mathcal S_\theta$ is bounded
below by $\lambda/2$. Since $\bg=\bLambda^\dagger\bb$ lies in
$\gG$, constrained optimality gives
\[
\langle\widehat\bLambda\widehat\bg-\widehat\bb,
\bg-\widehat\bg\rangle\geq0.
\]
Writing $h=\widehat\bg-\bg\in\mathcal S_\theta$ and using $\bLambda\bg=\bb$,
\[
\frac\lambda2\|h\|^2
\leq h^\top\widehat\bLambda h
\leq\|h\|\bigl(\|\widehat\bb-\bb\|
+G\|\widehat\bLambda-\bLambda\|_{\mathrm{op}}\bigr).
\]
Thus on $\mathcal E$, $\|h\|$ is at most $2/\lambda$ times the expression
in parentheses. On its complement, both coefficients remain in the radius
$G$ ball, so $\|h\|\leq2G$. Taking conditional expectations and applying
Cauchy--Schwarz proves
\[
\E_i\|h\|\leq
\frac{2(\sqrt{C_A}\beta_K+G\sqrt{M_F})}{\lambda\sqrt N}
+\frac{8GM_F}{\lambda^2N}=a_N.
\]
The positive spectrum controls the coefficient error on $\mathcal E$, while
the regression constraint controls the remaining batches. This is why the
argument remains valid even when an empirical covariance is singular.
\end{proof}


\begin{lem}[NPG Prediction Error for Oracle Features]
Let $\mathcal F_i$ be the pre-batch history and
\[
\mathrm{err}^{(t_i)}=
\E_{d_\rho^{\pi_{\rm cmp}}}
[A^{(t_i)}-\langle\widehat\bg^{(t_i)},\bphi^{(t_i)}\rangle].
\]
Then
$\E[\mathrm{err}^{(t_i)}\mid\mathcal F_i]
\leq\sqrt{\epsilon_{\mathrm{bias}}}+B_*a_N$.
Furthermore, with probability at least $1-\delta$,
\[
\sum_{i=1}^L\mathrm{err}^{(t_i)}
\leq L(\sqrt{\epsilon_{\mathrm{bias}}}+B_*a_N)
+2B_*G\sqrt{2L\log(1/\delta)}.
\]
\label{lem:npg-err-t-bound}
\end{lem}

% \begin{proof}
%     Consider 
%     \begin{align*}
%         \text{err}^{(t)} 
%         &=  \E_{d^*}\left[A^{(t)}(s,a) - \left\langle\widehat{\bg}^{(t)}, \bphi^{(t)}(s,a)\right\rangle\right] \\
%         &= \E_{d^*}\left[A^{(t)}(s,a) - \left\langle{\bg}^{(t)}, \bphi^{(t)}(s,a)\right\rangle\right] 
%         + \E_{d^*}\left[\left\langle{\bg}^{(t)} - \widehat{\bg}^{(t)}, \bphi^{(t)}(s,a)\right\rangle\right].
%     \end{align*}

%     The first term, corresponding to the asymptotic bias of the least-squares solution, can be bounded by an application of Cauchy-Schwarz and Assumption \ref{aspt:approx-realizability} as in \citet{agarwal2021policygradient}:
%     $$\E_{d^*}\left[A^{(t)}(s,a) - \left\langle{\bg}^{(t)}, \bphi^{(t)}(s,a)\right\rangle\right] \leq \sqrt{\E_{d^*}\left[\left(A^{(t)}(s,a) - \left\langle{\bg}^{(t)}, \bphi^{(t)}(s,a)\right\rangle\right)^2\right]} = \sqrt{\gL_A^*(\bg^{(t)}, \theta^{(t)})} \leq \sqrt{\epsilon_{\text{bias}}}.$$

    

%     We can now bound the second term corresponding to the error of the least-squares solution:
%     \begin{align*}
%         \E_{d^*}\left[\left\langle{\bg}^{(t)} - \widehat{\bg}^{(t)}, \bphi^{(t)}(s,a)\right\rangle\right]
%         &\leq  \E_{d^*}\left[\left\lVert{\bg}^{(t)} - \widehat{\bg}^{(t)}\right\rVert_2 \left\lVert \bphi^{(t)}\right\rVert_2\right] \\
%         &\leq \E_{d^*}\left[\frac{\beta_K^2}{\lambda_{\min}(\bLambda_\nu)^2}\sqrt{\frac{\log^2(pT/\delta)}{(1-\gamma)^6 N}} \cdot \beta_K \right] \\
%         &= \frac{\beta_K^3}{\lambda_{\min}(\bLambda_\nu)^2}\sqrt{\frac{\log^2(pT/\delta)}{(1-\gamma)^6 N}}.
%     \end{align*}

    
% \end{proof}

% \begin{lem}[NPG Prediction Error for Estimated Features]
% \end{lem}
\begin{proof}
    Consider 
    \begin{align}
        \text{err}^{(t)} 
        &= \E_{d_\rho^{\pi_{\rm cmp}}}\left[A^{(t)}(s,a) - \left\langle\widehat{\bg}^{(t)}, {\bphi}^{(t)}(s,a)\right\rangle\right] \\
        &= \E_{d_\rho^{\pi_{\rm cmp}}}\left[A^{(t)}(s,a) - \left\langle{\bg}^{(t)}, \bphi^{(t)}(s,a)\right\rangle\right] + \E_{d_\rho^{\pi_{\rm cmp}}}\left[\left\langle{\bg}^{(t)}, \bphi^{(t)}(s,a)\right\rangle - \left\langle\widehat{\bg}^{(t)}, {\bphi}^{(t)}(s,a)\right\rangle\right]\\
        &= \E_{d_\rho^{\pi_{\rm cmp}}}\left[A^{(t)}(s,a) - \left\langle{\bg}^{(t)}, \bphi^{(t)}(s,a)\right\rangle\right] 
        + \E_{d_\rho^{\pi_{\rm cmp}}}\left[\left\langle{\bg}^{(t)} - \widehat{\bg}^{(t)}, \bphi^{(t)}(s,a)\right\rangle\right].
    \end{align}

    The first term, corresponding to the asymptotic bias of the least-squares solution, can be bounded by an application of Cauchy-Schwarz and the comparator transfer condition of Theorem~\ref{thm:npg-regret} as in \citet{agarwal2021policygradient}:
    \begin{align}
        \E_{d_\rho^{\pi_{\rm cmp}}}\left[A^{(t)}(s,a) - \left\langle{\bg}^{(t)}, \bphi^{(t)}(s,a)\right\rangle\right] &\leq \sqrt{\E_{d_\rho^{\pi_{\rm cmp}}}\left[\left(A^{(t)}(s,a) - \left\langle{\bg}^{(t)}, \bphi^{(t)}(s,a)\right\rangle\right)^2\right]}\\
    &= \sqrt{\gL_A^{d_\rho^{\pi_{\rm cmp}}}(\bg^{(t)}, \theta^{(t)})} \leq \sqrt{\epsilon_{\text{bias}}}.
    \end{align}

    

    Let $\mu_i=\E_{d_\rho^{\pi_{\rm cmp}}}\bphi^{(t_i)}$, which is
$\mathcal F_i$-measurable. The second term is
$\mu_i^\top(\bg^{(t_i)}-\widehat\bg^{(t_i)})$, so its conditional mean is
at most $B_*a_N$ by Lemma~\ref{lem:least-squares-npg}. For the probability
bound, the centered prediction error is exactly
\[
\mathrm{err}^{(t_i)}-\E_i\mathrm{err}^{(t_i)}
=-\mu_i^\top(\widehat\bg^{(t_i)}-\E_i\widehat\bg^{(t_i)}).
\]
Its absolute value is at most $2B_*G$ since the fitted direction is constrained
to $\gG$. These centered errors form a martingale difference sequence.
Azuma's inequality proves the asserted sum bound.
For the batch-dependent transfer condition following
Theorem~\ref{thm:npg-regret}, the same Cauchy--Schwarz step bounds the
conditional bias in batch $i$ by $\sqrt{\epsilon_i}$.
The centered martingale is unchanged, so the proof gives
$\sum_i\sqrt{\epsilon_i}$ in place of $L\sqrt{\epsilon_{\mathrm{bias}}}$.
Substituting this bound in Lemma~\ref{lem:npg-regret-lemma} proves the
cumulative-transfer variant without changing the algorithm.
\end{proof}


\subsection{Finetuning only the task-relevant action coordinates}
\label{app:intrinsic-npg-lift}

A small initial KL divergence is only one way pretraining enters the NPG bound. The Fisher, score moments and advantage-approximation terms can still grow with action dimension. We now identify a setting where they remain unchanged: only $m$ coordinates affect rewards and transitions, and finetuning changes their distribution while preserving the conditional law of the other coordinates. A fixed decoder may mix both sets of coordinates in the displayed action. Differentiating the policy then removes the fixed conditional density and decoder Jacobian, reducing every learning term to the active $m$-dimensional problem.

\begin{prop}[NPG with frozen nuisance coordinates]
\label{prop:intrinsic-npg-lift}
Let an active-action MDP have action $x\in\mathcal X\subseteq\mathbb R^m$,
policy $\bar\pi_\theta(x\mid s)$, reward $\bar r(s,x)$ and transition
$\bar P(\cdot\mid s,x)$. For any $d\geq m$, let
$\kappa_s(z\mid x)$ be a fixed normalized conditional density on
$\mathcal Z\subseteq\mathbb R^{d-m}$, and let
$\Psi_s:\mathcal X\times\mathcal Z\to\mathcal A$ be a fixed smooth
bijection with nonzero Jacobian. Define
\[
 \pi_\theta(\Psi_s(x,z)\mid s)
 =\frac{\bar\pi_\theta(x\mid s)\kappa_s(z\mid x)}
 {|\det D\Psi_s(x,z)|},\quad
 r(s,\Psi_s(x,z))=\bar r(s,x),\quad
 P(\cdot\mid s,\Psi_s(x,z))=\bar P(\cdot\mid s,x).
\]
The conditional joint law of the reward and next state depends only on
$(s,x)$ and agrees with the active MDP. Suppose the actor has an evaluable
DDIM endpoint score and the active
problem satisfies Theorem~\ref{thm:npg-regret}. Lift its comparator by the
same conditional density and decoder. Then the full-action problem has
exactly the same values of
\[
 \mathcal K_0,\ \epsilon_*,\ \lambda,\ \beta_K,\ M_F,\ G,\
 L_\theta^{(K)},\ B_*,\ \epsilon_{\rm bias},\ C_A.
\]
The batchwise transfer errors can likewise be kept identical. Under a
coupling of the sampler randomness, the minimum-norm empirical NPG
updates and the regret are equal pathwise. Thus the entire regret bound,
not only its initial KL term, is independent of $d-m$.
\end{prop}
\begin{proof}
At a fixed state and action, $(x,z)=\Psi_s^{-1}(a)$ is independent of the
trainable parameter. Consequently
\[
 \nabla_\theta\log\pi_\theta(a\mid s)
 =\nabla_\theta\log\bar\pi_\theta(x\mid s),\qquad
 \nabla_\theta^2\log\pi_\theta(a\mid s)
 =\nabla_\theta^2\log\bar\pi_\theta(x\mid s).
\]
Integrating out $z$ leaves the same state transition kernel. The state
occupancies and values therefore agree, and
$A^{\pi_\theta}(s,\Psi_s(x,z))=A^{\bar\pi_\theta}(s,x)$.
The Fisher matrix, its positive spectrum, both score moment expressions,
and the population compatible-regression loss are identical. The
comparator integrals defining $B_*$ and the transferred squared error are
identical for the same reason. Cancellation of the fixed conditional
law and Jacobian gives
\[
 \KL(\widetilde\pi(\cdot\mid s)\|\pi_\theta(\cdot\mid s))
 =\KL(\widetilde{\bar\pi}(\cdot\mid s)\|\bar\pi_\theta(\cdot\mid s)).
\]
Every full-action policy induces an active-action marginal with the same
value, and every active-action policy has a lift. Hence the unrestricted
optimal values, and therefore $\epsilon_*$, agree as well.

Use the same initial states, geometric stopping times, active actions,
reward and transition randomness, and fair Q/V coins in the two reset
samplers, adding only the nuisance draws to the full-action sampler.
Each pair $(\widehat A,\bphi)$ is then equal, so the empirical losses and
their minimum-norm constrained minimizers are equal. Induction over
batches proves equality of the actual updates and regret. Thus $C_A$,
the coefficient bound $a_N$, and every term of Theorem~\ref{thm:npg-regret}
are unchanged when nuisance coordinates are appended.
\end{proof}

This also gives an intrinsic small-ball version without requiring a
full-dimensional density to have a lower-dimensional small-ball exponent.
Suppose $\bar\pi^{(1)}(B(x^*(s),r)\mid s)\geq c_0(r/R)^m$ and the
active optimal $Q$-function is locally $L_Q$-Lipschitz. Condition only
$x$ on this ball and leave $\kappa_s(\cdot\mid x)$ unchanged. The
conditioning event in the full action space is the cylinder
$\Psi_s(B(x^*(s),r)\times\mathcal Z)$, and its probability is exactly
the active small-ball probability. Writing $\mathcal K_0(r)$ for this comparator's initial KL cost, we obtain
\[
 \mathcal K_0(r)\leq\log(c_0^{-1})+m\log(R/r),\qquad
 \epsilon_*\leq L_Qr/(1-\gamma).
\]
For $r=R/\sqrt T$, the comparison-gap and optimization terms in Theorem~\ref{thm:npg-regret} are at most
\[
 \frac{L_QR\sqrt T+
 G\sqrt{2L_\theta^{(K)}[1+\log(c_0^{-1})+(m/2)\log T]}\,T^{3/4}}
 {1-\gamma}.
\]
All remaining terms have the active-problem constants. These constants
can depend on $m$, the comparator scale and the active learning problem;
they do not depend on the number of frozen nuisance coordinates.

For a fully verified nonlinear instance, take the one-dimensional
trainable actor in Proposition~\ref{prop:two-layer-shape-learning} and
append $a_j=\tanh Z_j$ for $j=2,\ldots,d$, with independent standard
Gaussian $Z_j$. Let the reward depend only on $a_1$. This is a one-step
DDIM actor with the same coefficients as that proposition and denoising
field $\epsilon_\theta(z)=\sqrt2z-(A_\theta(z_1),\tanh z_2,\ldots,
\tanh z_d)$. Its two trainable coordinates change the hidden and output
weights of the active nonlinear network. Every constant in that
proposition, including its cumulatively verified transfer error, remains
unchanged for all $d$. The result concerns statistical learning cost. Generating or decoding a
$d$-dimensional action can still require dimension-dependent computation.

\paragraph{Redundant parameter coordinates.}\label{app:redundant-parameter-lift}
The same learner can have a singular parameter Fisher. Write $\phi$ and $F$ for the original policy score and Fisher. For a fixed lift
$\theta=P\vartheta$ with $PP^\top=I$, the score is $P^\top\phi$,
and the Fisher and its pseudoinverse are $P^\top F P$ and
$P^\top F^\dagger P$, with the same positive spectrum as $F$.
A component of a fitted coefficient in $\ker P$ changes no prediction
and can only increase its norm. Minimum-norm constrained regression
therefore gives $\widehat g_\vartheta=P^\top\widehat g_\theta$.
Since $P^\top$ preserves norms and $PP^\top=I$, the lifted updates
reproduce the original updates and regret exactly when initialized with $P\vartheta^{(1)}=\theta^{(1)}$.

\subsection{NPG Analysis}

\begin{lem}[Modified NPG Regret Lemma (Lemma 6.2, \citet{agarwal2021policygradient})]
Let $\pi_{\rm cmp}$ and $\mathcal K_0$ be the comparison policy and initial KL cost in Theorem~\ref{thm:npg-regret}, and write $\epsilon_*:=V^*(\rho)-V^{\pi_{\rm cmp}}(\rho)$. For $L$ batches of $N$ samples and bounded
updates $\|\widehat\bg^{(t_i)}\|\leq G$, the score-smoothness and finite-KL
conditions imply
\[
\sum_{t=1}^T(V^*(\rho)-V^{(t)}(\rho))
\leq T\epsilon_*+\frac{N}{1-\gamma}
\left(\frac {\mathcal K_0}\eta+\frac{\eta L_\theta^{(K)}G^2L}{2}
+\sum_{i=1}^L\mathrm{err}^{(t_i)}\right).
\]
Here $\mathrm{err}^{(t_i)}$ is defined using $d_\rho^{\pi_{\rm cmp}}$ as in
Lemma~\ref{lem:npg-err-t-bound}.
\label{lem:npg-regret-lemma}
\end{lem}
\begin{proof}
    The proof is similar to that of Lemma 6.2 within \citet{agarwal2021policygradient}. However, to accommodate the diffusion setting, there is some additional bookkeeping needed to account for estimation error and batched updates. We therefore provide a proof here.
    
    By Assumption \ref{aspt:ddpm-bijection}, we know that for all $s, a \in \gS \times \gA$,
    \begin{equation}\left\lVert\nabla_\theta \log \pi_{\theta^{(t_i+1)}}(a|s) - \nabla_\theta \log \pi_{\theta^{(t_i)}}(a|s) \right\rVert_2 \leq L^{(K)}_\theta\left\lVert \theta^{(t_i+1)} - \theta^{(t_i)}\right\rVert_2.
    \end{equation}
    By Taylor's theorem, we then have that
    \begin{align}
        &\left| \log \pi_{\theta^{(t_i+1)}}(a|s) - \log \pi_{\theta^{(t_i)}}(a|s) - \left(\nabla_\theta \log \pi_{\theta^{(t_i)}}(a|s) \right)^\top \left(\theta^{(t_i+1)} - \theta^{(t_i)}\right) \right| \\
        &\leq \frac{L^{(K)}_\theta}{2} \left\lVert \theta^{(t_i+1)} - \theta^{(t_i)}\right\rVert_2^2.
    \end{align}
    This in turn implies
    \begin{align}
        \log\left(\frac{\pi_{\theta^{(t_i+1)}}(a|s)}{\pi_{\theta^{(t_i)}}(a|s)}\right) 
        &\geq \left(\nabla_\theta \log \pi_{\theta^{(t_i)}}\right)^\top \left(\theta^{(t_i+1)} - \theta^{(t_i)}\right) - \frac{L^{(K)}_\theta}{2} \left\lVert \theta^{(t_i+1)} - \theta^{(t_i)}\right\rVert_2^2 \\
        &= \eta \left\langle \bphi^{(t_i)}(a|s), \widehat{\bg}^{(t_i)} \right\rangle - \frac{\eta^2 L^{(K)}_\theta}{2} \left\lVert  \widehat{\bg}^{(t_i)} \right\rVert_2^2.
        \label{eq:npg-lemma-1}
    \end{align}

    Integrate Equation~\ref{eq:npg-lemma-1} under the fixed comparator
occupancy $d_\rho^{\pi_{\rm cmp}}$. Define
\[
\mathcal K_{i-1}=\E_{s\sim d_{\rho,S}^{\pi_{\rm cmp}}}
\KL(\pi_{\rm cmp}(\cdot\mid s)\|\pi^{(t_i)}(\cdot\mid s)).
\]
Here $\mathcal K_{i-1}$ is the potential before batch $i$, and $\mathcal K_L$ uses the policy after the last update. The Taylor bound above
makes each log-policy ratio integrable under the comparator, since its score
features are integrable and the update direction is bounded. Starting from
$\mathcal K_0<\infty$, the log-ratio identity therefore implies that all potentials $\mathcal K_0,\ldots,\mathcal K_L$ are finite.

    By the performance difference lemma from \citet{kakade2002approxoptimalrl} in Lemma \ref{lem:pdl} and Equation \ref{eq:npg-lemma-1}, we find that
    \begin{align}
        &\E_{d_\rho^{\pi_{\rm cmp}}}\left[\log\left(\frac{\pi_{\theta^{(t_i+1)}}(a|s)}{\pi_{\theta^{(t_i)}}(a|s)}\right)\right]  \\
        &\geq \eta \cdot \E_{d_\rho^{\pi_{\rm cmp}}}\left[\left\langle \bphi^{(t_i)}(s,a), \widehat{\bg}^{(t_i)} \right\rangle\right] - \frac{\eta^2 L^{(K)}_\theta}{2} \left\lVert  \widehat{\bg}^{(t_i)} \right\rVert_2^2 \\
        &= \eta \cdot \E_{d_\rho^{\pi_{\rm cmp}}}\left[A^{\pi_{\theta^{(t_i)}}}(s,a)\right] - \frac{\eta^2 L^{(K)}_\theta}{2} \left\lVert  \widehat{\bg}^{(t_i)} \right\rVert_2^2 + \eta \cdot \E_{d_\rho^{\pi_{\rm cmp}}}\left[\left\langle \bphi^{(t_i)}(s,a), \widehat{\bg}^{(t_i)}\right\rangle - A^{(t_i)}(s,a)\right] \\
        &= \eta(1-\gamma)\left(V^{\pi_{\rm cmp}}(\rho) - V^{(t_i)}(\rho)\right) - \frac{\eta^2 L^{(K)}_\theta}{2} \left\lVert  \widehat{\bg}^{(t_i)} \right\rVert_2^2 - \eta \cdot \text{err}^{(t_i)}.
    \end{align}

    The expected log-policy ratio is $\mathcal K_{i-1}-\mathcal K_i$. It is now time to account for the batched update. Summing over all $t=1,...,T$, noting that $\pi^{(t_{i-1}+1)} = ... = \pi^{(t_i)}$, we see that
    \begin{align}
        &\sum_{t=1}^T \left(V^{\pi_{\rm cmp}}(\rho) - V^{(t)}(\rho)\right) \\
        &= N \sum_{i=1}^L \left(V^{\pi_{\rm cmp}}(\rho) - V^{(t_i)}(\rho)\right) \\
        &\leq N \sum_{i=1}^L \frac{1}{1-\gamma}\left(\frac{\mathcal K_{i-1}-\mathcal K_i}{\eta} + \frac{\eta L^{(K)}_\theta G^2}{2} + \text{err}^{(t_i)}\right).
    \end{align}

    However, the sum of KL-divergences below telescopes:
    \begin{align}
        &\sum_{i=1}^L \left(\mathcal K_{i-1}-\mathcal K_i\right) \\
        &=\mathcal K_0-\mathcal K_L\\
        &\leq \mathcal K_0.
    \end{align}
    This then yields:
    \begin{align}
        \sum_{t=1}^T \left(V^{\pi_{\rm cmp}}(\rho) - V^{(t)}(\rho)\right)
        &\leq \frac{N\mathcal K_0}{\eta(1-\gamma)} + \frac{\eta L^{(K)}_\theta G^2 NL}{2(1-\gamma)} + \frac{N}{1-\gamma}\sum_{i=1}^L \text{err}^{(t_i)}.
    \end{align}
Adding $T\epsilon_*$ gives the result.
\end{proof}

To obtain the pretrained comparator described after Theorem~\ref{thm:npg-regret}, write $E_r(s)=\{a:\|a-a^*(s)\|\leq r\}$ and $Z_r(s)=\pi^{(1)}(E_r(s)\mid s)$. We condition the initial policy on this set:
\[
\begin{gathered}
\pi_{{\rm cmp},r}(da\mid s)
=\frac{\mathbf 1_{E_r(s)}(a)\pi^{(1)}(da\mid s)}{Z_r(s)},\\
\KL(\pi_{{\rm cmp},r}(\cdot\mid s)\|\pi^{(1)}(\cdot\mid s))
=-\log Z_r(s)
\leq\log(c_0^{-1})+m\log(R/r).
\end{gathered}
\]
Averaging under $d_{\rho,S}^{\pi_{{\rm cmp},r}}$ bounds $\mathcal K_0$. The performance-difference identity and local Lipschitz condition give
\[
V^*(\rho)-V^{\pi_{{\rm cmp},r}}(\rho)
=\frac{1}{1-\gamma}\E_{s,a\sim d_\rho^{\pi_{{\rm cmp},r}}}
[Q^*(s,a^*(s))-Q^*(s,a)]
\leq\frac{L_Qr}{1-\gamma}.
\]
For $r=R/\sqrt T$, set $\overline{\mathcal K}_0(T)=1+\log(c_0^{-1})+(m/2)\log T$. With $L=N=\sqrt T$ and $\eta=\sqrt{2\overline{\mathcal K}_0(T)/(LL_\theta^{(K)}G^2)}$, the comparator and optimization terms in Theorem~\ref{thm:npg-regret} are at most
\[
\frac{L_QR\sqrt T+G\sqrt{2L_\theta^{(K)}\overline{\mathcal K}_0(T)}\,T^{3/4}}{1-\gamma}.
\]
The remaining estimation, transfer, and high-probability terms are unchanged, with their constants evaluated for this comparator and the resulting updates. The comparator is only used in the analysis and need not be represented by the actor or known to the learner.


\subsection{Proof of Theorem \ref{thm:npg-regret}}
The reset proof below uses Lemma~\ref{lem:least-squares-npg}. For mixing and hybrid access, the coefficient bounds in Appendices~\ref{app:npg-mixing} and \ref{app:hybrid-reset-dist} replace $a_N$ under the corresponding regression law. The prediction-error and KL-telescope steps are otherwise the same.
% \kevin{Choose $N=T^{1/2}$, $L=T^{1/2}$. Weirdly enough we do not need to address a log-covering number issue -- we don't need to cover the policy class. We are not off-policy, and don't need to bound the error over a policy class. Same as Agarwal.}
\begin{proof}
Substitute Lemma~\ref{lem:npg-err-t-bound} into
Lemma~\ref{lem:npg-regret-lemma}. This gives
\begin{align}
\mathrm{Reg}(T)\leq{}&T\epsilon_*+
\frac N{1-\gamma}\left(\frac {\mathcal K_0}\eta+
\frac{\eta L_\theta^{(K)}G^2L}{2}\right)
+\frac T{1-\gamma}(\sqrt{\epsilon_{\mathrm{bias}}}+B_*a_N)\nonumber\\
&+\frac{2NB_*G}{1-\gamma}\sqrt{2L\log(1/\delta)}.
\label{eq:ddim-regret-1}
\end{align}
Taking expectations in the two lemmas instead removes the last term.
For $\mathcal K_0>0$, the stated stepsize balances $\mathcal K_0/\eta$ with
$\eta L_\theta^{(K)}G^2L/2$. With $L=N=\sqrt T$, this term and the
$N^{-1/2}$ coefficient-estimation term scale as $T^{3/4}$, while the
$N^{-1}$ term is $O(\sqrt T)$. The two approximation terms account for the
comparison policy's value gap and the transferred advantage error.
\end{proof}
% \begin{lem}[DDIM NPG Regret Bound]
%     Let  $L = \sqrt{T}$, $N = \sqrt{T}$, $\eta = \Theta\left(\frac{d_A \log(d_A R_a K)}{L\beta_K G^2}\right)$, and choose some $K \geq K^*$. With probability at least $1-\delta$, Algorithm \ref{alg:dppo-reinforce} with the DDIM sampler yields
%     $$ \text{Reg}(T)
%     \leq O\left(\sqrt{\frac{\beta_K^{\mathsf{ddim}} G^2 T d_A \log(d_A R_a K)}{(1-\gamma)^2}}  +  \frac{\beta_K^3 T^{3/4} \log(pT/\delta)}{\lambda_{\min}(\bLambda_\nu)^2 (1-\gamma)^4}\right).$$
% \end{lem}
% \begin{proof}
%     Choosing $L = \sqrt{T}$, $N = \sqrt{T}$, $\eta = \Theta\left(\frac{d_A \log(d_A R_a K)}{L\beta_K G^2}\right)$ and plugging this choice into Lemma \ref{lem:npg-regret-lemma} yields:
%     \begin{align}
%         \sum_{t=1}^T \left(V^*(\rho) - V^{(t)}(\rho)\right)
%         &\leq \sqrt{\frac{2\beta_K^{\mathsf{ddim}} G^2 T d_A \log(d_A R_a K)}{(1-\gamma)^2}} + \frac{1}{1-\gamma}\sum_{t=1}^T \text{err}^{(t)}.
%         \label{eq:ddim-regret-1}
%     \end{align}
    
%     By Lemma \ref{lem:npg-err-t-bound} and our choice of $N=\sqrt{T}$, we have that 
%     $$\text{err}^{(t)} \leq \sqrt{\epsilon_{\text{bias}}} + \frac{\beta_K^3 \log(pT/\delta)}{\lambda_{\min}(\bLambda_\nu)^2 (1-\gamma)^3 T^{1/4}} + \beta_K \cdot \E_{d^*}\left[ \left\lVert\widehat{\bg}^{(t)} - \widehat{\bg}^{(t)}\right\rVert_2 \right].$$

%     By Lemma \ref{lem:feature-concentration-ddim}, we have that $\widehat{\bphi}^{(t)} = \bphi^{(t)}$, and so it must hold that $\widehat{\bg}^{(t)} = \widehat{\bg}^{(t)}$. This observation and Equation \ref{eq:ddim-regret-1} yields:
%     \begin{align}
%         \sum_{t=1}^T \left(V^*(\rho) - V^{(t)}(\rho)\right)
%         &\leq \sqrt{\frac{2\beta_K^{\mathsf{ddim}} G^2 T d_A \log(d_A R_a K)}{(1-\gamma)^2}}  +  \frac{\beta_K^3 T^{3/4} \log(pT/\delta)}{\lambda_{\min}(\bLambda_\nu)^2 (1-\gamma)^4} + \frac{\sqrt{\epsilon_{\text{bias}}}\cdot T}{1-\gamma}.
%     \end{align}

%     Our expression for $\beta_K^{\mathsf{ddim}}$ in Lemma \ref{lem:lipschitz-constant-policy}, as well as setting $K$ large enough such that $K \geq K^*$, yields the result.

% \end{proof}

\subsection{Natural policy gradients with mixing access}
\label{app:npg-mixing}
\paragraph{NPG coefficient and regret bounds.}
Apply the successive-return lemma in Appendix~\ref{app:vanilla-pg-mixing}
to $Y=\phi\phi^\top$ in Frobenius space and to
$Y=\widehat A\phi$ in parameter space. For a stationary-start batch,
the existing coefficient proof changes only by replacing $M_F/N$ and
$C_A\beta_K^2/N$ with $\tau_{\rm mix}M_F/N$ and
$\tau_{\rm mix}C_A\beta_K^2/N$. The existing label moment bound is
conditional on accepted state, action and path, so it applies to these
complete return outputs.
For an actual batch, condition on all its prior history and frozen
parameter. Its post-burn-in state is within
$\varepsilon_b=c_{\rm mix}\kappa^b$ of $m_\theta$. Maximal coupling of
this state, followed by the same complete-sample randomness, couples the
whole batch with failure probability at most $\varepsilon_b$.
The stationary reference coefficient obeys the preceding bound. Both
empirical coefficients have norm at most $G$, including on failed couplings,
so no range condition is needed for transient samples. This yields
\[
 \E_i\|\widehat g_i-g_i\|\leq a_N^{\rm mix}:=
 \frac{2\sqrt{\tau_{\rm mix}}(\sqrt{C_A}\beta_K+G\sqrt{M_F})}
      {\lambda\sqrt N}
 +\frac{8G\tau_{\rm mix}M_F}{\lambda^2N}+2G\varepsilon_b.
\]
The prediction-error and KL proofs now apply unchanged, with
$a_N^{\rm mix}$ replacing $a_N$. Taking $L=N=\sqrt T$ and
$\varepsilon_b\leq T^{-1}$ preserves the $T^{3/4}$ exponent. The burn-in
count is $Lb=\sqrt T\,b$.


\begin{thm}[DDIM NPG Mixing Regret]\label{thm:npg-mixing-regret}
Under the policy-dependent state-mixing case of
Assumption~\ref{aspt:coverage}, use the epochwise sampler in
Appendix~\ref{app:vanilla-pg-mixing} and impose the NPG moment, coefficient,
and comparator conditions under $q_\theta=m_\theta\pi_\theta$.
Theorem~\ref{thm:npg-regret} holds with $a_N^{\rm mix}$ in place of $a_N$.
For $L=N=\sqrt T$ and $c_{\rm mix}\kappa^b\leq T^{-1}$, its statistical
terms are $O(T^{3/4})$, with $Lb$ additional burn-in transitions and the
return-length accounting of Appendix~\ref{app:vanilla-pg-mixing}.
\end{thm}
The conditional coefficient bound above is the only change needed in the
prediction-error and KL proofs, and therefore proves the theorem.

\subsection{Natural policy gradients with offline trajectories}
\label{app:hybrid-reset-dist}

Write $F_{{\rm on},i}$ and $F_{{\rm off},i}$ for the score
covariances under the online law and $\mu$, respectively. Their equally
weighted mixture is $F_i=(F_{{\rm on},i}+F_{{\rm off},i})/2$, with
identifiable span $\mathcal R_i=\operatorname{range}(F_i)$.

\begin{lem}[Hybrid covariance on an identifiable span]
\label{lem:hybrid-reset-dist}
If the two component covariances have the same range $\mathcal R_i$, then
\[
 \lambda_{\min}(F_i|_{\mathcal R_i})
 \geq\frac{\lambda_{\min}^+(F_{{\rm on},i})+
                 \lambda_{\min}^+(F_{{\rm off},i})}{2}
 \geq\frac{\max\{\lambda_{\min}^+(F_{{\rm on},i}),
                         \lambda_{\min}^+(F_{{\rm off},i})\}}2.
\]
More generally, if the offline covariance has range $\mathcal R_i$, then
$\lambda_{\min}(F_i|_{\mathcal R_i})
\geq\lambda_{\min}^+(F_{{\rm off},i})/2$, even when the online covariance
is singular on this span.
\end{lem}
\begin{proof}
For every unit vector $v\in\mathcal R_i$, average the two inequalities
$v^\top F_{x,i}v\geq\lambda_{\min}^+(F_{x,i})$ when both ranges equal
$\mathcal R_i$. When only the offline range equals $\mathcal R_i$, use
$v^\top F_{{\rm on},i}v\geq0$ instead. Minimize over these unit vectors.
\end{proof}
The same-range condition is needed to use the maximum of the component
positive eigenvalues. Without it, the theorem uses the positive spectrum
of their sum directly. For the online law $d_\rho^{\pi_i}$, if the initial-state covariance
$F_{\rho,i}=\E_{s\sim\rho,a\sim\pi_i}\phi_i\phi_i^\top$ and the offline
covariance also have range $\mathcal R_i$, the discounted-occupancy comparison
$F_{{\rm on},i}\succeq(1-\gamma)F_{\rho,i}$ gives
\[
 \lambda_{\min}(F_i|_{\mathcal R_i})
 \geq\tfrac12\max\{(1-\gamma)\lambda_{\min}^+(F_{\rho,i}),
                         \lambda_{\min}^+(F_{{\rm off},i})\}.
\]
Thus offline features can supply conditioning that the online samples lack.
The comparator transfer condition is still imposed on the mixed regression coefficient, as in Theorem~\ref{thm:npg-regret}.


\paragraph{Current-policy Monte Carlo labels from logged trajectories.}
Let the offline state-action law be
$\mu(\mathrm ds,\mathrm da)=\mu_S(\mathrm ds)\beta(\mathrm da\mid s)$.
The offline records used below contain behavior-policy continuations, not
only individual transitions. We use per-decision importance sampling
\citep{precup2000eligibility}. Reserve independent trajectory blocks before
learning and reveal one block at each update. The current policy may depend
on all previously revealed blocks, but not the unused block. In particular,
unused blocks are excluded from initialization, pretraining and model
selection. The behavior
policy can be evaluated and sampled computationally at a stored state.
Assume $\pi_i(\cdot\mid s)\ll\beta(\cdot\mid s)$ on the relevant
trajectory states. Evaluating the current policy density and endpoint score
at stored actions is also required. For an invertible DDIM actor this entails
inverse-map computation, whose efficient implementation is a separate issue.

Use the geometric stopping rule in Algorithm~\ref{alg:dppo-sampler}:
the offline record has independent length $H\geq1$ with
$\mathbb{P}(H>t)=\gamma^t$. The record contains all rewards and
transitions up to this stopping time. Its initial pair has law $\mu$, and
subsequent actions follow $\beta$. With the current policy frozen, define
\[
 w_0=\frac{\pi_i(A_0\mid S_0)}{\beta(A_0\mid S_0)},\qquad
 W_0=1,\qquad
 W_t=\prod_{k=1}^t
             \frac{\pi_i(A_k\mid S_k)}{\beta(A_k\mid S_k)},\qquad
 \widehat Q_i=\sum_{t=0}^{H-1}W_tR_t.
\]
The stopping mechanism is independent of the underlying behavior trajectory.
It may already be part of data collection, or be applied to a record that
contains the requisite continuation. A fixed finite record cannot implement
an unbounded stopping time without truncation.

Flip a fair coin, independently of the record. On its first branch use the
feature pair $(S_0,A_0)$ and label $2\widehat Q_i$. On its second branch draw
$A'\sim\beta(\cdot\mid S_0)$ independently of the record conditional on
$S_0$, and use the feature pair $(S_0,A')$ with label
$-2w_0\widehat Q_i$. The additional action draw is a policy computation,
not an environment query. Both branches have feature law $\mu$.

\begin{lem}[Importance-weighted advantage labels]\label{lem:hybrid-is-label}
The label $\widehat A_i$ above satisfies
\[
 \E[\widehat A_i\mid S_0=s,\text{ feature action }a]
       = A^{\pi_i}(s,a).
\]
No environment restart or fitted critic is used to obtain an offline label.
\end{lem}
\begin{proof}
Change of measure at times $1,\ldots,t$, and the independent stopping coin,
give
\[
 \E[\mathbf1\{H>t\}W_tR_t\mid S_0=s,A_0=a]
    =\gamma^t\E_{\pi_i}[R_t\mid S_0=s,A_0=a].
\]
The rewards are nonnegative and bounded, so monotone convergence gives
$\E[\widehat Q_i\mid s,a]=Q^{\pi_i}(s,a)$.
Similarly $\E[w_0\widehat Q_i\mid S_0=s]=V^{\pi_i}(s)$.
On the second branch the feature action is independent of this return
conditional on the state. Since the two branch feature laws coincide, the
coin remains fair after conditioning on the feature pair. Averaging the
two conditional label means proves the identity.
\end{proof}

A state baseline can be included without requiring an accurate critic.
For any $B_i(s)$ fixed before the fresh block is revealed, replace the
two labels by $2(\widehat Q_i-B_i(s))$ and
$-2w_0(\widehat Q_i-B_i(s))$, respectively. Their conditional means
still average to $(Q^{\pi_i}-B_i)-(V^{\pi_i}-B_i)=A^{\pi_i}$.
Its score-weighted second moment must still be controlled, and an arbitrary
baseline need not reduce variance. The familiar cancellation of a state
baseline under on-policy score expectations \citep{schulman2018gae} does
not by itself give this identity under the offline feature law $\mu$.

\paragraph{Moment condition and a sufficient likelihood-ratio bound.}
The regression proof only needs the conditional, score-weighted moment
\[
 \E\|\phi_i(S_0,A)\widehat A_i\|^2\leq M_{b,\rm off}
 \tag{IS moment}\label{eq:hybrid-is-moment}
\]
uniformly over the histories at which a fresh block is used. This condition
can hold without a uniform bound on the likelihood ratios. A sufficient
condition that makes the dependence explicit is
\[
 \sup_s\E_{a\sim\beta(\cdot\mid s)}
      \left(\frac{\pi_i(a\mid s)}{\beta(a\mid s)}\right)^2\leq R,
       \qquad \gamma R<1.
\]
For the uncentered labels ($B_i=0$), $\E[W_t^2\mid S_0,A_0]\leq R^t$,
and Minkowski's inequality gives
\[
 \bigl(\E[\widehat Q_i^2\mid S_0,A_0]\bigr)^{1/2}
 \leq\sum_{t\geq0}(\gamma R)^{t/2}
 =\frac1{1-\sqrt{\gamma R}}=:\sqrt{C_Q}.
\]
The two branches then satisfy
\[
 \E[\widehat A_i^2\mid S_0=s,\text{ feature action }a]
 \leq 2C_Q(1+R)=:C_{A,\rm off}.
\]
Thus $M_{b,\rm off}=C_{A,\rm off}\beta_{\rm off}^2$ is sufficient when
$\E_\mu\|\phi_i\|^2\leq\beta_{\rm off}^2$.
The result below uses this score-weighted moment condition. The uniform
ratio bound is one sufficient way to verify it.

\paragraph{Hybrid regression and cost.}
Use $N$ online Monte Carlo labels from the original online sampling law
$q_{{\rm on},i}$ and $N$ independent offline records per update, retaining
the equal source weights in the scalar squared loss. Define
\[
 q_i=\tfrac12q_{{\rm on},i}+\tfrac12\mu,\qquad
 F_i=\E_{q_i}[\phi_i\phi_i^\top],\qquad
 b_i=\E_{q_i}[\phi_iA^{\pi_i}],\qquad g_i=F_i^\dagger b_i.
\]
Retain the coefficient bound and positive spectrum
\[
\|g_i\|\leq G,\qquad \lambda_{\min}^+(F_i)\geq\lambda,
\]
and all regularity and comparator conditions of
Theorem~\ref{thm:npg-regret}, now for this same mixed design.
For $x\in\{{\rm on},{\rm off}\}$, suppose the two sources have
covariance-estimation and target moments
\[
\begin{aligned}
\E_x\|\phi_i\phi_i^\top-F_{x,i}\|_F^2&\leq M_{F,x},\\
\E_x\|\phi_i\widehat A_i\|^2&\leq M_{b,x},
\end{aligned}
\]
For conditionally independent online labels set
\[
\begin{gathered}
 u_F^2=\frac{M_{F,\rm on}+M_{F,\rm off}}{4N},\qquad
 u_b^2=\frac{M_{b,\rm on}+M_{b,\rm off}}{4N},\\
 a_N^{\rm IS}=\frac{2(u_b+Gu_F)}\lambda
                    +\frac{8Gu_F^2}{\lambda^2}.
\end{gathered}
\]
\begin{thm}[Hybrid NPG with importance-weighted Monte Carlo labels]
\label{thm:hybrid-npg-regret}
Under the preceding conditions, Theorem~\ref{thm:npg-regret}
holds with $a_N$ replaced by $a_N^{\rm IS}$. Here $T=LN$ counts online
regression labels and $N$ offline records are used per batch. In particular, with
$L=N=\sqrt T$ and fixed moment and regularity constants, its statistical
terms retain order $T^{3/4}$, with the same comparator and transfer terms.

The offline budget is $T$ independent behavior records, containing
$T/(1-\gamma)$ stored transitions in expectation. For
\[
h_{\rm cost}=\lceil\log(T/\delta_{\rm cost})/\log(1/\gamma)\rceil,
\]
all record lengths are
at most $h_{\rm cost}$ with probability at least $1-\delta_{\rm cost}$, so their total
length is at most $Th_{\rm cost}$. Offline labeling uses no online environment
interactions. Online return sampling retains its own interaction cost.
\end{thm}
\begin{proof}
Condition on the revealed history preceding an update. The unused offline
records, their branch coins, and the auxiliary behavior actions are fresh.
Lemma~\ref{lem:hybrid-is-label} makes the empirical target unbiased for
$b_i$. Independence and the source second moments give the displayed
$u_F,u_b$. Every feature lies almost surely in $\operatorname{range}(F_i)$:
a vector in the null space has zero squared feature projection under
both positive-weight source laws. The same constrained normal-equation
argument as in Lemma~\ref{lem:least-squares-npg} gives
\[\E_i\|\widehat g_i-g_i\|\leq a_N^{\rm IS}.\] The NPG KL telescope and
bounded-update prediction martingale then apply directly. No event about
future importance weights is conditioned upon. The length bound follows
from \[
\mathbb{P}(H>h_{\rm cost})=\gamma^{h_{\rm cost}}\leq\delta_{\rm cost}/T
\]
and a union bound.
\end{proof}

\paragraph{Online sampling and adaptive data reuse.}
The original hybrid statement uses $q_{{\rm on},i}=d_\rho^{\pi_i}$ from
ordinary initial-state interaction. This construction adds no operation
that restarts the environment at a stored state. If online interaction
must be a single continuing trajectory, use the stationary complete-return
block construction with $q_{{\rm on},i}=m_{\pi_i}\pi_i$. Its proved mixing
bound multiplies the online terms in $u_F^2,u_b^2$ by $\tau_{\rm mix}$,
and a per-batch burn-in coupling adds $2G\epsilon_b$ to $a_N^{\rm IS}$.
This continuing version explicitly requires that mixing condition and
changes the online design law.

Reserved independent offline blocks justify the conditional moments above.
Sampling fresh indices from one adaptively reused fixed trajectory pool
does not by itself do so: conditional on the pool, its rewards and
transitions are fixed. Such reuse requires a uniform concentration or
stability argument, or a separate bound on its empirical-target bias.
Individual transition tuples without their
continuations and behavior likelihoods do not supply these Monte Carlo
labels. These data and moment conditions are additional to covariance
coverage, rather than consequences of it.


\subsection{A stochastic multimodal DDIM example}
\label{app:joint-stochastic-example}

We now ask whether a multimodal actor can satisfy the learning conditions jointly and approach the optimal value when its updates use noisy returns. An exponential family makes the relation between reward and policy-score features explicit: its advantage is exactly compatible, and reparameterizing by Fisher arc length keeps its positive Fisher eigenvalue equal to one. We can then follow the empirical NPG updates as the policy concentrates near the rewarding modes.

For this actor we use the comparator-averaged curvature condition already derived for neural denoisers in Appendix~\ref{app:practical-network-regularity}. The NPG proof lower-bounds the expected log-density increment under a fixed comparison law $q$. Writing $f_q(\theta)=\mathbb E_q\log\pi_\theta(a)$, the needed inequality is
\begin{equation}\label{eq:joint-average-curvature}
 f_q(\theta+d)-f_q(\theta)
 \geq f_q'(\theta)d-\frac{L_q}{2}d^2
\end{equation}
along each update segment. Pointwise score smoothness implies this inequality, but an integrable bound under $q$ is sufficient. The construction below realizes the actor by inverse-CDF sampling. Appendix~\ref{app:two-layer-shape-example} gives a finite neural actor and instead verifies the pointwise condition, with its compatible-approximation error controlled along the run.

\paragraph{An unknown-reward problem and its actor.}
There is one state and the action space is $(0,1)$. Fix
$h(a)=\sin(4\pi a)$. The unknown environment parameter is
$\xi\in\{-1,1\}$, and each reward, conditionally independently given its
action, is Bernoulli with mean $(1+\xi h(a))/2$. Thus the learner observes
noisy rewards and does not know which two action regions are optimal.
The discount factor $\gamma\in(0,1)$ is fixed. The optimal discounted value
is $V^*=1/(1-\gamma)$, attained by actions with $\xi h(a)=1$.

Consider the exponential family
\[
 \pi_x(a)=\exp\{xh(a)-\psi(x)\},\qquad
 \psi(x)=\log\int_0^1e^{xh(a)}\,da,\qquad x\in\mathbb R.
\]
For every nonzero $x$, this density has two modes. Its sufficient statistic $h(a)$ is also the reward-relevant quantity. Write $\mu(x)=\psi'(x)$ and $\sigma(x)^2=\psi''(x)$. The Fisher information in the $x$ coordinate is $\sigma(x)^2$, so we rescale distance by $\sigma(x)$ to obtain unit Fisher information. The resulting parameter is
\[
 \theta(x)=\int_0^x\sigma(u)\,du.
\]
Below we show that this is a smooth bijection of $\mathbb R$ onto itself.
We write $x(\theta)$ for its inverse and $\pi_\theta=\pi_{x(\theta)}$.
The endpoint score, Fisher information, advantage and population
coefficient are, respectively,
\begin{equation}\label{eq:joint-identities}
 \phi_\theta(a)=\frac{h(a)-\mu(x)}{\sigma(x)},\qquad
 F_\theta=1,\qquad
 A^{\pi_\theta}(a)=\frac{\xi}{2}(h(a)-\mu(x)),\qquad
 g_\theta=\frac{\xi\sigma(x)}2.
\end{equation}
In particular, compatibility is exact at every action and every parameter,
and $|g_\theta|\leq G:=1/2$.

This actor has an exact one-step DDIM realization. Let $H_x$ be its action
CDF and let $\Phi$ be the standard Gaussian CDF. Set
$D_\theta(z)=H_{x(\theta)}^{-1}(\Phi(z))$. This is a smooth diffeomorphism
from $\mathbb R$ onto the fixed domain $(0,1)$. In the convention
$a_0=\alpha_1^{-1/2}(z+\eta_1\epsilon_\theta(z))$, choose
$\alpha_1=1/2$, $\eta_1=-1/\sqrt2$, and
$\epsilon_\theta(z)=\sqrt2z-D_\theta(z)$. Then $a_0=D_\theta(z)$.
The score in \eqref{eq:joint-identities} can equivalently be evaluated from
the endpoint density or the DDIM change-of-variables identity. The integrals, CDF inversion, and scalar coordinate inversion specify this smooth actor and its score independently of the unknown reward sign.

For a unit vector $u\in\mathbb R^p$, $p\geq2$, the redundant family
$\pi_\vartheta=\pi_{u^\top\vartheta}$ has
$F_\vartheta=F_\vartheta^\dagger=uu^\top$.
The lift in Appendix~\ref{app:redundant-parameter-lift} shows that its
minimum-norm empirical updates, moment constants and regret agree with
the scalar learner, although its Fisher is singular.

\paragraph{Stochastic update.}
Initialize $\theta_0=0$. For $L\geq1$ batches, take $N=L$ independent
executions of Algorithm~\ref{alg:dppo-sampler}, with its accepted action, fair
Q/V branch, and geometrically terminated return. Fit its constrained
least-squares coefficient $\widehat g_i$ over $[-G,G]$, with minimum-norm
tie-breaking, and update
\[
 \theta_{i+1}=\theta_i+\eta\widehat g_i,\qquad \eta=L^{-1/2}.
\]
No parameter projection is used. The algorithm is not given $\xi$, the
population coefficient, the population reward mean, or the analysis
comparator. Let $T=LN=L^2$, and define the same batch pseudo-regret as in
the NPG theorem,
\[
 \operatorname{Reg}(T)=N\sum_{i=0}^{L-1}(V^*-V^{\pi_{\theta_i}}).
\]

\begin{prop}[Stochastic NPG for a two-mode actor]\label{prop:joint-stochastic-example}
For every fixed $\gamma\in(0,1)$ and $\delta\in(0,1)$, the stochastic
algorithm above satisfies, for either unknown reward sign,
\[
 \mathbb{P}\{\operatorname{Reg}(T)
       \leq C_{\gamma,\delta}T^{3/4}(1+\log T)\}\geq1-\delta,
\]
where $C_{\gamma,\delta}<\infty$ is independent of $L$.
The required second and fourth score moments, Fisher conditioning and compatible-coefficient
conditions hold globally. On an event of probability at least $1-\delta/2$,
a fixed stochastic comparison policy chosen before the samples are drawn has
bounded comparator score mean and satisfies
\eqref{eq:joint-average-curvature} along every actual update segment,
with constants independent of $L$. Its optimality gap is $O_\gamma(L^{-1/2})$
and its initial KL divergence is $O_{\gamma,\delta}(1+\log L)$.
\end{prop}

\begin{proof}
We first control the noise in the empirical NPG direction. This will bound the region visited by the iterates and let us choose a comparator whose curvature is controlled throughout that region.

\emph{Score and return moments.}
To obtain uniform score moments as the policy concentrates, we need the tail behavior of $\mu(x)$ and $\sigma(x)$. Under the uniform action distribution, $U=1-h(a)$ has density
$[\pi\sqrt{u(2-u)}]^{-1}$ on $(0,2)$. For $x>0$ define
\[
 I_k(x)=\int_0^2e^{-xu}u^{k-1/2}(2-u)^{-1/2}\,du.
\]
Then $\mathbb E_{\pi_x}U^k=I_k(x)/I_0(x)$. Substituting $t=xu$ and
expanding $(2-t/x)^{-1/2}$ on $0\leq t\leq x$ gives, for
$k=0,\ldots,4$,
\[
 I_k(x)=\frac{\Gamma(k+1/2)}{\sqrt2\,x^{k+1/2}}
 \left(1+\frac{k+1/2}{4x}+O(x^{-2})\right).
\]
For completeness, the expansion error on this interval is bounded by a
constant times $(t/x)^2$. Its integral is bounded by the corresponding
Gamma integral. On $x\leq t\leq2x$, the exponential contributes $e^{-x}$
and the integrable endpoint factor contributes only a polynomial in $x$.
This proves the displayed remainder. The raw moments through order four
therefore give
\begin{align*}
 1-\mu(x)&=\frac1{2x}+O(x^{-2}),&
 \sigma(x)^2&=\frac1{2x^2}+O(x^{-3}),\\
 \mathbb E_{\pi_x}(h-\mu(x))^3&=-\frac1{x^3}+O(x^{-4}),&
 \frac{\mathbb E_{\pi_x}(h-\mu(x))^4}{\sigma(x)^4}&\longrightarrow15.
\end{align*}
The exponential-family derivative identity
$2\sigma\sigma'=\mathbb E(h-\mu)^3$ gives
$\sigma'(x)=-1/(\sqrt2x^2)+O(x^{-3})$.
The law of $h$ is symmetric, so $\mu$ is odd, $\sigma$ is even, and the
same fourth-moment bound holds as $x\to-\infty$. On every compact interval,
$\sigma$ is positive and all these moments are continuous. Consequently
\[
 K_4:=\sup_{x\in\mathbb R}
 \frac{\mathbb E_{\pi_x}(h-\mu(x))^4}{\sigma(x)^4}<\infty.
\]
Thus $\beta_K=\lambda=1$, $M_F=K_4-1$ is valid, and
$\theta(x)\to\pm\infty$. Differentiation under the integral and likelihood
regularity hold locally at every finite parameter since $h$ is bounded.
The DDIM map is smooth at every finite $\theta$ and $z$ by the implicit
function theorem. This verifies the likelihood, invertibility and moment conditions in Assumption~\ref{aspt:ddpm-bijection}. We check comparator-averaged curvature below, after bounding the iterates.

The Q/V labels from Algorithm~\ref{alg:dppo-sampler} obey
$\mathbb E[\widehat A\mid a]=A^{\pi_\theta}(a)$ and
$\mathbb E[\widehat A^2\mid a]\leq C_A$ with
$C_A=4(1+\gamma)/(1-\gamma)^2$. Indeed $|\widehat A|\leq2H$ for a
geometric return length $H\geq1$, and
$\mathbb EH^2=(1+\gamma)/(1-\gamma)^2$. The initial occupancy sampling
phase does not change the accepted action law in the one-state problem.
Hence Lemma~\ref{lem:least-squares-npg} applies conditionally at
every parameter, even if an iterate later leaves the region used below:
\begin{equation}\label{eq:joint-coefficient}
 \mathbb E_i|e_i|\leq a_N:=a/\sqrt N+b/N,\qquad
 e_i=\widehat g_i-g_i,\quad
 a=2(\sqrt{C_A}+G\sqrt{M_F}),\quad b=8GM_F.
\end{equation}
Also $|e_i|\leq2G$ deterministically.

\emph{Controlling the iterates.}
The population direction points toward the rewarding modes. We bound the total perturbation from fitting it to noisy returns, then use that bound to control the stochastic trajectory. Let $\mathcal F_i$ be the history before batch $i$.
The variables $|e_i|-\mathbb E_i|e_i|$ are martingale differences bounded
in absolute value by $2G$. Azuma's inequality and
\eqref{eq:joint-coefficient} show that, with probability at least
$1-\delta/2$,
\begin{equation}\label{eq:joint-noise-event}
 \eta\sum_{i=0}^{L-1}|e_i|
 \leq a+b+2G\sqrt{2\log(2/\delta)}=:B.
\end{equation}
Denote this event by $\mathcal E$. The bound is uniform over the iterates because the preceding moments are global. To separate progress toward the rewarding modes from estimation noise, use the signed coordinate $y_i=\xi\theta_i$ and write
\[
 y_i=z_i+w_i,\qquad
 z_i=\frac\eta2\sum_{j<i}\sigma(x(\theta_j)),\qquad
 w_i=\eta\sum_{j<i}\xi e_j.
\]
On $\mathcal E$, $z_i\geq0$ and $|w_i|\leq B$ for every $i\leq L$.
Since $\theta$ is odd and $\sigma$ is even,
$\sigma(x(\theta_i))=\sigma(x(y_i))$.
Choose a fixed $X_0\geq1$ such that, for $x\geq X_0$,
\[
 \frac1{4x}\leq\sigma(x)\leq\frac2x,
 \qquad \sigma'(x)\leq0.
\]
Let $\theta_0^+=\theta(X_0)$ and
$A=e^{\theta_0^+/2}\max\{1,2/X_0\}$.
The upper variance bound implies
$\theta(x)\leq\theta_0^++2\log(x/X_0)$ for $x\geq X_0$.
Together with $\sigma\leq1$, this yields the global bound
\[
 \sigma(x(y))\leq A e^{-y/2},\qquad y\in\mathbb R.
\]
Therefore, on $\mathcal E$,
\[
 0\leq\Delta z_i\leq\frac\eta2,
 \qquad
 \Delta z_i\leq\frac{\eta A e^{B/2}}2 e^{-z_i/2}.
\]
Using $e^u-1\leq ue^u$ and $\eta\leq1$ gives
\[
 e^{z_{i+1}/2}-e^{z_i/2}
 \leq\frac{\eta A e^{B/2+1/4}}4=:\eta C_B.
\]
After summing, every actual endpoint obeys
\begin{equation}\label{eq:joint-envelope}
 -B\leq y_i\leq U:=B+2\log(1+C_B\sqrt L),\qquad 0\leq i\leq L.
\end{equation}
On $\mathcal E$, every update segment stays inside this deterministic interval. For the martingale argument below, we enlarge it to $I=[-B-G,U+G]$. Since $|\eta\widehat g_i|\leq G$, this larger interval also contains the whole first step that could leave $[-B,U]$.

\emph{Choosing the comparator.}
We need a comparator concentrated enough to have a small optimality gap, but with controlled log-density curvature along the iterates. Choose
\[
 X_*:=\max\{X_0,\sqrt L,x(U+G)\},\qquad q=\pi_{\xi X_*}.
\]
This comparator is fixed by $L,\gamma,\delta$ and the environment sign before the observations are drawn. Only the proof uses it.
The lower variance bound implies
$x(y)\leq X_0e^{4(y-\theta_0^+)}$ for $y\geq\theta_0^+$.
Consequently $X_*$ is at most a fixed power of $L$ times a finite
$(\gamma,\delta)$-dependent constant. The moment estimates imply a finite
constant $C_0$, independent of $L,\gamma,\delta$, such that
\[
 \epsilon_*:=V^*-V^q
 =\frac{1-\mu(X_*)}{2(1-\gamma)}
 \leq\frac{C_0}{(1-\gamma)\sqrt L}.
\]
Moreover $e^{\psi(x)}=e^xI_0(x)/\pi$ and the Laplace estimate give
\[
 \mathcal K_0=\mathrm{KL}(q\|\pi_0)
 =X_*\mu(X_*)-\psi(X_*)
 =\tfrac12\log X_*+O(1)
 =O_{\gamma,\delta}(1+\log L).
\]

For $y=\xi\theta\in I$, write $v=x(y)\leq X_*$ and
$v_-=x(-B-G)$. The comparator score mean satisfies
\[
 |\mathbb E_q\phi_\theta|
 =\frac{\mu(X_*)-\mu(v)}{\sigma(v)}
 \leq B_q,
\]
where the following constant is finite and independent of $L$:
\[
 B_q:=\max\left\{
 \sup_{v\geq X_0}\frac{1-\mu(v)}{\sigma(v)},
 \frac2{\min_{v\in[v_-,X_0]}\sigma(v)}\right\}.
\]
To bound curvature, differentiating twice under $q$ gives
\[
 f_q''(\theta)
 =-1-\frac{(\mu(X_*)-\mu(v))\sigma'(v)}{\sigma(v)^3}.
\]
For $v\geq X_0$ the second term is nonnegative. On the remaining compact
interval it is bounded below, so throughout $I$,
\[
 f_q''(\theta)\geq-L_q,\qquad
 L_q:=1+2\max_{v\in[v_-,X_0]}\frac{|\sigma'(v)|}{\sigma(v)^3}<\infty.
\]
Integrating this second-derivative bound proves \eqref{eq:joint-average-curvature} on every segment in $I$. The same formulas and compact-interval bounds give the required integrability under $q$.

\emph{Summing policy improvement.}
The KL argument applies throughout $I$. To use martingale concentration before knowing whether the iterates stay there, we stop the martingale at the first exit. Let $\tau$ be the first index at which $y_i\notin[-B,U]$, with
$\tau=L+1$ if there is no such index up to $L$. Set
$m_i=\mathbb E_q\phi_{\theta_i}$. The variables
\[
 Z_i=\mathbf1\{i<\tau\}\,
 m_i(\mathbb E_i\widehat g_i-\widehat g_i)
\]
are martingale differences with $|Z_i|\leq2B_qG$. The indicator is
$\mathcal F_i$-measurable. Thus, with probability at least $1-\delta/2$,
\[
 \sum_{i=0}^{L-1}Z_i
 \leq2B_qG\sqrt{2L\log(2/\delta)}.
\]
Intersecting this event with $\mathcal E$ gives probability at least $1-\delta$. On the intersection no exit occurs, all update segments satisfy the curvature inequality, and exact compatibility gives
\[
 (1-\gamma)(V^q-V^{\pi_{\theta_i}})=m_i g_i.
\]
Summing \eqref{eq:joint-average-curvature}, using nonnegative terminal KL,
and then adding and subtracting $\mathbb E_i\widehat g_i$ yields
\begin{align*}
 \operatorname{Reg}(T)\leq{}&T\epsilon_*+
 \frac{N}{1-\gamma}\left(\frac {\mathcal K_0}\eta+
                   \frac{\eta L_qG^2L}{2}\right)
 +\frac{TB_q a_N}{1-\gamma}\\
 &+\frac{2NB_qG}{1-\gamma}\sqrt{2L\log(2/\delta)}.
\end{align*}
Substituting $N=L$, $\eta=L^{-1/2}$, $T=L^2$, and the bounds on $\mathcal K_0$ and $\epsilon_*$ proves the stated high-probability regret bound. All displayed constants are finite for fixed $\gamma,\delta$ and independent of $L$.
\end{proof}

\paragraph{Environment interactions.}\label{app:example-interactions}
Each label uses the sampler's geometric occupancy phase (mean
$\gamma/(1-\gamma)$ transitions) and geometric return phase (mean
$1/(1-\gamma)$ reward transitions). The expected environment cost is
therefore $T(1+\gamma)/(1-\gamma)$ transitions, and the geometric coins give
a total cost $O_\gamma(T+\log(1/\delta))$ with high probability. The displayed
regret follows the paper's labeled-sample/batch convention. The same
exponent, up to logarithms, also holds when each auxiliary transition is
charged its deployed policy's discounted-value gap. Indeed, with probability
$1-\delta'$, both geometric phases of every label have length at most
$h=\lceil\log(2T/\delta')/\log(1/\gamma)\rceil$, by a union bound.
The transition count is then at most $2hT$, and the corresponding
transition pseudo-regret is at most $2h\operatorname{Reg}(T)$. Intersecting
this event with the proposition's event gives probability at least
$1-\delta-\delta'$. This is discounted-value pseudo-regret, not a claim
about undiscounted realized-reward regret.


\subsection{A two-layer actor with trainable hidden slopes and output weights}
\label{app:two-layer-shape-example}

We next verify the learning conditions for a finite neural actor whose hidden and output weights both change during training. This addresses two issues left by the analytic two-mode example: its inverse-CDF map is not a finite network, and its exact compatibility avoids having to control a nonzero advantage-approximation error. Here a three-unit tanh network satisfies the original pointwise smoothness and positive-Fisher conditions globally, and its cumulative transferred approximation error has the rate required for sublinear regret.

We use a location coordinate to move toward the rewarding action and a shape coordinate to alter the centered distribution. The hidden slopes and output amplitudes are linked through a periodic parameterization, keeping the shape derivatives and Fisher conditioning uniform for unrestricted parameter updates. Thus the result concerns this specified neural family, rather than arbitrary independent changes to every raw weight.

For $z\in\mathbb R$, define
\[
 B_w(z)=\tanh(wz+1)+\tanh(wz-1)-2\tanh(wz),
\]
and parameterize the hidden slope and output amplitude by
\[
 w(\varphi)=1+\tfrac14\cos\varphi,\qquad
 v(\varphi)=\tfrac1{64}+\tfrac1{128}\sin\varphi.
\]
The two-layer residual actor, with two unconstrained trainable coordinates
$\theta=(c,\varphi)\in\mathbb R^2$, is
\begin{equation}\label{eq:two-layer-actor}
 T_\varphi(z)=z+v(\varphi)B_{w(\varphi)}(z),\qquad
 A_\theta(z)=\tanh(c+T_\varphi(z)),\qquad Z\sim N(0,1).
\end{equation}
It has three hidden tanh units, a residual connection, and a tanh output.
The hidden slopes $w(\varphi)$ and hidden-to-output weights
$(v(\varphi),v(\varphi),-2v(\varphi))$ are trainable. It is exactly a
one-step DDIM actor with $\alpha_1=1/2$, $\eta_1=-1/\sqrt2$ and
$\epsilon_\theta(z)=\sqrt2z-A_\theta(z)$. This gives a finite two-layer residual MLP denoising field for the policy-optimization parameterization.

The environment is the one-state discounted MDP with actions $[-1,1]$
and conditionally independent Bernoulli rewards of unknown mean
$r_\xi(a)=(1+\xi a)/2$, where $\xi\in\{-1,1\}$. Its optimal value
is $1/(1-\gamma)$ at the available boundary action $a=\xi$. The actor
and analysis comparator have common open support $(-1,1)$; all likelihood
derivatives below are on that support.

\paragraph{Regularity and Fisher conditioning.}
Write $R_\varphi=T_\varphi-\mathrm{id}$ and $Q_\varphi=T_\varphi'$.
Since $|B_w'|\leq4w$,
\[
 |R_\varphi'|\leq4\cdot\frac3{128}\cdot\frac54
 =\frac{15}{128}<\frac14,
 \qquad \frac34<Q_\varphi<\frac54.
\]
Thus every $T_\varphi$ is a smooth diffeomorphism of the line, and every
$A_\theta$ is a diffeomorphism onto $(-1,1)$. Set
$D_0=\sup_{\varphi,z}|R_\varphi(z)|<\infty$.
The residual and all mixed derivatives of any fixed finite order in
$z,\varphi$ decay as a polynomial in $|z|$ times
$e^{-3|z|/2}$, uniformly in $\varphi$. To see this, expand the three tanh
terms at $z\to\pm\infty$. Their constant terms cancel, $w\geq3/4$,
and differentiating with respect to $w$ introduces only polynomial
factors in $z$. Derivatives in $\varphi$ have bounded coefficients.

For a fixed action let $x=\operatorname{atanh}(a)$ and
$z=T_\varphi^{-1}(x-c)$. The endpoint parameter scores are
\begin{align*}
 s_c(z,\varphi)&=\frac z{Q_\varphi(z)}
                  +\frac{Q_\varphi'(z)}{Q_\varphi(z)^2},\\
 s_\varphi(z,\varphi)&=z u_\varphi(z)-u_\varphi'(z),
 \qquad u_\varphi(z)=\frac{\partial_\varphi T_\varphi(z)}{Q_\varphi(z)}.
\end{align*}
These follow by differentiating
$-z^2/2-\log Q_\varphi(z)-\log(1-a^2)-\log\sqrt{2\pi}$ at fixed $a$,
using $\partial_cz=-1/Q_\varphi$ and $\partial_\varphi z=-u_\varphi$.
The preceding tail bounds give universal finite constants $C_s,L_s$ with
\[
 |s_c(z,\varphi)|\leq C_s(1+|z|),\qquad
 |s_\varphi(z,\varphi)|\leq C_s,
 \qquad \|\nabla_\theta^2\log\pi_\theta(a)\|_{\rm op}\leq L_s
\]
uniformly in $\theta,a$. For the last statement, differentiate these
score formulas with the two fixed-action inverse derivatives above:
every nonconstant term is a polynomial in $z$ times an exponentially
decaying residual derivative, divided by a power of $Q_\varphi\geq3/4$.
The remaining derivative is the constant Gaussian location curvature.
This verifies the original pointwise smoothness premise globally.

Oddness of $T_\varphi$ makes $s_c$ odd and $s_\varphi$ even as functions
of $z$. Therefore, under the sampled standard Gaussian,
\[
 F_\theta=\operatorname{diag}(I_c(\varphi),I_\varphi(\varphi)),\qquad
 I_c=\mathbb E s_c^2,\quad I_\varphi=\mathbb E s_\varphi^2.
\]
Both diagonal entries are uniformly positive. For location, integration
by parts gives $\mathbb E[T_\varphi(Z)s_c(Z,\varphi)]=1$, so
$I_c\geq1/\mathbb E T_\varphi(Z)^2>0$ uniformly. If $I_\varphi=0$,
continuity of the score would give $u_\varphi'=z u_\varphi$ everywhere.
Its only exponentially decaying solution is zero, which would force
\[
 v'(\varphi)B_w+v(\varphi)w'(\varphi)\partial_wB_w=0.
\]
The two functions here are linearly independent: as $z\to+\infty$,
$B_w(z)=C e^{-2wz}(1+o(1))$ with
$C=-4(\cosh2-1)\neq0$, whereas
$\partial_wB_w(z)/B_w(z)=-2z+o(1)$. Also $v>0$ and $v',w'$ never
vanish simultaneously. This is a contradiction. Continuity and
$2\pi$-periodicity now give constants
\[
 0<\lambda\leq I_c(\varphi),I_\varphi(\varphi)\leq\Lambda<\infty.
\]
The Gaussian score bounds above give uniform finite second and fourth
moments. Fix any resulting bounds $\mathbb E\|s\|^2\leq\beta^2$ and
$\mathbb E\|ss^\top-F\|_F^2\leq M_F$.
All constants in this paragraph are determined by the displayed fixed
network; they are not additional assumptions on the learning trajectory.

At $\varphi=0$, the derivative of the latent variance satisfies
\[
 \left.\frac{\mathrm d}{\mathrm d\varphi}\operatorname{Var}(T_\varphi(Z))\right|_{\varphi=0}
 =2v'(0)\mathbb E[T_0(Z)B_{w(0)}(Z)]<0.
\]
Indeed, for $t=\tanh(wz)$
and $r=\tanh1$,
\[
 B_w(z)=-\frac{2t r^2(1-t^2)}{1-t^2r^2},
\]
which has sign opposite to $z$, while $T_0$ has the sign of $z$.
Hence changing the trained shape coordinate changes the centered latent
law, not merely its location.

\paragraph{Stochastic NPG update.}
Initialize $(c_0,\varphi_0)=(0,0)$. In each of $L$ batches take $N=L$
independent executions of the original reset sampler, including its fair
Q/V coin and geometric noisy returns. Fit the minimum-norm least-squares
coefficient over $\|g\|\leq G:=\lambda^{-1/2}$ and update
$\theta_{i+1}=\theta_i+\eta\widehat g_i$. The population coefficient
belongs to this ball because $|A^{\pi_\theta}|\leq1$ in the one-state
problem and the normal equations give
$\|g_\theta\|^2\leq\mathbb E(A^{\pi_\theta})^2/\lambda$.
There is no projection of $c$ or $\varphi$; the periodic weight
parameterization is part of the fixed actor definition.

Let $T=L^2$, $C_A=4(1+\gamma)/(1-\gamma)^2$, and set
\[
\begin{gathered}
 A_0=\frac{2(\sqrt{C_A}\beta+G\sqrt{M_F})}{\lambda},\quad
 B_0=\frac{8GM_F}{\lambda^2},\quad
 C_\delta=2G\sqrt{2\log(2/\delta)},\\
 \eta=\min\left\{1,\frac1{(A_0+C_\delta)\sqrt L+B_0}\right\}.
\end{gathered}
\]
The original coefficient lemma applies globally and conditionally at
every actual iterate:
\[
 \mathbb E_i\|\widehat g_i-g_i\|\leq a_N:=A_0/\sqrt N+B_0/N,
 \qquad \|\widehat g_i-g_i\|\leq2G.
\]
The return moment used here is the same $C_A$ bound, since
$|\widehat A|\leq2H$ and
$\mathbb EH^2=(1+\gamma)/(1-\gamma)^2$. These statements do not
condition on the future path event. Azuma gives an event $\mathcal E$
of probability at least $1-\delta/2$ on which
\begin{equation}\label{eq:two-layer-noise-budget}
 \eta\sum_{i<L}\|\widehat g_i-g_i\|\leq1.
\end{equation}

\begin{prop}[Stochastic NPG for a two-layer neural actor]\label{prop:two-layer-shape-learning}
For the specified network and either unknown reward sign, the actual
stochastic NPG algorithm satisfies
\[
 \mathbb{P}\left\{\operatorname{Reg}(T)\leq
 C\left[\frac1{(1-\gamma)^2}
       +\frac{\sqrt{\log(2/\delta)}}{1-\gamma}\right]
 T^{3/4}(1+\log T)^2\right\}\geq1-\delta,
\]
where $C<\infty$ depends only on the displayed fixed architecture.
The original pointwise score smoothness, sampling moment and positive
Fisher spectrum conditions hold globally. The transferred squared
advantage error is controlled cumulatively along the actual run.
The same example has a singular Fisher after any fixed redundant
orthonormal-row lift of its two coordinates.
\end{prop}

\begin{proof}
Block diagonal Fisher information is crucial. The population location
direction is exactly
\[
 (g_i)_c=\frac{\xi}{2I_c(\varphi_i)}
 \mathbb E\operatorname{sech}^2(c_i+T_{\varphi_i}(Z)).
\]
Thus it always points toward the unknown rewarding boundary. Write
$y_i=\xi c_i=z_i+w_i$, where
$z_i=\eta\sum_{j<i}\xi(g_j)_c\geq0$ and
$w_i=\eta\sum_{j<i}\xi(\widehat g_j-g_j)_c$.
On $\mathcal E$, $|w_i|\leq1$ for every $i\leq L$.
Since $|T_\varphi(Z)-Z|\leq D_0$, the elementary sech bounds give
universal constants $a,b>0$ such that, uniformly over every shape path,
\[
 1+a\eta i\leq e^{2z_i}\leq1+b\eta i,
 \qquad e^{-2y_i}\leq\frac{e^2}{1+a\eta i},\qquad
 -1\leq y_i\leq U:=1+\tfrac12\log(1+b\eta L).
\]
For explicit choices, take $c_G=\mathbb E e^{-2|Z|}>0$,
$K_-=c_Ge^{-2D_0-2}/(2\Lambda)$,
$K_+=2e^{4+2D_0}/\lambda$,
$a=2K_-$ and $b=2K_+e^{1/\lambda}$.
Indeed $\eta K_-e^{-2z_i}\leq\Delta z_i\leq\eta K_+e^{-2z_i}$,
and $\Delta z_i\leq\eta/(2\lambda)$; exponentiating and summing gives
the inequalities. No bound on the angular iterate is needed.

The value gap is at most $C e^{-2y_i}/(1-\gamma)$. In addition the
shape population direction obeys
\[
 |(g_i)_\varphi|
 =\frac{|\mathbb E[\operatorname{sech}^2(c_i+T_{\varphi_i}(Z))
                  \partial_\varphi T_{\varphi_i}(Z)]|}
        {2I_\varphi(\varphi_i)}
 \leq C e^{-2y_i},
\]
because $\partial_\varphi T$ is uniformly bounded and both exponential
moments of $T_\varphi(Z)$ are bounded. The same upper bound holds for
the location direction. Hence $\|g_i\|\leq C e^{-2y_i}$ on the event.

Choose the fixed comparator $q=\pi_{(\xi t_*,0)}$ with
$t_*=1+\tfrac12\log[L(1+b)]\geq U$. Its optimality gap is
$\epsilon_*:=V^*-V^q\leq C/[(1-\gamma)L]$.
To bound initial KL, write a sample before its outer squash as
$x=\xi t_*+T_0(Z)$, and put $z_0=T_0^{-1}(x)$.
Then $|z_0-\xi t_*-Z|\leq2D_0$ and
\[
 \log\frac{q(A)}{\pi_{(0,0)}(A)}
 =\tfrac12(z_0^2-Z^2)+\log Q_0(z_0)-\log Q_0(Z).
\]
Consequently $\mathcal K_0=\mathrm{KL}(q\|\pi_{(0,0)})\leq C(1+t_*^2)$.

At an actual parameter with $-1\leq y_i\leq U$, the latent inverse
under a comparator sample obeys
$|z_i-\xi(t_*-y_i)-Z|\leq2D_0$. Thus the score formulas give
\[
 \|s_{\theta_i}\|_{L^2(q)}\leq C(t_*+2),\qquad
 \|\mathbb E_qs_{\theta_i}\|\leq B_*:=C(t_*+2).
\]
For the actual squared transfer error
$\epsilon_i=\mathbb E_q(A^{\pi_{\theta_i}}-s_{\theta_i}^\top g_i)^2$,
use
\[
 \|1-r_\xi(A)\|_{L^2(q)}\leq C e^{-2t_*},\qquad
 1-\mathbb E_{\pi_{\theta_i}}r_\xi\leq C e^{-2y_i}.
\]
The first bound follows by integrating the squared loss using Gaussian exponential moments. Combining these bounds with the
coefficient and comparator-score estimates yields
\[
 \sqrt{\epsilon_i}\leq
 C[(t_*+3)e^{-2y_i}+e^{-2t_*}],\qquad
 \sum_{i<L}\sqrt{\epsilon_i}\leq
 C\left[(t_*+3)\left(1+\frac{\log(1+a\eta L)}{a\eta}\right)+1\right].
\]
The transferred error therefore has the cumulative bound required by the NPG analysis.

Stop the analysis at the first endpoint with $y_i\notin[-1,U]$.
The predictable stopped comparator-score martingale has increments
bounded by $2B_*G$ and hence sum at most
$2B_*G\sqrt{2L\log(2/\delta)}$ with probability $1-\delta/2$.
Intersect with $\mathcal E$, on which no stop occurs. The unchanged
pointwise-smooth KL proof, retaining the batchwise squared errors, gives
\begin{align*}
 \operatorname{Reg}(T)\leq{}&T\epsilon_*+
 \frac N{1-\gamma}\left(\frac {\mathcal K_0}\eta+
       \frac{\eta L_sG^2L}{2}+\sum_{i<L}\sqrt{\epsilon_i}\right)\\
 &+\frac{TB_*a_N}{1-\gamma}
 +\frac{2NB_*G}{1-\gamma}\sqrt{2L\log(2/\delta)}.
\end{align*}
Here $\eta^{-1}=O([(1-\gamma)^{-1}+\sqrt{\log(2/\delta)}]\sqrt L+1)$
and $t_*=O(1+\log L)$, proving the proposition. The direct value-gap
sum also gives the same rate with one fewer logarithm.
\end{proof}

\paragraph{Shape learning and redundant parameters.}
The population update changes the actor's shape as well as its location. At $\varphi=0$ and
$\xi c>0$, pairing positive and negative Gaussian arguments shows
$\partial_\varphi\mathbb E r_\xi(A)>0$: $v'(0)>0$, $B_w(z)<0$ for
$z>0$, and the difference of the two sech-squared factors has the
opposite sign. Thus policy improvement changes the shape coordinate.
Both $v(\varphi)$ and $w(\varphi)$ are nonconstant, so the learnable
manifold changes hidden and output weights, not merely the final bias.
For $(c,\varphi)^\top=P\vartheta$ with $PP^\top=I_2$ and
$P\in\mathbb R^{2\times p}$, $p>2$, Appendix~\ref{app:redundant-parameter-lift}
gives the same empirical updates with a rank-two Fisher. Thus the
construction also permits redundant trainable coordinates.

The environment-transition cost and discounted-value pseudo-regret bounds
in Appendix~\ref{app:example-interactions} apply unchanged to this learner.

\subsection{Miscellaneous Lemmas}

\begin{lem}[Diffusion Policy Score]
For a DDPM path with a parameter-independent Gaussian initial latent and
positive transition densities,
\[
\nabla_\theta\log\pi_\theta(a_0\mid s)
=\E\left[\left.\sum_{k=1}^K
\nabla_\theta\log\pi_\theta(a_{k-1}\mid a_k,s)
\right|s,a_0\right].
\]
The expectation is under the posterior law of the denoising path conditional
on the endpoint, not its unconditional forward sampling law.
\label{lem:diffusion-policy-score}
\end{lem}

\begin{proof}
Let $p_\theta(a_{0:K}\mid s)=\varphi(a_K)
\prod_{k=1}^K\pi_\theta(a_{k-1}\mid a_k,s)$, where $\varphi$ is the standard
Gaussian density. Recall that we can write $\pi_\theta(a_0\mid s)$ as the equivalent integral over all possible paths, integrating $p_\theta$ against Lebesgue measure over $a_{1:K}$. Under the differentiation and integrability conditions of Assumption~\ref{aspt:ddpm-bijection}, using the log-derivative rule, we obtain that
\begin{align*}
\nabla_\theta\log\pi_\theta(a_0\mid s)
&=\frac1{\pi_\theta(a_0\mid s)}
\int p_\theta(a_{0:K}\mid s)\nabla_\theta\log p_\theta(a_{0:K}\mid s)
\,\mathrm da_{1:K}\\
&=\int p_\theta(a_{1:K}\mid s,a_0)
\sum_{k=1}^K\nabla_\theta\log\pi_\theta(a_{k-1}\mid a_k,s)
\,\mathrm da_{1:K}.
\end{align*}
This is the stated conditional expectation. The initial Gaussian contributes
zero to the parameter derivative.
\end{proof}


% As we have defined 
%     $$d_\nu^\pi(s, a):=(1-\gamma) \mathbb{E}_{s_0, a_0 \sim \nu}\left[\sum_{t=0}^{\infty} \gamma^t \mathbb{P}^\pi\left(s_t=s, a_t=a \mid s_0, a_0\right)\right] \geq (1-\gamma)\nu,$$
%     it follows that $(1-\gamma)\nu \leq d_\nu^{(t)} = d^{(t)}$ for all $t=1,...,T$. We can then see that 
%     $$\left\lVert \frac{d^*}{d^{(t)}} \right\rVert_{\infty} \leq \frac{1}{1-\gamma}\left\lVert \frac{d^*}{\nu} \right\rVert_{\infty} \leq \frac{c_{\text{npg}}^*}{1-\gamma}, \text{ by Assumption \ref{aspt:npg-concentrability}}.$$  
    
    % \kevin{We could use this if we bound the prediction error under the current occupancy measure. But seems like we don't need to do so. }

% We will be sampling on-policy from the diffusion policy. This incurs sampling error from the sampling target because we cannot sample from the ``score'' directly -- we have to use DDPM or DDIM to do so. Additionally, within the NPG update, we will be using $\widehat{\pi}^{(t)}$ instead of $\pi^{(t)}$. This is simply policy gradient on the family of diffusion policies where we denoise for $K$ timesteps. 

% Suspect that $\epsilon_{\text{bias}}$ is decreasing in $K$. Policy class/advantage function estimation complexity should show up in $B$ and $W$. 



% \subsection{Diffusion actor-critic with NPG updates}

% \kevin{Agarwal and friends give us an inexact update for NPG. Fit a critic, sample 1/eps2 on-policy samples. Don't need log-linear policy class, but need smoothness. Look at section 6.3, for the sample-based algorithm, in the Agarwal paper. }


% \begin{algorithm}[h]
%     \caption{Optimistic Natural Diffusion Policy Actor-Critic}
%     \begin{algorithmic}[1]
%             \STATE {\bfseries Input:} Pretrained diffusion policy $\pi_\theta$ with offline dataset $\gD_{\off}$ via behavior cloning, batch size $N$, within-batch learning rate $\alpha$, between-batch learning rate $\eta$, denoising steps $K$, discount factor $\gamma$. Function class $\gF$. 
%             \STATE {\bfseries Initialize:} Rollout buffer $\gD_{\text{on}} \gets \emptyset$, $t_{\last} \gets 0$, $\bg^{(0)} \gets 0$.  
%             \FOR{episode $t = 1, 2, \dots, T$}
%             \STATE Set $\widehat{Q}^{\pi_{\theta^{(t)}}}, \widehat{V}^{\pi_{\theta^{(t)}}} = 0$.
%             \STATE Sample $s_1^{(t)} \sim \nu$, $a_1^{(t)} \sim \nu(\cdot | s_1^{(t)})$.
%             \FOR{horizon $h=1,...$}
%             \IF[\textcolor{blue}{sample from discounted occupancy measure}]{$\text{Ber}(\gamma) = 1$}
%             \STATE Execute $a_h^{(t)}$, draw $s_{h+1}^{(t)}, r_h^{(t)}$.
%             \STATE Sample next denoised action $a_{h+1,k}^{(t)} \sim \pi_{\theta^{(t)}}(\cdot | a_{h+1,k+1}^{(t)}, s_{h+1}^{(t)})$ for $k=K-1,...,0$, set $a_{h+1}^{(t)} \gets a_{h+1,0}^{(t)}$.
%             \ELSE 
%             \STATE Accept $(s_h^{(t)}, a_h^{(t)})$ as a sample from discounted occupancy measure $d_\nu^{\pi_{\theta^{(t)}}}$, append to $\gD_{\on}.$
%             \ENDIF
%             \STATE Set  $\widehat{A}^{\pi_{\theta^{(t)}}} \gets f^{(t)}(s_h^{(t)}, a_h^{(t)}) - \E_{a \sim \pi_{\theta^{(t)}}} f^{(t)}(s_h^{(t)}, a) $.
%             \STATE Update for $\gG = \left\{g\in \R^p \mid ||g||_2 \leq G\right\}$: \COMMENT{\textcolor{blue}{SGD-type update on the running gradient estimate}}
%             $$\bg^{(t)} \gets \text{Proj}_{\gG}\left(\bg^{(t-1)} - 2\alpha\left(\left\langle \bg^{(t-1)}, \nabla_\theta \log\pi_{\theta^{(t)}}(a_h^{(t)} \mid s_h^{(t)})\right\rangle - \widehat{A}^{\pi_{\theta^{(t)}}}\right) \nabla_\theta \log\pi_{\theta^{(t)}}(a_h^{(t)} \mid s_h^{(t)})\right).$$
%             \ENDFOR
%             \IF{$t \mod N = 0$}
%             \STATE Set $\widehat{\bg}^{(t)} \gets \frac{1}{N}\sum_{t'=t_{\last}}^t \bg^{(t)}$, update $\theta^{(t+1)} \gets \theta^{(t)} + \eta \cdot \widehat{\bg}^{(t)}$, set $t_{\last} \gets t$. 
%             \STATE Update critic by computing confidence set $\gF^{(t)}$:
%             \begingroup
%             \addtolength\jot{-1ex}
%             \begin{align*}
%                 &\mathcal{F}^{(t)} \leftarrow\left\{f \in \mathcal{F}: \mathcal{L}^{(t)}\left(f_h, f_{h+1}\right)-\min _{f_h^{\prime} \in \mathcal{F}_h} \mathcal{L}^{(t)}\left(f^{\prime}, f\right) \leq (1-\gamma)^2\beta \right\},\\
%                 &\mathcal{L}_h^{(t)}\left(f, f'\right) \gets\sum_{\left(s, a, r, s^{\prime}\right) \in \mathcal{D}}\left(f(s, a)-r-\gamma f^{\prime}(s^{\prime}, \pi_{\theta^{(t)}}(s'))\right)^2.
%             \end{align*}
%             \endgroup
%         \vspace{-2ex}
%             \ELSE 
%             \STATE Set $\theta^{(t+1)} \gets \theta^{(t)}$, and $f^{(t+1)} \gets f^{(t)}$.
%             \ENDIF
%         \ENDFOR
%     \end{algorithmic}
% \label{alg:dppo-ac}
% \end{algorithm}



\section{Analysis of \aigen{Algorithm~\ref{alg:dipo-bc} (Continuous Score-Matching Analysis)}}
\label{app:dipo-bc}
In this section, we analyze \aigen{Algorithm~\ref{alg:dipo-bc}} by decomposing the regret of the algorithm into two parts: diffusion policy behavior cloning error and critic estimation error. Then, diffusion policy behavior cloning error and critic estimation error are analyzed separately and combined into the final regret bound. Initial-state expressions below are averaged over $s_1^{(t)}\sim\rho$, so that the regret agrees with the initial-distribution definition. The critic $f^{(t)}$ uses data from episodes $i<t$. Its confidence set is therefore $\gF^{(t-1)}$, with $\gF^{(0)}=\gF$.

\subsection{Regret Decomposition}
\begin{lem}[Regret Decomposition for Algorithm \ref{alg:dipo-bc}]
    The following regret decomposition holds for Algorithm \ref{alg:dipo-bc}:
    \begin{align*}
        \text{Reg}(T) &= \sum_{t=1}^T\left(V_1^{*}(s_1^{(t)})-V_1^{\pi^{(t)}}(s_1^{(t)})\right) \\
        &\leq \underbrace{2\sum_{t=1}^T\sum_{h=1}^H \E_{d_h^{(t)}}\left[\left\langle f_h^{(t)}(s_h, \cdot), \pi_h^{f^{(t)}}(\cdot \mid s_h)-\pi_h^{(t)}(\cdot \mid s_h)\right\rangle\right]}_{\textcolor{red}{\text{Diffusion policy behavior cloning error.}}} + \underbrace{\sum_{t=1}^T \sum_{h=1}^H \mathbb{E}_{d_h^{(t)}}\left[f_h^{(t)}-\gT_h f_{h+1}^{(t)}\right]}_{\textcolor{blue}{\text{Critic estimation error.}}} .
    \end{align*}
    \label{lem:regret-decomp-dipo}
\end{lem}
\begin{proof}
We first add and subtract $f_1^{(t)}(s_1^{(t)}, \pi_1^{f^{(t)}}(s_1^{(t)}))$, obtaining
    \begin{align*}
        \text{Reg}(T) &= \sum_{t=1}^T\left(V_1^{*}(s_1^{(t)})-V_1^{\pi^{(t)}}(s_1^{(t)})\right) \\
        &= \sum_{t=1}^T \left(V_1^{*}(s_1^{(t)})-f_1^{(t)}(s_1^{(t)}, \pi_1^{f^{(t)}}(s_1^{(t)})\right) + \sum_{t=1}^T\left(f_1^{(t)}(s_1^{(t)}, \pi_1^{f^{(t)}}(s_1^{(t)})-V_1^{\pi^{(t)}}(s_1^{(t)})\right) .
    \end{align*}

    The first term is non-positive by optimistic selection and $Q^* \in \gF^{(t-1)}$ from Lemma~\ref{lem:optimism-in-sample-error}. It remains to address the second term. Applying the generalized policy difference lemma from \citep{efroni2020optimisticpolicyoptimizationbandit, cai2024provablyefficientexplorationpolicy} (presented in Lemma \ref{ch05:lem:policy-value-difference-lemma} for completeness) with $\pi = \pi^{f^{(t)}}, \pi' = \pi^{(t)}$, and with $f^{(t)}$ as the Q-function, followed by $\gT_h^{\pi^{f^{(t)}}}f_{h+1}^{(t)}\ge\gT_h^{\pi^{(t)}}f_{h+1}^{(t)}$, yields the inequality:
    \begin{align*}
        &\sum_{t=1}^T\left(f_1^{(t)}(s_1^{(t)}, \pi_1^{f^{(t)}}(s_1^{(t)})-V_1^{\pi^{(t)}}(s_1^{(t)})\right)  \\
        &\leq \underbrace{\sum_{t=1}^T\sum_{h=1}^H \E_{d_h^{(t)}}\left[\left\langle f_h^{(t)}\left(s_h, \cdot\right), \pi_h^{f^{(t)}}\left(\cdot \mid s_h\right)-\pi_h^{(t)}\left(\cdot \mid s_h\right)\right\rangle\right]}_{\textcolor{red}{\text{Diffusion policy behavior cloning error.}}} + \underbrace{\sum_{t=1}^T \sum_{h=1}^H \mathbb{E}_{d_h^{(t)}}\left[f_h^{(t)}-\gT_h^{\pi^{(t)}}f_{h+1}^{(t)}\right]}_{\textcolor{blue}{\text{Critic estimation error.}}} .
    \end{align*}
    The first term corresponds to the behavior cloning error, while the second corresponds to critic estimation error. However, we are not done yet. The second term, at least in this specific form, is difficult to deal with in general without making a very strong value completeness-like assumption on the function class \citep{liu2023optimisticnaturalpolicygradient}. We therefore further decompose the term:
    \begin{align*}
        &\sum_{t=1}^T \sum_{h=1}^H \mathbb{E}_{d_h^{(t)}}\left[f_h^{(t)}-\gT_h^{\pi^{(t)}}f_{h+1}^{(t)}\right]
        \\
        &= \sum_{t=1}^T \sum_{h=1}^H \mathbb{E}_{d_h^{(t)}}\left[f_h^{(t)}-\gT_hf_{h+1}^{(t)}\right] + \sum_{t=1}^T \sum_{h=1}^H \mathbb{E}_{d_h^{(t)}}\left[\gT_h f_{h+1}^{(t)}-\gT_h^{\pi^{(t)}}f_{h+1}^{(t)}\right] \\
        &= \sum_{t=1}^T \sum_{h=1}^H \mathbb{E}_{d_h^{(t)}}\left[f_h^{(t)}-\gT_hf_{h+1}^{(t)}\right] + \sum_{t=1}^T\sum_{h=2}^H \E_{d_h^{(t)}}\left[\left\langle f_h^{(t)}(s_h,\cdot),\pi_h^{f^{(t)}}(\cdot\mid s_h)-\pi_h^{(t)}(\cdot\mid s_h)\right\rangle\right],
    \end{align*}
    where expectations involving $s_{h+1}$ include the transition from $(s_h,a_h)$. The last line follows from the observation that
\begin{align*}
& \mathbb{E}_{\pi^{(t)}}\left[\left(\gT_h f_{h+1}^{(t)} -\gT_h^{\pi^{(t)}}f_{h+1}^{(t)}\right)\left(s_h, a_h\right)\right]\\
        &= \mathbb{E}_{d_h^{(t)}}\left[r_h(s_h, a_h) + f_{h+1}^{(t)}\left(s_{h+1}, \pi_{h+1}^{f^{(t)}}(s_{h+1})\right) - r_h(s_h, a_h) - f_{h+1}^{(t)}\left(s_{h+1}, \pi_{h+1}^{(t)}(s_{h+1})\right) \right] \\
        &= \mathbb{E}_{d_h^{(t)}}\left[f_{h+1}^{(t)}\left(s_{h+1}, \pi_{h+1}^{f^{(t)}}(s_{h+1})\right) - f_{h+1}^{(t)}\left(s_{h+1}, \pi_{h+1}^{(t)}(s_{h+1})\right) \right] \\
        &=  \mathbb{E}_{d_h^{(t)}}\left[\left\langle f_{h+1}^{(t)}(s_{h+1},\cdot), \pi_{h+1}^{f^{(t)}} (\cdot | s_{h+1}) -\pi_{h+1}^{(t)}(\cdot | s_{h+1}) \right\rangle\right].
    \end{align*}
    
    Each greedy-action gap is nonnegative. Thus the sum over $h=2,\ldots,H$ is at most the sum over $h=1,\ldots,H$, and combining the bounds gives:
    \begin{align*}
        &\sum_{t=1}^T\left(f_1^{(t)}(s_1^{(t)}, \pi_1^{f^{(t)}}(s_1^{(t)})-V_1^{\pi^{(t)}}(s_1^{(t)})\right)  \\
        &\leq \underbrace{2\sum_{t=1}^T\sum_{h=1}^H \E_{d_h^{(t)}}\left[\left\langle f_h^{(t)}\left(s_h, \cdot\right), \pi_h^{f^{(t)}}\left(\cdot \mid s_h\right)-\pi_h^{(t)}\left(\cdot \mid s_h\right)\right\rangle\right]}_{\textcolor{red}{\text{Diffusion policy behavior cloning error.}}} + \underbrace{\sum_{t=1}^T \sum_{h=1}^H \mathbb{E}_{d_h^{(t)}}\left[f_h^{(t)}-\gT_h f_{h+1}^{(t)}\right]}_{\textcolor{blue}{\text{Critic estimation error.}}} .
    \end{align*}

\end{proof}

\subsection{Critic Estimation Error}

\begin{lem}[Critic Estimation Error]
    With probability at least $1-\delta$, 
    $$\sum_{t=1}^T \sum_{h=1}^H\mathbb{E}_{d_h^{(t)}}\left[\left(f_h^{(t)}-\gT_h f_{h+1}^{(t)}\right)\left(s_h, a_h\right)\right] \leq O(\sqrt{\beta H^4 T\cdot \mathsf{SEC}(\gF, \Pi, T)}).$$
    \label{lem:critic-estimation-error}
\end{lem}
\begin{proof}
    Write $\delta_h^{(t)}=f_h^{(t)}-\gT_hf_{h+1}^{(t)}$. The argument follows the proof of Theorem 12 in \citet{xie2022role}:
\begin{align*}
    &\sum_{t=1}^T \sum_{h=1}^H\mathbb{E}_{d_h^{(t)}}\left[\left(f_h^{(t)}-\gT_h f_{h+1}^{(t)}\right)\left(s_h, a_h\right)\right] \\
    &= \sum_{t=1}^T \sum_{h=1}^H \E_{d_h^{(t)}}[\delta_h^{(t)}]\\
    &= \sum_{t=1}^T \sum_{h=1}^H \E_{d_h^{(t)}}[\delta_h^{(t)}] \left(\frac{H^2 \vee \sum_{i=1}^{t-1}\E_{d_h^{(i)}}[(\delta_h^{(t)})^2]}{H^2 \vee \sum_{i=1}^{t-1}\E_{d_h^{(i)}}[(\delta_h^{(t)})^2]}\right)^{1/2} \\
    &\leq \sqrt{\sum_{t=1}^T\sum_{h=1}^H \frac{\E_{d_h^{(t)}}[\delta_h^{(t)}]^2}{H^2 \vee \sum_{i=1}^{t-1}\E_{d_h^{(i)}}[(\delta_h^{(t)})^2]}}\sqrt{\sum_{t=1}^T\sum_{h=1}^H \left(H^2 \vee \sum_{i=1}^{t-1}\E_{d_h^{(i)}}[(\delta_h^{(t)})^2]\right)} \\
    &\overset{\text{(a)}}{\lesssim} \sqrt{H\cdot\mathsf{SEC}(\gF, \Pi, T)}\sqrt{\beta H^3 T}\\
    &\lesssim \sqrt{\beta H^4 \mathsf{SEC}(\gF, \Pi, T)\cdot T}.
\end{align*}
    In (a), the first term is bounded by the definition of SEC and the second by Lemma \ref{lem:optimism-in-sample-error}.
\end{proof}



\subsection{Diffusion Policy Behavior Cloning Error}
\label{app:bc-sampler-constants}
The analysis of diffusion policy behavior cloning error mainly relies on the following lemma on sampling error of diffusion models. The greedy policy is deterministic, whereas the deployed sampler retains noise. We compare them through the projected noised target $\widetilde\pi_h^{f^{(t)}}$, which appears only in the analysis. Denote $W_1(\cdot, \cdot)$ as the Wasserstein-1 distance between two distributions, i.e.
\[
 W_1(P,Q)=\inf_{\Gamma\in\mathcal C(P,Q)}\int\|x-y\|_2\,d\Gamma(x,y),
\]
where $\mathcal C(P,Q)$ is the set of couplings of $P$ and $Q$. We first specify the sampler coefficients and constants used in the comparison.

For this cloning result, take $K\ge2$ and $0<\alpha_k<1$. In Equation~\ref{eq:DDPM} use
\[
 \eta_k^{\mathsf{ddpm}}=-\frac{1-\alpha_k}{\sqrt{v_k}},\qquad
 (\sigma_k^{\mathsf{ddpm}})^2=\frac{\alpha_k(1-\alpha_k)v_{k-1}}{v_k},\quad k=K,\ldots,2.
\]
These coefficients give the exact Gaussian reverse kernels of the conditional target. Stop at $Y_1$ and project to $\gA$. Define the smoothing, initialization, and schedule terms
\begin{align*}
 b_K&=(1-\sqrt{\alpha_1})R_A+\sqrt{d_A(1-\alpha_1)},\\
 I_K&=\tfrac12[\bar\alpha_KR_A^2+d_A(-\bar\alpha_K-\log(1-\bar\alpha_K))],\qquad
 \kappa_K=K\max_{2\le k\le K}\frac{(\eta_k^{\mathsf{ddpm}})^2}{2(\sigma_k^{\mathsf{ddpm}})^2}.
\end{align*}
For DDIM use $\eta_k^{\mathsf{ddim}}=\sqrt{\alpha_kv_{k-1}}-\sqrt{v_k}$ and the same stopping/projection rule. Assume its learned maps are $\ell_k$-Lipschitz, uniformly in states and iterates. Put
\[
 J_K=\sqrt{\bar\alpha_K}R_A+(1-\sqrt{v_K})\sqrt{d_A},\quad
 P_K=\prod_{k=2}^K\ell_k,\quad
 C_K=\left[K\sum_{k=2}^K\frac{(\eta_k^{\mathsf{ddim}})^2}{\alpha_k}
              \left(\prod_{j=2}^{k-1}\ell_j\right)^2\right]^{1/2}.
\]
Empty products equal one. These DDIM factors measure uniform stability of the learned maps. The $\widetilde O(\sqrt T)$ specialization in Section~\ref{sec:behavior-clone} assumes nonexpansiveness, $\ell_k\le1$. The Wasserstein argument below uses this stability and prediction error instead of a separate average Jacobian-matching error. An average Jacobian-error bound alone does not control $P_K$ or $C_K$. Lemma~\ref{lem:behavior-cloning-error} bounds the per-stage actor error by
\[
 B_t^{\mathsf{ddpm}}=L_f b_K+H\sqrt{[I_K+\kappa_K(\epsscore^{(t)})^2]/2},\qquad
 B_t^{\mathsf{ddim}}=L_f[b_K+P_KJ_K+C_K\epsscore^{(t)}].
\]
The two bounds explain how the sampler affects cloning accuracy. For DDPM, the matched Gaussian reverse kernels let us accumulate prediction error through the KL chain rule. For DDIM, we follow how each prediction error propagates through the remaining denoising maps, which introduces the factors $P_K$ and $C_K$. In both cases, the choice of $K$ and its noise schedule controls how closely the initial Gaussian matches the noised target and how closely the final output approximates the greedy action.

\begin{lem}[Diffusion Sampler Error]
Let $\widetilde\pi_h^{f^{(t)}}(\cdot\mid s)$ be the projection to $\gA$ of $q_1(\cdot\mid s)$, with the sampler choices and constants above. Then
\[
 W_1(\pi_h^{f^{(t)}}(\cdot\mid s),\widetilde\pi_h^{f^{(t)}}(\cdot\mid s))\le b_K.
\]
For DDPM, the state-averaged sampling error satisfies
\[
 \E_{d_{h,S}^{(t)}}d_{\rm TV}(\widetilde\pi_h^{f^{(t)}}(\cdot\mid s),\pi_h^{(t)}(\cdot\mid s))
 \le\sqrt{[I_K+\kappa_K(\epsscore^{(t)})^2]/2}.
\]
For DDIM, under the stated learned-map Lipschitz bounds,
\[
 \E_{d_{h,S}^{(t)}}W_1(\widetilde\pi_h^{f^{(t)}}(\cdot\mid s),\pi_h^{(t)}(\cdot\mid s))
 \le P_KJ_K+C_K\epsscore^{(t)}.
\]
Consequently, in either case,
\[
 \E_{d_h^{(t)}}\langle f_h^{(t)}(s,\cdot),\pi_h^{f^{(t)}}(\cdot\mid s)-\pi_h^{(t)}(\cdot\mid s)\rangle
 \le B_t^{\mathsf{sampler}}.
\]
Since this holds at every stage, $B_t\le B_t^{\mathsf{sampler}}$ for the cloning loss in Equation~\ref{eq:bc-greediness-loss}.
\label{lem:behavior-cloning-error}
\end{lem}
The proof of Lemma~\ref{lem:behavior-cloning-error} is provided below.
To translate the sampling bounds into behavior cloning error, we add and subtract the intermediate measure:
\begin{align*}
 \langle f,\pi^f-\pi\rangle
 &=\langle f,\pi^f-\widetilde\pi^f\rangle+\langle f,\widetilde\pi^f-\pi\rangle\\
 &\le L_fW_1(\pi^f,\widetilde\pi^f)+H d_{\rm TV}(\widetilde\pi^f,\pi).
\end{align*}
The last inequality uses Lipschitzness and $0\le f\le H$, with $d_{\rm TV}(P,Q)=\sup_B|P(B)-Q(B)|$. For DDIM we instead bound the second term by $L_fW_1(\widetilde\pi^f,\pi)$. Averaging over the deployed states and applying the lemma gives
\[
 2\sum_{t=1}^T\sum_{h=1}^H\E_{d_h^{(t)}}
 \langle f_h^{(t)},\pi_h^{f^{(t)}}-\pi_h^{(t)}\rangle
 \le 2H\sum_{t=1}^T B_t
 \le 2H\sum_{t=1}^T B_t^{\mathsf{sampler}}.
\]
Combining this bound with Lemmas~\ref{lem:regret-decomp-dipo} and~\ref{lem:critic-estimation-error} proves Theorem~\ref{thm:dipo-bc-regret}. The smoothing term quantifies the value lost by replacing the deterministic greedy action with its slightly noised version, with sensitivity controlled by $L_f$.
\begin{proof}[Proof of Lemma~\ref{lem:behavior-cloning-error}]
Fix $t,h,s$ and suppress these indices. Couple the deterministic greedy action $a^*=\pi_h^{f^{(t)}}(s)$ with $X_1=\sqrt{\alpha_1}a^*+\sqrt{1-\alpha_1}Z$. Projection fixes $a^*$ and is nonexpansive, so
\begin{align}
 W_1(\pi^f,\widetilde\pi^f)&\le\E\|a^*-\operatorname{Proj}_{\gA}(X_1)\|_2\nonumber\\
 &\le(1-\sqrt{\alpha_1})R_A+\sqrt{1-\alpha_1}\E\|Z\|_2\le b_K.\label{eq:w1}
\end{align}
Thus the smoothing error decreases as the final noise variance $1-\alpha_1$ decreases.

\paragraph{DDPM score and initialization errors.}
The conditional forward chain $X_0=a^*$, $X_k=\sqrt{\alpha_k}X_{k-1}+\sqrt{1-\alpha_k}Z_k$ is Gaussian. Its marginal is $q_k$, and Gaussian conditioning gives for $k\ge2$
\[
 X_{k-1}\mid X_k=x\sim\gN\left(\frac{x+\eta_k^{\mathsf{ddpm}}\epsilon_{\theta^*}(x,k,s)}{\sqrt{\alpha_k}},
                       \frac{(\sigma_k^{\mathsf{ddpm}})^2}{\alpha_k}I\right).
\]
Thus the coefficients specified above give the exact reverse kernels for this deterministic conditional target. Let $Q_s$ denote the reversed forward-path law on $X_{K:1}$ and $P_{\theta,s}$ the learned chain on $Y_{K:1}$ initialized from $N(0,I)$, with the same positive covariances. The KL chain rule gives
\begin{align}
 \KL(Q_s\Vert P_{\theta,s})
 &=\KL(q_K(\cdot\mid s)\Vert N(0,I))\nonumber\\
 &\quad+\sum_{k=2}^K\frac{(\eta_k^{\mathsf{ddpm}})^2}{2(\sigma_k^{\mathsf{ddpm}})^2}
 \E_{q_k(\cdot\mid s)}\|\epsilon_{\theta_h^{(t)}}-\epsilon_{\theta^*}\|_2^2.\label{eq:score-error-KL-TV}
\end{align}
The factors $1/\alpha_k$ in the mean difference and covariance cancel. The initial KL is exactly
\[
 \tfrac12[\bar\alpha_K\|a^*\|_2^2+d_A(-\bar\alpha_K-\log(1-\bar\alpha_K))]\le I_K.
\]
Data processing through the endpoint and projection, followed by Pinsker's inequality, bounds the conditional TV distance by $\sqrt{\KL(Q_s\Vert P_{\theta,s})/2}$. Average over $s\sim d_{h,S}^{(t)}$ and use Jensen's inequality and Assumption~\ref{aspt:score-matching-error}. The weighted error is at most $\kappa_K(\epsscore^{(t)})^2$, giving the DDPM claim. Averaging over the deployed states matches the measure appearing in the regret decomposition. Since the conditional target is Gaussian, the chosen reverse kernels are exact, and their error separates into initialization and prediction terms. Stopping at time one keeps the reverse covariances positive, with the remaining distance to the greedy action accounted for by the smoothing bound.

For the DDPM rate in Equation~\ref{eq:bc-sampler-rates}, take $q=T+1$, $K=\lceil\sqrt q\rceil+1$, and $r_k=v_k/(1-v_k)=q^{-1+2(k-1)/(K-1)}$. Direct substitution gives
\[
 \frac{(\eta_k^{\mathsf{ddpm}})^2}{2(\sigma_k^{\mathsf{ddpm}})^2}
 =\frac12\left(\frac{r_k}{r_{k-1}}-1\right),\qquad
 \kappa_K=\frac K2\left(q^{2/(K-1)}-1\right)\le8\log q.
\]
Here $2\log q/(K-1)\le4/e$ and $e^x-1\le e^{4/e}x$ on this interval. Also $v_1=\bar\alpha_K=1/(q+1)$, so
\[
 b_K\le\frac{R_A}{q+1}+\sqrt{\frac{d_A}{q+1}},\qquad
 I_K\le\frac{R_A^2}{2(q+1)}+\frac{3d_A}{8(q+1)^2}.
\]
Substituting these bounds into $2H\sum_t B_t^{\mathsf{ddpm}}$ gives the stated DDPM rate.

\paragraph{DDIM sampling error.}
Write $m_k=\sqrt{\bar\alpha_k}a^*$ and define the exact Gaussian transport
\[
 T_k(x)=m_{k-1}+\sqrt{v_{k-1}/v_k}(x-m_k).
\]
It maps $q_k$ to $q_{k-1}$ and equals the DDIM update with the stated coefficient and exact predictor. Couple $X_K=m_K+\sqrt{v_K}Z$ with $Y_K=Z$, so that $\E\|X_K-Y_K\|_2\le J_K$. Evolve $X_{k-1}=T_k(X_k)$ and $Y_{k-1}=\Phi_k^\theta(Y_k)$. The learned-map Lipschitz bound gives
\[
 \E\|Y_{k-1}-X_{k-1}\|_2\le\ell_k\E\|Y_k-X_k\|_2+
 \frac{|\eta_k^{\mathsf{ddim}}|}{\sqrt{\alpha_k}}
 \E_{q_k}\|\epsilon_{\theta_h^{(t)}}-\epsilon_{\theta^*}\|_2.
\]
Iterate from $K$ to $2$ and project both endpoints. Averaging states and applying Cauchy--Schwarz to the weighted sum over $k$ yields $P_KJ_K+C_K\epsscore^{(t)}$. Finally use the smoothing bound and the actor-error inequalities above. For this deterministic conditional target, the transport comparison controls cloning error directly through prediction accuracy and the stability of the learned maps, without requiring inverse maps or a separate Jacobian-matching error.

The DDIM bound also gives an explicit square-root sampler contribution under
a sufficient stability condition. Suppose that every learned denoising map is
nonexpansive, so that $\ell_k\le1$, and use the preceding geometric noise-odds
schedule. Put $q=T+1$, $n=K-1=\lceil\sqrt q\rceil$, and
$u=q^{1/n}$. Since $v_k/(1-v_k)=u^2v_{k-1}/(1-v_{k-1})$, we have
\[
 \frac{|\eta_k^{\mathsf{ddim}}|}{\sqrt{\alpha_k}}
 =\sqrt{v_k/\alpha_k}-\sqrt{v_{k-1}}
 =(u-1)\sqrt{v_{k-1}}\le u-1.
\]
Consequently,
\[
 P_K\le1,\qquad C_K\le\sqrt{K(K-1)}(u-1)
 \le\sqrt{3/2}\,e^{2/e}\log q.
\]
The last inequality uses $K/(K-1)\le3/2$ for $q\ge2$,
$\log q/n\le2/e$, and $e^x-1\le e^{2/e}x$ on this interval.
Moreover,
\[
 J_K\le \frac{R_A}{\sqrt{q+1}}+\frac{\sqrt{d_A}}{q+1},
 \qquad b_K\le\frac{R_A}{q+1}+\sqrt{\frac{d_A}{q+1}}.
\]
Substitution yields
\[
 2H\sum_{t=1}^T B_t^{\mathsf{ddim}}
 =O\!\left(\sqrt T+\log(T+1)\sum_t\epsscore^{(t)}\right)
\]
for fixed $H,L_f,R_A,d_A$. Nonexpansiveness is a genuine additional condition
on the learned maps. The exact conditional Gaussian maps satisfy it because
$\operatorname{Lip}(T_k)=\sqrt{v_{k-1}/v_k}<1$. A sufficient perturbation
condition is
\[
 \operatorname{Lip}_a(\epsilon_{\theta_h^{(t)}}-\epsilon_{\theta^*})
 \le\frac{\sqrt{\alpha_k}}{|\eta_k^{\mathsf{ddim}}|}
          \left(1-\sqrt{v_{k-1}/v_k}\right),\qquad k\ge2.
\]
This follows by writing the learned map as the exact Gaussian map plus
$\eta_k^{\mathsf{ddim}}(\epsilon_{\theta_h^{(t)}}-\epsilon_{\theta^*})/
\sqrt{\alpha_k}$ and adding their Lipschitz bounds. An average Jacobian error
alone does not imply this uniform condition.
\end{proof}

% We have a TV distance bound from Gen. Only randomness is expectation over score error. Assume action space is bounded and we don't need to make any assumptions over the state space.

\subsection{Optimism and In-Sample Error Control}

\begin{lem}[Optimism and In-sample Error Control, Analogue of Lemmas 39 and 40 of \citet{jin2021bellman}]
    With probability at least $1-\delta$, simultaneously for all $t\in[T]$ and $h\in[H]$, 
    \begin{enumerate}[label=(\roman*)]
        \item $Q^*\in\gF^{(t-1)}$,
        \item $\sum_{i=1}^{t-1}\E_{d_h^{(i)}}[(f_h-\gT_hf_{h+1})^2] \leq O(\beta H^2)$ for all $f\in\gF^{(t-1)}$.
    \end{enumerate}
    \label{lem:optimism-in-sample-error}
\end{lem}
\begin{proof}
The proof follows analogously to Lemmas 39 and 40 of \citet{jin2021bellman}. The states and actions are drawn from the deployed diffusion policies, so the concentration statement must be uniform over the critic and its target. For $g\in\gF_h$, $u\in\gF_{h+1}$, put $m=\gT_hu$, $Y_i=r_h^{(i)}+\max_a u(s_{h+1}^{(i)},a)$, and
\[
 X_i(g,u)=[g(s_h^{(i)},a_h^{(i)})-Y_i]^2-[m(s_h^{(i)},a_h^{(i)})-Y_i]^2.
\]
Given the history before episode $i$, its mean is $\mu_i=\E_{d_h^{(i)}}(g-m)^2$, $|X_i|\le4H^2$, and $\E[X_i^2\mid\mathcal H_{i-1}]\le16H^2\mu_i$. Here $g,m\in[0,H]$ by completeness and $Y_i\in[0,H+1]$. For $D_i=X_i-\mu_i$ and $\lambda=(32H^2)^{-1}$, $|\lambda D_i|\le1/4$ and
\[
 \E[e^{\pm\lambda D_i}\mid\mathcal H_{i-1}]\le\exp(16\lambda^2H^2\mu_i).
\]
Multiplying these bounds and applying Markov's inequality, then a union bound over times, stages, and a uniform-norm net of pairs $(g,u)$, gives
\[
 \left|\sum_{i<t}X_i(g,u)-\sum_{i<t}\mu_i\right|
 \le\tfrac12\sum_{i<t}\mu_i+C_0H^2\ell,
 \quad\ell=\log\!\left(\frac{2H(T+1)\gN_\gF(1/(T+1))^2}{\delta}\right).
\]
The bound is simultaneous for all pairs: extending from the net changes the accumulated losses and squared residuals by $O(TH/(T+1))$, absorbed into the constant. Choose $\beta\ge2C_0\ell$. For target $u=Q_{h+1}^*$, completeness gives $m=Q_h^*$, and the lower bound places $Q^*$ within the empirical confidence threshold. For a selected critic, completeness permits comparison against $m=\gT_hf_{h+1}^{(t)}$ in the empirical minimum, and the same lower bound gives the claimed in-sample error. Optimistic maximization over $\gF^{(t-1)}$ completes the argument.
\end{proof}

\section{Analysis of Algorithm \ref{alg:dipo-bc} with Hybrid Data}
\label{app:dipo-fqi}




The analysis framework we use is similar to the online-data analysis, where we decompose the regret into behavior cloning error and critic estimation error.
The behavior cloning error analysis is the same as before, so we omit it here and focus on critic estimation error.
\subsection{Regret Decomposition}

All initial-state expressions are averaged over $\rho$, as in the online analysis.

\begin{lem}[Hybrid Regret Decomposition for Algorithm \ref{alg:dipo-bc}]
    The following regret decomposition holds for Algorithm \ref{alg:dipo-bc} with hybrid data:
    \begin{align*}
        \text{Reg}(T) &= \sum_{t=1}^T\left(V_1^{*}(s_1^{(t)})-V_1^{\pi^{(t)}}(s_1^{(t)})\right) \\
        &\leq \underbrace{2\sum_{t=1}^T\sum_{h=1}^H \E_{d_h^{(t)}}\left[\left\langle f_h^{(t)}(s_h, \cdot), \pi_h^{f^{(t)}}(\cdot \mid s_h)-\pi_h^{(t)}(\cdot \mid s_h)\right\rangle\right]}_{\textcolor{red}{\text{Diffusion policy behavior cloning error.}}} \\
        &\quad+ \underbrace{\sum_{t=1}^T \sum_{h=1}^H \mathbb{E}_{\pi^*}\left[\mathcal{T}_h f_{h+1}^{(t)}- f_h^{(t)}\right] + \sum_{t=1}^T \sum_{h=1}^H \mathbb{E}_{d_h^{(t)}}\left[f_h^{(t)}-\gT_h f_{h+1}^{(t)}\right]}_{\textcolor{blue}{\text{Critic estimation error.}}} .
    \end{align*}
    \label{lem:regret-decomp-dipo-hybrid}
\end{lem}


\begin{proof}
As in the proof of Lemma \ref{lem:regret-decomp-dipo}, we first add and subtract $f_1^{(t)}(s_1^{(t)}, \pi_1^{f^{(t)}}(s_1^{(t)}))$, obtaining
    \begin{align}
        \text{Reg}(T) &= \sum_{t=1}^T\left(V_1^{*}(s_1^{(t)})-V_1^{\pi^{(t)}}(s_1^{(t)})\right) \\
        &= \sum_{t=1}^T \left(V_1^{*}(s_1^{(t)})-f_1^{(t)}(s_1^{(t)}, \pi_1^{f^{(t)}}(s_1^{(t)})\right) + \sum_{t=1}^T\left(f_1^{(t)}(s_1^{(t)}, \pi_1^{f^{(t)}}(s_1^{(t)})-V_1^{\pi^{(t)}}(s_1^{(t)})\right).
        \label{eq:regret-decomp-alg-4-1}
    \end{align}

    Unlike the proof of Lemma \ref{lem:regret-decomp-dipo}, the first term in Equation \ref{eq:regret-decomp-alg-4-1} as-is is not necessarily non-positive as we do not have optimism. Still, we can address this by applying the generalized policy difference lemma from \citep{efroni2020optimisticpolicyoptimizationbandit, cai2024provablyefficientexplorationpolicy} (presented in Lemma \ref{ch05:lem:policy-value-difference-lemma} for completeness) with $\pi = \pi^{f^{(t)}}, \pi' = \pi^*$, and with $f^{(t)}$ as the Q-function, yielding the decomposition:
    \begin{align}
        &\sum_{t=1}^T \left(V_1^{*}(s_1^{(t)})-f_1^{(t)}(s_1^{(t)}, \pi_1^{f^{(t)}}(s_1^{(t)})\right)\\
        &= \sum_{t=1}^T\sum_{h=1}^H \E_{\pi^*}\left[\left\langle f_h^{(t)}\left(s_h, \cdot\right), \pi_h^{*}\left(\cdot \mid s_h\right)-\pi_h^{f^{(t)}}\left(\cdot \mid s_h\right)\right\rangle\right] + \sum_{t=1}^T \sum_{h=1}^H \mathbb{E}_{\pi^*}\left[\gT_hf_{h+1}^{(t)} - f_h^{(t)}\right].
        \label{eq:first-term-decomp-alg-4-2}
    \end{align}
    Here $\gT_h^{\pi^{f^{(t)}}}f_{h+1}^{(t)}=\gT_hf_{h+1}^{(t)}$ because $\pi^{f^{(t)}}$ is greedy. However, by definition,
\[f_h^{(t)}(s, \pi_h^*(s)) \leq \max_a f_h^{(t)}(s, a) = f_h^{(t)}(s, \pi_h^{f^{(t)}}(s)),\]
so we have that the first term in Equation \ref{eq:first-term-decomp-alg-4-2} is non-positive:
    \begin{equation}
        \sum_{t=1}^T\sum_{h=1}^H \E_{\pi^*}\left[\left\langle f_h^{(t)}\left(s_h, \cdot\right), \pi_h^{*}\left(\cdot \mid s_h\right)-\pi_h^{f^{(t)}}\left(\cdot \mid s_h\right)\right\rangle\right] \leq 0.
        \label{eq:nonpositive-alg-4-3}
    \end{equation}
    
    We therefore have that
    \begin{equation}
        \sum_{t=1}^T \left(V_1^{*}(s_1^{(t)})-f_1^{(t)}(s_1^{(t)}, \pi_1^{f^{(t)}}(s_1^{(t)})\right) \leq \sum_{t=1}^T \sum_{h=1}^H \mathbb{E}_{\pi^*}\left[\mathcal{T}_h f_{h+1}^{(t)}- f_h^{(t)}\right]
        \label{eq:first-term-final-regret-decomp-alg-4-6}
    \end{equation}

    
    It remains to address the original second term in Equation \ref{eq:regret-decomp-alg-4-1}. Repeating the same arguments in the proof of Lemma \ref{lem:regret-decomp-dipo} yields:
    \begin{align}
        &\sum_{t=1}^T\left(f_1^{(t)}(s_1^{(t)}, \pi_1^{f^{(t)}}(s_1^{(t)})-V_1^{\pi^{(t)}}(s_1^{(t)})\right)  \\
        &\leq \underbrace{2\sum_{t=1}^T\sum_{h=1}^H \E_{d_h^{(t)}}\left[\left\langle f_h^{(t)}(s_h, \cdot), \pi_h^{f^{(t)}}(\cdot \mid s_h)-\pi_h^{(t)}(\cdot \mid s_h)\right\rangle\right]}_{\textcolor{red}{\text{Diffusion policy behavior cloning error.}}} + \underbrace{\sum_{t=1}^T \sum_{h=1}^H \mathbb{E}_{d_h^{(t)}}\left[f_h^{(t)}-\gT_h f_{h+1}^{(t)}\right]}_{\textcolor{blue}{\text{Critic estimation error.}}}.
        \label{eq:second-term-final-regret-decomp-alg-4-7}
    \end{align}
    Putting Equations \ref{eq:first-term-final-regret-decomp-alg-4-6} and \ref{eq:second-term-final-regret-decomp-alg-4-7} together yields the final regret decomposition:
    \begin{align}
        \text{Reg}(T) &= \sum_{t=1}^T\left(V_1^{*}(s_1^{(t)})-V_1^{\pi^{(t)}}(s_1^{(t)})\right) \\
        &\leq \underbrace{2\sum_{t=1}^T\sum_{h=1}^H \E_{d_h^{(t)}}\left[\left\langle f_h^{(t)}(s_h, \cdot), \pi_h^{f^{(t)}}(\cdot \mid s_h)-\pi_h^{(t)}(\cdot \mid s_h)\right\rangle\right]}_{\textcolor{red}{\text{Diffusion policy behavior cloning error.}}} \\
        &\quad+ \underbrace{\sum_{t=1}^T \sum_{h=1}^H \mathbb{E}_{\pi^*}\left[\mathcal{T}_h f_{h+1}^{(t)}- f_h^{(t)}\right] + \sum_{t=1}^T \sum_{h=1}^H \mathbb{E}_{d_h^{(t)}}\left[f_h^{(t)}-\gT_h f_{h+1}^{(t)}\right]}_{\textcolor{blue}{\text{Critic estimation error.}}} .
    \end{align}
\end{proof}

\subsection{Controlling Error Under $\pi^*$}

We control the diffusion policy behavior cloning error (Lemma \ref{lem:behavior-cloning-error}) and the critic estimation error under $d_h^{(t)}$ (Lemma \ref{lem:critic-estimation-error}) in the exact same way as we do for Algorithm \ref{alg:dipo-bc}. The only difference is that we use Lemma \ref{lem:in-sample-error-control-hybrid-dipo} within the proof of Lemma \ref{lem:critic-estimation-error} instead of Lemma \ref{lem:optimism-in-sample-error}, but otherwise the arguments are verbatim the same. We therefore omit the proofs for Lemmas \ref{lem:behavior-cloning-error} and Lemma \ref{lem:critic-estimation-error} in the case of Algorithm \ref{alg:dipo-bc} with hybrid data to avoid redundancy.

It remains to control the error under $\pi^*$. To do so, we use the offline data. The argument is verbatim the same as the proof of Theorem 3 in \citet{tan2025actorcriticsachieveoptimalsample}.

Concretely, using the squared signed Bellman-error expectation in Definition~\ref{defn:spc},
\begin{align*}
    &\sum_{t=1}^T \sum_{h=1}^H\mathbb{E}_{\pi^*}\left[\left(\gT_hf_{h+1}^{(t)} - f_h^{(t)}\right)\left(s_h, a_h\right)\right] \\ 
    &\leq \sum_{t=1}^T \sum_{h=1}^H \E_{\pi^*}\left[\left(\gT_hf_{h+1}^{(t)} - f_h^{(t)}\right)\right] \left(\frac{N_{\off}\E_{\mu_h}\left[\left(\gT_hf_{h+1}^{(t)} - f_h^{(t)}\right)^2\right] + \sum_{i=1}^{t-1}\E_{d_h^{(i)}}\left[\left(\gT_hf_{h+1}^{(t)} - f_h^{(t)}\right)^2\right]}{N_{\off}\E_{\mu_h}\left[\left(\gT_hf_{h+1}^{(t)} - f_h^{(t)}\right)^2\right] + \sum_{i=1}^{t-1}\E_{d_h^{(i)}}\left[\left(\gT_hf_{h+1}^{(t)} - f_h^{(t)}\right)^2\right]}\right)^{1/2} \\
    &\leq \sqrt{\sum_{t=1}^T\sum_{h=1}^H \frac{\E_{\pi^*}\left[\gT_hf_{h+1}^{(t)} - f_h^{(t)}\right]^2}{N_{\off}\E_{\mu_h}\left[\left(\gT_hf_{h+1}^{(t)} - f_h^{(t)}\right)^2\right] + \sum_{i=1}^{t-1}\E_{d_h^{(i)}}\left[\left(\gT_hf_{h+1}^{(t)} - f_h^{(t)}\right)^2\right]}}\\
    &\qquad \cdot\sqrt{\sum_{t=1}^T\sum_{h=1}^H \left(N_{\off}\E_{\mu_h}\left[\left(\gT_hf_{h+1}^{(t)} - f_h^{(t)}\right)^2\right] + \sum_{i=1}^{t-1}\E_{d_h^{(i)}}\left[\left(\gT_hf_{h+1}^{(t)} - f_h^{(t)}\right)^2\right]\right)} \\
    &\leq \sqrt{\sum_{t=1}^T\sum_{h=1}^H \frac{\E_{\pi^*}\left[\gT_hf_{h+1}^{(t)} - f_h^{(t)}\right]^2}{N_{\off}\E_{\mu_h}\left[\left(\gT_hf_{h+1}^{(t)} - f_h^{(t)}\right)^2\right]}}\\
    &\qquad \cdot\sqrt{\sum_{t=1}^T\sum_{h=1}^H \left(N_{\off}\E_{\mu_h}\left[\left(\gT_hf_{h+1}^{(t)} - f_h^{(t)}\right)^2\right] + \sum_{i=1}^{t-1}\E_{d_h^{(i)}}\left[\left(\gT_hf_{h+1}^{(t)} - f_h^{(t)}\right)^2\right]\right)} \\
    &\lesssim \sqrt{c_{\off}^*(\gF)HT/N_{\off}}\sqrt{\beta H^3T} \\
    &\lesssim \sqrt{H^4\beta c_{\off}^*(\gF) T^2/N_{\off}}.
\end{align*}
We use the single-policy concentrability coefficient to bound the squared residual ratio $$\frac{\E_{\pi^*}\left[\gT_hf_{h+1}^{(t)} - f_h^{(t)}\right]^2}{\E_{\mu_h}\left[\left(\gT_hf_{h+1}^{(t)} - f_h^{(t)}\right)^2\right]} \leq \max_{h\in[H]}\sup_{f \in \gF} \frac{\E_{\pi^*}\left[f_h-\gT_hf_{h+1} \right]^2}{\E_{\mu_h}[\left(f_h-\gT_h f_{h+1}\right)^2]} = c_{\off}^*(\gF),$$
The sign change disappears after squaring. Terms with zero denominator are interpreted using the convention in Definition~\ref{defn:spc}.



\subsection{In-Sample Error Control}

\begin{lem}[In-Sample Error Control for Algorithm \ref{alg:dipo-bc} with Hybrid Data]
With probability at least $1-\delta$, simultaneously for all $t\in[T]$ and $h\in[H]$, Algorithm~\ref{alg:dipo-bc} with hybrid data enjoys the guarantee that
    $$N_{\off} \E_{\mu_h}\left[\left(f_h^{(t)} - \gT_h f_{h+1}^{(t)}\right)^2\right] + \sum_{i=1}^{t-1} \E_{d_h^{(i)}}\left[\left(f_h^{(t)} - \gT_h f_{h+1}^{(t)}\right)^2\right] \leq O(H^2\beta),$$
    for the $\beta$ and offline sampling model in the hybrid regret theorem.
    \label{lem:in-sample-error-control-hybrid-dipo}
\end{lem}
\begin{proof}
The proof is analogous to that of Lemma 39 in \citet{jin2021bellman}. We append the independent offline episodes to the start of the sequence of online episodes, and let $\mathcal H_{i-1}$ be the history before episode $i$ in this combined sequence. For a fixed stage $h$ and pair $g\in\gF_h$, $u\in\gF_{h+1}$, let $m=\gT_hu$ and $Y_i=r_h^{(i)}+\max_a u(s_{h+1}^{(i)},a)$. As in Lemma~\ref{lem:optimism-in-sample-error}, define
\begin{align}
 X_i(g,u)
 &=[g(s_h^{(i)},a_h^{(i)})-Y_i]^2
   -[m(s_h^{(i)},a_h^{(i)})-Y_i]^2.
 \label{eq:err-control-alg-4-1}
\end{align}
Now, this has conditional expectation
\[
 \E[X_i(g,u)\mid\mathcal H_{i-1}]
 =\E[(g-m)^2(s_h^{(i)},a_h^{(i)})\mid\mathcal H_{i-1}],
\]
and variance bounded by
\begin{align*}
 \operatorname{Var}(X_i(g,u)\mid\mathcal H_{i-1})
 &\leq\E[X_i(g,u)^2\mid\mathcal H_{i-1}]\\
 &\leq16H^2\E[X_i(g,u)\mid\mathcal H_{i-1}].
\end{align*}
For a fixed pair, offline conditional expectations are under $\mu_h$, while the $j$-th online episode has conditional state--action law $d_h^{(j)}$. Applying the concentration argument in Lemma~\ref{lem:optimism-in-sample-error} and a union bound over times, stages, and a net of pairs with radius $(N_{\off}+T+1)^{-1}$ gives
\[
 \left|\sum X_i(g,u)-B(g,u)\right|\le\tfrac12 B(g,u)+C_0H^2\ell,
\]
where $B(g,u)=N_{\off}\E_{\mu_h}(g-\gT_hu)^2+\sum_{i<t}\E_{d_h^{(i)}}(g-\gT_hu)^2$ and $\ell$ is the logarithmic factor in the hybrid theorem. Here the sum of $X_i$ includes all offline episodes and the first $t-1$ online episodes. The uniform concentration event is simultaneous in both $g$ and $u$, so it applies to targets fitted backwards on the same data.

This is where the proof differs. As we choose the empirical risk minimizer, and completeness places $\gT_hf_{h+1}^{(t)}$ in $\gF_h$,
\[
 \sum X_i(f_h^{(t)},f_{h+1}^{(t)})
 =\widehat{\gL}_h^{(t-1)}(f_h^{(t)},f_{h+1}^{(t)})
  -\widehat{\gL}_h^{(t-1)}(\gT_hf_{h+1}^{(t)},f_{h+1}^{(t)})\le0.
\]
The lower concentration bound therefore gives $B(f_h^{(t)},f_{h+1}^{(t)})\le2C_0H^2\ell\le O(H^2\beta)$. Reusing the offline data over iterations is permitted because this event is uniform over all pairs.
\end{proof}

\subsection{Cloning a critic-tilted pretrained policy}
\label{app:pretrained-gibbs-bc}
Cloning one greedy action can discard other good actions already represented by a pretrained diffusion policy. We can instead reweight that policy toward actions favored by the current critic. This preserves its stochastic support while allowing us to quantify the loss relative to a greedy action. Let $\pi_{0,h}(\cdot\mid s)$ be the law of the pretrained diffusion sampler, supported on $\gA$, and define the cloning target
\[
 q_{\tau,h}^f(da\mid s)
 =\frac{e^{f_h(s,a)/\tau}\pi_{0,h}(da\mid s)}{Z_{\tau,h}^f(s)},
 \qquad Z_{\tau,h}^f(s)=\int e^{f_h(s,a)/\tau}\pi_{0,h}(da\mid s),\quad \tau>0.
\]
The critic updates and transition data are the same as in the optimistic or backward-FQI analysis. We change the actor target, and account for the resulting greedy-action gap in the regret decomposition. The temperature $\tau$ determines how strongly the critic reweights the prior. The next proposition bounds this gap using the prior mass near a greedy action and gives a rejection sampler. Appendix~\ref{app:gibbs-neural-fitting} then fits a denoising actor to this generally multimodal target, whose reverse kernels differ from those for the deterministic target.

\ReserveSpace{15\baselineskip}
\begin{prop}[Pretrained Gibbs targets and a sampling realization]
\label{prop:gibbs-cloning}
Retain the critic, action, and sampling assumptions of
Theorem~\ref{thm:dipo-bc-regret}, except for its actor score-error and sampler
conditions. For every selected critic, stage, and deployed state, suppose that,
at some radius $0<r\le R$,
\[
 \pi_{0,h}(B(a_h^f(s),r)\mid s)\ge c_0(r/R)^m,
 \qquad a_h^f(s)\in\argmax_a f_h(s,a),
\]
where $0<c_0\le1$, $m>0$, and the ball is intersected with $\gA$.
Write $D_r=\log(c_0^{-1})+m\log(R/r)$ and $g_{\tau,r}=L_fr+\tau D_r$.
If the deployed actor satisfies
\[
 \E_{s\sim d_{h,S}^{(t)}}d_{\rm TV}
   (q_{\tau,h}^{f^{(t)}}(\cdot\mid s),\pi_h^{(t)}(\cdot\mid s))
 \le e_{t,h},
\]
then, with the same high-probability critic event,
\[
 \mathrm{Reg}(T)\le \mathcal C_T+HTg_{\tau,r}
                   +H\sum_{t=1}^T\sum_{h=1}^H e_{t,h}.
\]
In online RL the critic contribution is
\[\mathcal C_T=O(H^2\sqrt{\beta T\mathsf{SEC}(\gF,\Pi,T)}).\]
In hybrid RL it additionally contains
\[O(H^2T\sqrt{\beta_{\rm hyb}c_{\off}^*(\gF)/N_{\off}}),\]
with the confidence parameters defined in the respective critic theorems. In both
cases $\Pi$ contains the policies actually deployed in this extension.
Alternatively, a state-averaged Wasserstein error $w_{t,h}$ replaces the last
term by $L_f\sum_{t,h}w_{t,h}$.

There is a concrete implementation of this bound using the pretrained sampler
and the same greedy-action oracle as Algorithm~\ref{alg:dipo-bc}. At state $s$,
compute $M=f_h^{(t)}(s,a_h^{f^{(t)}}(s))$. Repeatedly propose
$A\sim\pi_{0,h}(\cdot\mid s)$ independently and accept with probability
$\exp((f_h^{(t)}(s,A)-M)/\tau)$. Unlimited rejection sampling returns exactly
$q_{\tau,h}^{f^{(t)}}$ and uses at most $p_0^{-1}$ proposals in expectation,
where
\[
 p_0=c_0(r/R)^m e^{-L_fr/\tau}.
\]
For $0<\varepsilon<1$, if we stop after $J\ge p_0^{-1}\log(1/\varepsilon)$ proposals and return the last
proposal when none is accepted, then $e_{t,h}\le\varepsilon$ uniformly. This uses at
most $J$ pretrained-policy samples and critic evaluations per environment
action, besides the greedy-action maximization. If each pretrained sample
uses $K_0$ denoising evaluations, its denoising cost is at most $JK_0$.

For $\tau=T^{-1/2}$ and $r=\min\{R,m\tau/L_f\}$ when $L_f>0$,
$HTg_{\tau,r}=\widetilde O(Hm\sqrt T+H\sqrt T\log(c_0^{-1}))$ for fixed
$L_f,R,m$. Taking $\varepsilon=(T+1)^{-1}$ contributes at most $H^2$ to regret.
When $m\tau\le L_fR$, the expected proposal bound and the capped proposal
bound are respectively
\[
 p_0^{-1}=c_0^{-1}\left(\frac{eL_fR}{m\tau}\right)^m,
 \qquad J=\left\lceil p_0^{-1}\log(T+1)\right\rceil.
\]
Thus the target-bias term has intrinsic small-ball dependence,
with a sampling cost polynomial in $1/\tau$ for fixed $m$, but exponential
in $m$. If $L_f=0$, the target is already $\pi_{0,h}$, and no rejection step is
needed.
\end{prop}

\begin{proof}
Fix $t,h,s$ and suppress these indices, writing $\pi_0$ for the pretrained law. Restrict $\pi_0$ to
$E=B(a^f(s),r)$ and normalize, obtaining $\nu_r$. Then
$\KL(\nu_r\Vert\pi_0)=\log(1/\pi_0(E))\le D_r$, and
$\E_{\nu_r}f\ge M-L_fr$. The Gibbs identity, valid for every
$\nu\ll\pi_0$ of finite relative entropy, is
\[
 \tau\log Z=\E_\nu f-\tau\KL(\nu\Vert\pi_0)
                          +\tau\KL(\nu\Vert q_\tau^f).
\]
In particular, the variational inequality and $\KL(q_\tau^f\Vert\pi_0)\ge0$
give
\[
 \E_{q_\tau^f}f\ge\tau\log Z
       \ge\E_{\nu_r}f-\tau\KL(\nu_r\Vert\pi_0)
       \ge M-L_fr-\tau D_r.
\]
Since $0\le f\le H$, the deployed greedy gap therefore satisfies
\[
 M-\E_\pi f\le g_{\tau,r}+H d_{\rm TV}(q_\tau^f,\pi).
\]
Lipschitzness gives the stated Wasserstein alternative.

For completeness, let $v_h^f(s)=\max_a f_h(s,a)$ and $v_{H+1}^f=0$.
A direct telescope along any deployed policy yields the exact identity
\[
 \E_\rho v_1^f-V^\pi
 =\sum_{h=1}^H\E_{d_h^\pi}
   [v_h^f(s_h)-f_h(s_h,a_h)+f_h(s_h,a_h)-\gT_hf_{h+1}(s_h,a_h)].
\]
Optimism makes $V^*-\E_\rho v_1^{f^{(t)}}\le0$. For backward FQI, the
same telescope along $\pi^*$ and greedy maximization instead give
\[
 V^*-\E_\rho v_1^{f^{(t)}}
 \le\sum_h\E_{d_h^{\pi^*}}[\gT_hf_{h+1}^{(t)}-f_h^{(t)}].
\]
The existing uniform critic concentration and SEC arguments apply to either
actor construction, and the offline Bellman-error coverage argument controls
the latter comparator term. Summing the greedy-gap bounds proves the regret
claim. The direct telescope counts each greedy gap once.

The probability of accepting a proposal is
\[
 p=e^{-M/\tau}Z=\E_{\pi_0} e^{(f-M)/\tau}
   \ge\pi_0(E)e^{-L_fr/\tau}\ge p_0.
\]
The accepted law is $q_\tau^f$ by the definition of rejection sampling.
Conditional on acceptance within $J$ proposals, it remains exactly
$q_\tau^f$: for each possible first acceptance index the accepted proposal
has this same law. The probability of no acceptance is
$(1-p)^J\le e^{-Jp_0}\le\varepsilon$. The capped law is consequently a mixture of
$q_\tau^f$ and a fallback law with mixing weight at most $\varepsilon$, proving the TV
bound. The sample and cost bounds follow by the geometric waiting-time
formula and substitution of the specified radius and temperature.
\end{proof}

The small-ball condition is needed uniformly near the greedy actions of the
selected critics, at the radius used in the bound. It is stronger than
coverage near only the optimal action. The exponent $m$ need not describe a
manifold, and for a fixed smooth full-dimensional density it cannot remain
below $d_A$ at arbitrarily small radii with fixed $c_0$. Exponential tilting preserves the prior's null sets and assigns positive mass to each prior mode of positive mass.

\paragraph{Probability of favorable actions.}\label{app:favorable-actions}
The stochastic target can preserve several promising action regions while increasing the probability of regions preferred by the critic. To quantify this at a fixed state, consider a measurable region $B$ with $p=\pi_0(B)$. If
$\inf_B f-\sup_{B^c}f\ge\Delta\ge0$, then
\[
 q_\tau(B)\ge\frac{p}{p+(1-p)e^{-\Delta/\tau}}.
\]
For $0<p<1$ and $\Delta>0$, this exceeds $p$.
Without separation, a region satisfying $\inf_Bf\ge M-\xi$ obeys
$q_\tau(B)\ge e^{-\xi/\tau}\pi_0(B)$.
If $\|f-Q\|_\infty\le\zeta$, true-value separation $\Delta_Q$ can be
used with $\Delta=\Delta_Q-2\zeta\ge0$.
An actor within TV distance $e$ satisfies the same lower bounds minus $e$.
For $N$ independent draws at this fixed state, the chance to discover $B$
is at least $1-(1-\max\{0,q_{\rm lower}-e\})^N$.
Indeed, with $Z_B=\int_B e^{f/\tau}\,d\pi_0$ and
$Z_{B^c}=\int_{B^c}e^{f/\tau}\,d\pi_0$, the separated-region bound follows from
$Z_{B^c}/Z_B\le(1-p)e^{-\Delta/\tau}/p$. For the near-maximizing
region use $Z_B\ge p e^{(M-\xi)/\tau}$ and $Z\le e^{M/\tau}$. They compare a specified favorable region with the same prior;
they do not imply smaller SEC or better environment-state coverage.

\subsection{Learning the critic-tilted target by denoising regression}
\label{app:gibbs-neural-fitting}

Rejection sampling implements the critic-tilted target directly, but its proposal cost grows as the target concentrates. We can instead fit a denoising actor to samples from the pretrained policy, using importance weights for the selected critic. The following regression and reverse-sampler bounds convert its fitted prediction error into the Wasserstein error charged by Proposition~\ref{prop:gibbs-cloning}. We give this construction for a known finite state space, with fresh synthetic samples at every state and finite-class optimization. A finite neural approximation family makes the procedure explicit, at the substantial computational cost stated below. This supplies one way to control the actor contribution while leaving practical replay-based SGD and the critic complexity conditions separate.

Fix a selected critic and stage, and write $R=R_A>0$, $d=d_A$.
The action set is compact and convex with $\|a\|\le R$.
At state $s$, set
\[
 M_s=\max_a f(s,a),\quad w_s(a)=e^{(f(s,a)-M_s)/\tau},\quad
 p_s=\mathbb E_{\pi_0(\cdot\mid s)}w_s(A),\quad
 q_s(da)=p_s^{-1}w_s(a)\pi_0(da\mid s).
\]
The maximum $M_s$ is evaluated with the same exact greedy-action oracle as
Proposition~\ref{prop:gibbs-cloning}, and the pretrained sampler can be
called at every known state. Suppose $p_s\ge p_0>0$. The small-ball hypothesis in Proposition~\ref{prop:gibbs-cloning}
supplies $p_0=c_0(r/R_0)^m e^{-L_f r/\tau}$, where $R_0$ denotes its reference
radius, distinct from the action bound $R$.
Let $K\ge2$, $v_k=1-\bar\alpha_k$, and $0<v_1<\cdots<v_K<1$.
The forward coefficients are $\alpha_1=1-v_1$ and
$\alpha_k=(1-v_k)/(1-v_{k-1})$ for $k\ge2$. Put
\[
 \lambda_k=\frac12\left(\frac1{v_{k-1}}-\frac1{v_k}\right),\quad
 \Lambda=\sum_{k=2}^K\lambda_k,\quad \rho_k=\lambda_k/\Lambda.
\]
For $A\sim q_s$, $Z\sim N(0,I_d)$ independent and
$X_k=\sqrt{1-v_k}A+\sqrt{v_k}Z$, let
$m_{s,k}(x)=\mathbb E[A\mid X_k=x]$.

\paragraph{Finite-data fitting.}
Let $\mathcal G$ be a finite family of neural clean-action predictors
$g(x,k)$ with values in the radius-$R$ ball. Their associated noise predictors
are
\[
 \epsilon_g(x,k)=\frac{x-\sqrt{1-v_k}g(x,k)}{\sqrt{v_k}}.
\]
All network weights may range over a finite grid; the theorem does not
restrict training to the output layer. For each state draw $n$ independent
triples
\[
 A_i\sim\pi_0(\cdot\mid s),\quad k_i\sim\rho,\quad Z_i\sim N(0,I_d),
 \qquad X_i=\sqrt{1-v_{k_i}}A_i+\sqrt{v_{k_i}}Z_i.
\]
Fit $\widehat g_s$ to within additive $\eta\ge0$ of the minimum of
\[
 \widehat L_s(g)=\frac1n\sum_{i=1}^n
       w_s(A_i)\|g(X_i,k_i)-A_i\|^2.
\]
This is denoising score matching with clean-action prediction: indeed
$\|g-A\|^2=v_k(1-v_k)^{-1}\|\epsilon_g-Z\|^2$ on each training sample.
Thus noise prediction uses the factor $v_k/(1-v_k)$ together with the time weights $\rho_k$ and critic weights $w_s(A_i)$. Enumeration of $\mathcal G$ is an executable optimizer with
$\eta=0$, costing $n|\mathcal G|$ network-loss evaluations per fitted state.
Define the approximation error, under the actual Gibbs target, by
\[
 a_s=\min_{g\in\mathcal G}\sum_{k=2}^K\rho_k
       \mathbb E_{X_k\sim q_{s,k}}\|g(X_k,k)-m_{s,k}(X_k)\|^2.
\]
\begin{prop}[Weighted neural denoising regression]
\label{prop:gibbs-neural-fit}
For $S=|\gS|$ states, $T$ rounds and $H$ stages, take
$L=\log(2|\mathcal G|SHT/\delta)$. The finite class and schedule may
change between fits, provided they are fixed before the fresh actor samples
are drawn, with $|\mathcal G|$ denoting the size for that fit. Conditional on each history before drawing
its fresh actor data, and hence jointly with probability at least $1-\delta$,
\begin{equation}\label{eq:gibbs-neural-fit}
 \sum_k\rho_k\mathbb E_{q_{s,k}}
       \|\widehat g_s-m_{s,k}\|^2
 \le 3a_s+\frac{96R^2L/n+2\eta}{p_s}
 \le 3a_s+\frac{96R^2L/n+2\eta}{p_0}
.
\end{equation}
The bound holds for every fitted state, stage, and round.
No deployed-state coverage coefficient is needed when every state is fitted.
All states must be available as inputs to the pretrained sampler and critic.
This requires no environment reset.
\end{prop}

\begin{proof}
Write $m=m_{s,k}(X)$ and
$D_g=w_s(A)(\|g-A\|^2-\|m-A\|^2)$.
Weighted conditional expectation makes $m$ the regression function, so
$r_g=\mathbb E D_g=p_s\mathbb E_{q_s,\rho}\|g-m\|^2\ge0$.
Since $g,m,A$ are in the radius-$R$ ball,
$|D_g|\le4R^2$ and
\[
 \mathbb E D_g^2\le16R^2\mathbb E[w_s(A)\|g-m\|^2]
                  =16R^2r_g.
\]
Bernstein's inequality and a union bound give, simultaneously for the finite
class,
\[
 |\widehat r_g-r_g|
 \le \sqrt{32R^2r_gL/n}+\frac{16R^2L}{3n}
 \le \frac{r_g}{2}+\frac{22R^2L}{n}.
\]
Here we used $|D_g-\mathbb E D_g|\le8R^2$ and the conservative
Bernstein remainder $2(8R^2)L/(3n)$.
Empirical minimization therefore gives
$r_{\widehat g}\le3\min_g r_g+88R^2L/n+2\eta$.
Enlarging $88$ to $96$ proves the claim.
Fresh conditional data and a union bound over $SHT$ fits handle changing
critics without a uniform complexity bound over the critic sequence.
\end{proof}

\begin{lem}[Reverse-sampler error for a bounded target]
\label{lem:gibbs-neural-sampler}
Run the standard Gaussian DDPM reverse kernels for $k=K,\ldots,2$, initialized
at $N(0,I_d)$, with the learned $\epsilon_{\widehat g_s}$, and project the
result at time one onto $\mathcal A$. The reverse covariance is
$\sigma_k^2 I_d$ with
$\sigma_k^2=(1-\alpha_k)v_{k-1}/v_k$; this is the effective covariance
\emph{after} the outer $1/\sqrt{\alpha_k}$ in the manuscript's sampler.
Write $\widehat\pi_s$ for its endpoint law and define
\[
 z_k=2R^2\lambda_k,\qquad
 B_K=\frac12\sum_{k=2}^K e^{4z_k}z_k^2,
\]
\[
 I_K=\frac12[(1-v_K)R^2+d(v_K-1-\log v_K)],\qquad
 b_K=(1-\sqrt{1-v_1})R+\sqrt{dv_1}.
\]
Then, on the fitting event,
\begin{equation}\label{eq:gibbs-neural-W1}
 W_1(q_s,\widehat\pi_s)
 \le b_K+2R\sqrt{\frac12\left[I_K+B_K+
    \Lambda\left(3a_s+\frac{96R^2L/n+2\eta}{p_0}\right)\right]}.
\end{equation}
This bound allows singular and multimodal Gibbs targets.
\end{lem}

\begin{proof}
The exact reverse kernel conditional on $X_k=x$ is the posterior mixture
\[
 Q_k(\cdot\mid x)=\mathbb E_{A\mid X_k=x}
  N(a_kx+c_kA,\sigma_k^2I_d),\quad
 a_k=\frac{\sqrt{\alpha_k}v_{k-1}}{v_k},\quad
 c_k=\frac{\sqrt{\bar\alpha_{k-1}}(1-\alpha_k)}{v_k}.
\]
Here $c_k^2/(2\sigma_k^2)=\lambda_k$.
Let $G_m=N(a_kx+c_km_{s,k}(x),\sigma_k^2 I_d)$.
For independent posterior copies of
$\Delta=c_k(A-m_{s,k}(x))$, direct integration of Gaussian densities gives
\[
 \chi^2(Q_k\Vert G_m)=
    \mathbb E_{\Delta,\Delta'}e^{\langle\Delta,\Delta'\rangle/\sigma_k^2}-1.
\]
The exponent has mean zero, absolute value at most $4z_k$, and second
moment at most $z_k^2$: the posterior variance
$\mathbb E\|A-m\|^2$ is at most $R^2$.
The inequality $e^u-1-u\le e^{|u|}u^2/2$ therefore gives
$\mathrm{KL}(Q_k\Vert G_m)\le\frac12e^{4z_k}z_k^2$.
Matching the mean in the Gaussian cross entropy gives the exact identity
\[
 \mathrm{KL}(Q_k\Vert G_{\widehat g_s})
 =\mathrm{KL}(Q_k\Vert G_m)+
                 \lambda_k\|m_{s,k}(x)-\widehat g_s(x,k)\|^2.
\]
Convexity bounds the initial divergence from $q_{s,K}$ to $N(0,I_d)$ by
$I_K$. The chain rule on reverse paths, followed by projection and
Pinsker's inequality, bounds the TV distance between the projected time-one
laws by the square root in \eqref{eq:gibbs-neural-W1}.
Their support has diameter at most $2R$, so TV times $2R$ bounds $W_1$.
Finally, couple $A\sim q_s$ to
$X_1=\sqrt{1-v_1}A+\sqrt{v_1}Z$ and project; nonexpansiveness gives
$W_1(q_s,\operatorname{Proj}_{\mathcal A}\#q_{s,1})\le b_K$.
\end{proof}

\paragraph{Explicit accuracy and cost.}
For $0<\varepsilon\le1$, take
\[
 v_1=\min\left\{\frac14,\frac{\varepsilon^2}{16(R+\sqrt d)^2}\right\},
 \quad u=1-v_K=\min\left\{\frac14,
 \frac{\varepsilon^2}{12R^2(R^2+d)}\right\},\quad
 D=v_1^{-1}-v_K^{-1}.
\]
Space $v_k^{-1}$ equally between $v_1^{-1}$ and $v_K^{-1}$ and choose
\[
 K-1\ge\max\{4R^2D,12eR^6D^2/\varepsilon^2\}.
\]
Then $b_K\le\varepsilon/4$ and
$I_K,B_K\le\varepsilon^2/(24R^2)$.
If
\[
 a_s\le\frac{\varepsilon^2}{216R^2\Lambda},\qquad
 n\ge\frac{6912R^4\Lambda L}{p_0\varepsilon^2},\qquad
 \eta\le\frac{p_0\varepsilon^2}{144R^2\Lambda},
\]
then the learned DDPM has $W_1(q_s,\widehat\pi_s)\le\varepsilon$ at every state.
For fixed $R,d$, these conservative bounds have
$K=O(\varepsilon^{-6})$ and
$n=O(p_0^{-1}\varepsilon^{-4}\log(|\mathcal G|SHT/\delta))$,
besides the explicit approximation requirement.
Corollary~\ref{cor:gibbs-softmax-approximation} constructs a finite neural
class satisfying the displayed approximation requirement for every bounded
target, using $a=\min\{R^2,\varepsilon^2/(216R^2\Lambda)\}$ there. Its class size must be included in $L$ and in the optimization cost.
This construction uses exhaustive search, not a convergence guarantee for
stochastic gradient descent on a standard MLP or U-Net.

\begin{cor}[A finite neural family with uniform approximation]
\label{cor:gibbs-softmax-approximation}
Fix $R>0$, $d\ge1$, $K\ge2$, $0<v_1<\cdots<v_K<1$, and a compact action set
$\mathcal A\subseteq\{a:\|a\|\le R\}$. For any $0<a\le R^2$, the
following finite family of softmax neural predictors satisfies
\[
 \sup_{\nu:\,\operatorname{supp}(\nu)\subseteq\mathcal A}
 \min_{g\in\mathcal G}\max_{2\le k\le K}
 \mathbb E_{X_k\sim q_k^\nu}\|g(X_k,k)-\mathbb E_\nu[A\mid X_k]\|^2
 \le a,
 \qquad X_k=\sqrt{1-v_k}A+\sqrt{v_k}Z.
\]
Set
\[
 \zeta=\frac{a}{200R^2},\quad
 U=R+\sqrt{2d\log(2d/\zeta)},\quad L_x=\frac{U+R}{v_1},\quad
 h=\min\left\{\frac1{2L_x},\frac{\sqrt a}{8(1+4RL_x)}\right\}.
\]
Supply an $h$-net $\{a_1,\ldots,a_J\}\subseteq\mathcal A$.
A maximal $h$-separated subset gives $J\le(1+2R/h)^d$.
This is an input requirement for computation; for example, a net can be
constructed for an explicitly given box or ball.
Let
\[
 B=\max\left\{\frac R{v_1},\log(J/\zeta)+\frac{R^2}{2v_1}\right\},
 \qquad \delta_w=\frac{\sqrt a}{16R(U\sqrt d+1)}.
\]
The family has a separate head for each $k$. Each head is a convex combination
of at most $J$ output atoms from the supplied net, with weights given by the
softmax of affine functions of $x$. Every slope coordinate and bias ranges
over a finite grid of $[-B,B]$ that approximates each coordinate to within
$\delta_w$. Units may be masked, with at least one active unit. In particular,
\[
 |\mathcal G|\le
 \left[2J\left(2+\frac{2B}{\delta_w}\right)^{d+1}\right]^{J(K-1)}.
\]
All output atoms, slopes, and biases can vary in the finite optimization.
Choosing this class in the weighted-ERM theorem verifies its approximation
premise for every Gibbs target, with the displayed class size included in its
statistical and computational costs. The resulting actor is a normalized-exponential neural family fitted by exhaustive search.
\end{cor}

\begin{proof}
Fix any target $\nu$. Gaussian tails give
$\mathbb{P}(\|X_k\|>U)\le\zeta$, uniformly over $k$ and over all laws supported
on the action ball. Map each action to a nearest net point $Q(A)$, and let
$\nu^D=Q_\#\nu$ have atom weights $p_j$.
On $\|x\|\le U$, the Gaussian likelihoods obey
\[
 e^{-L_xh}\le
 \frac{\exp(-\|x-\sqrt{1-v_k}A\|^2/(2v_k))}
      {\exp(-\|x-\sqrt{1-v_k}Q(A)\|^2/(2v_k))}
 \le e^{L_xh}.
\]
Thus their normalized posterior density ratio is between
$e^{-2L_xh}$ and $e^{2L_xh}$. Writing $m_k,m_k^D$ for the two posterior
means, first replacing $A$ by $Q(A)$ and then changing posterior weights gives
\[
 \|m_k(x)-m_k^D(x)\|
 \le h+R(e^{2L_xh}-1)\le\sqrt a/4.
\]
The last step uses $2L_xh\le1$ and $e^{2L_xh}-1\le4L_xh$.
All posterior means have norm at most $R$, so
\[
 \mathbb E_{q_k^\nu}\|m_k-m_k^D\|^2\le a/16+4R^2\zeta.
\]

Drop the atoms with $p_j<\zeta/J$ and renormalize, obtaining $\nu^-$.
The deleted mass is at most $\zeta$, and every retained normalized weight
$p_j^-$ lies in $[\zeta/J,1]$.
For any joint laws $P,Q$, disintegration and the triangle inequality give
\[
 \int P_X(dx)\,d_{\rm TV}(P(\cdot\mid x),Q(\cdot\mid x))
 \le2d_{\rm TV}(P,Q).
\]
Apply this to the joint laws of the atom and its Gaussian corruption under
$\nu^D$ and $\nu^-$. Conditional mean differences are at most $2R$ times
conditional TV. Therefore
$\mathbb E_{q_k^{\nu^D}}\|m_k^D-m_k^-\|^2\le8R^2\zeta$.
On the input ball $q_k^\nu/q_k^{\nu^D}\le e^{L_xh}$; outside it use the
$2R$ bound. This yields
\[
 \mathbb E_{q_k^\nu}\|m_k^D-m_k^-\|^2
 \le8R^2\zeta e^{L_xh}+4R^2\zeta.
\]

For the retained atoms the posterior mean is exactly
\[
 m_k^-(x)=\sum_j a_j\,\operatorname{softmax}_j(b_{jk}+u_{jk}^\top x),
 \quad u_{jk}=\frac{\sqrt{1-v_k}}{v_k}a_j,
 \quad b_{jk}=\log p_j^--\frac{1-v_k}{2v_k}\|a_j\|^2.
\]
Every parameter lies in the specified box. Round each coordinate to the
finite grid. On $\|x\|\le U$, each logit changes by at most
$c=\delta_w(U\sqrt d+1)=\sqrt a/(16R)\le1/16$.
Normalized softmax weights change by density ratios in $[e^{-2c},e^{2c}]$,
so the resulting mean $g(x,k)$ differs from $m_k^-(x)$ by at most
$R(e^{2c}-1)\le\sqrt a/4$. Outside the input ball both outputs remain
bounded by $R$. Hence
\[
 \mathbb E_{q_k^\nu}\|m_k^--g\|^2\le a/16+4R^2\zeta.
\]
Combining the three errors gives
\[
 \mathbb E_{q_k^\nu}\|m_k-g\|^2
 \le\frac{3a}{8}+24R^2\zeta e^{1/2}+36R^2\zeta<a.
\]
The construction works simultaneously for all heads. A unit has two mask
choices, $J$ output-atom choices, and at most
$(2+2B/\delta_w)^{d+1}$ affine-logit choices. Counting $J$ units in each of
$K-1$ heads proves the cardinality bound. A Euclidean packing of radius-$h/2$
balls proves the stated bound on $J$.
\end{proof}

At each round and stage, fitting all $S$ states costs $Sn$ pretrained samples
and critic evaluations, $SnK_0$ pretrained denoising calls if a pretrained
sample uses $K_0$ calls, $S$ greedy oracle calls, and at most
$Sn|\mathcal G|$ candidate network evaluations for exhaustive fitting.
Deployed actions cost $K-1$ learned denoising calls each. These synthetic
calls do not create additional environment transitions.
Setting $\varepsilon_t=t^{-1/2}$, $\tau_t=t^{-1/2}$ and the Gibbs radius
$r_t=\min\{R_0,m\tau_t/L_f\}$ gives
\[
 \mathrm{Reg}(T)\le\mathcal C_T+
 H\sum_{t=1}^T g_{\tau_t,r_t}
 +L_fH\sum_{t=1}^T\varepsilon_t.
\]
The two latter sums are $\widetilde O(\sqrt T)$ for fixed problem constants.
With a uniform class-size bound $\mathcal N_T$ for the first $T$ rounds,
the synthetic sample count is at most
$\widetilde O(SH T^{3+m/2}\log(2\mathcal N_T))$.
The explicit class in Corollary~\ref{cor:gibbs-softmax-approximation}
can make this factor large and dimension-dependent. Deployed actions use
$O(HT^4)$ learned denoising calls, each evaluating a head of width at most $J$.
The critic bound must use the policy class containing these deployed actors.
This finite-state construction does not certify that its SEC is small.

The Wasserstein guarantee for the fitted actor also gives an action-region
bound corresponding to Appendix~\ref{app:favorable-actions}, although it
does not supply TV control for an arbitrary region:
\[
 \widehat\pi(B)\ge q(B^{-r})-W_1(q,\widehat\pi)/r,
 \quad B^{-r}=\{a:\operatorname{dist}(a,B^c)>r\},\quad r>0.
\]
This follows by a coupling and Markov's inequality. Combined with the
earlier bound for the target mass $q(B^{-r})$, it quantifies which favorable
regions remain likely under the learned actor.

\section{Miscellaneous Lemmas}


\begin{lem}[Value Difference/Generalized Policy Difference Lemma, \citep{cai2024provablyefficientexplorationpolicy, efroni2020optimisticpolicyoptimizationbandit}]
    Let $\pi,\pi'$ be two policies and $f\in\gF$ with $f_{H+1}=0$. We interpret $f_h(s,\pi_h(s))$ as $\E_{a\sim\pi_h(\cdot\mid s)}f_h(s,a)$. Then
$$
\begin{aligned}
& f_1\left(s_1, \pi_1(s_1)\right)-V_1^{\pi^{\prime}}\left(s_1\right) \\
& \quad=\sum_{h=1}^H \mathbb{E}_{\pi^{\prime}}\left[\left\langle f_h\left(s_h, \cdot\right), \pi_h\left(\cdot \mid s_h\right)-\pi_h^{\prime}\left(\cdot \mid s_h\right)\right\rangle\right]+\sum_{h=1}^H \mathbb{E}_{\pi^{\prime}}\left[f_h\left(s_h, a_h\right)-\gT_h^{\pi}f_{h+1}\left(s_h, a_h\right)\right] .
\end{aligned}
$$
\label{ch05:lem:policy-value-difference-lemma}
\end{lem}

\begin{lem}[Performance Difference Lemma \citep{kakade2002approxoptimalrl}]
    For all policies $\pi, \pi'$, and states $s_0 \in \gS$,
    $$V^\pi(s_0) - V^{\pi'}(s_0) = \frac{1}{1-\gamma} \E_{d_{s_0}^\pi}\left[A^{\pi'}(s,a)\right].$$
    \label{lem:pdl}
\end{lem}

%%%%%%%%%%%%%%%%%%%%%%%%%%%%%%%%%%%%%%%%%%%%%%%%%%%%%%%%%%%%

%%%%%%%%%%%%%%%%%%%%%%%%%%%%%%%%%%%%%%%%%%%%%%%%%%%%%%%%%%%%





